% Le Droit et la Justice — Philosophie Tle Tech — Cours signé OnBuch+
% Programme officiel MINESEC. Compiler avec : tectonic N06-droit-justice.tex
\newcommand{\DOCMATIERE}{Philosophie}
\newcommand{\DOCNIVEAU}{Tle Tech}
\newcommand{\DOCMODULE}{Philosophie – classe de Terminale technique}
\newcommand{\DOCLECON}{N06}
\newcommand{\DOCTITRE}{Le Droit et la Justice}
% OnBuch+ — préambule « cours signés » ENRICHI (compilé avec tectonic / XeLaTeX)
% Système visuel « Pop » d'OnBuch+ : fond crème (jamais blanc), bordures encre
% épaisses, ombres « dures » décalées, titres Archivo Black en capitales.
% Version MATHÉMATIQUES : graphes (pgfplots/tikz), boîtes pédagogiques
% (définition, propriété, méthode, attention, le savais-tu, exercice résolu,
% exercice, TP, bilan), page de garde et signature OnBuch+ ancrée en bas.
%
% Macros attendues dans main.tex AVANT \input{preamble} :
%   \DOCMATIERE  \DOCNIVEAU  \DOCMODULE  \DOCLECON (numéro)  \DOCTITRE
\documentclass[11pt]{article}

\usepackage{fontspec}
\usepackage[french]{babel}
\usepackage[a4paper,margin=2cm,top=3.4cm,bottom=2.8cm,headheight=1.4cm,footskip=1.3cm]{geometry}
\usepackage{amsmath,amssymb,amsfonts}
\usepackage{mathtools} % \xmapsto, \coloneqq... souvent utilisés par les modèles
\usepackage{cancel} % simplifications barrées (bilans de forces)
% \abs{z} (module, valeur absolue) et \norm{u} (norme), utilisés par les modèles
\providecommand{\abs}{}\let\abs\relax\DeclarePairedDelimiter{\abs}{\lvert}{\rvert}
\providecommand{\norm}{}\let\norm\relax\DeclarePairedDelimiter{\norm}{\lVert}{\rVert}
\usepackage[table]{xcolor} % [table] : fournit aussi \rowcolors (alternance de lignes)
\usepackage{pagecolor}
\usepackage{tikz}
\usetikzlibrary{arrows.meta,positioning,calc,shapes.geometric,shapes.misc,shapes.arrows,decorations.pathmorphing,decorations.pathreplacing,decorations.markings,patterns,shadows,mindmap,trees,fit,backgrounds,matrix,intersections,angles,quotes}
\usepackage{pgfplots}
\usepackage[version=4]{mhchem} % équations nucléaires \ce{^{235}_{92}U}
\usepackage[european,siunitx]{circuitikz} % schémas électriques
\pgfplotsset{compat=1.18}
\usepgfplotslibrary{fillbetween,groupplots} % aires entre courbes (calcul intégral)
\usepackage{enumitem}
\usepackage{fancyhdr}
% [most] charge toutes les bibliothèques tcolorbox (listings, minted, raster,
% documentation...) et gonfle inutilement la mémoire TeX sur un document long
% (~300 boîtes) : on ne charge que ce qui est réellement utilisé.
\usepackage[skins,breakable]{tcolorbox}
\usepackage{graphicx}
\usepackage{colortbl}
\usepackage{booktabs}
\usepackage{tabularx}
\usepackage{multirow}
\usepackage{diagbox} % case d'en-tête diagonale des tableaux à double entrée
% Colonnes extensibles centrée / gauche / droite (C, L, R), souvent utilisées par les modèles
\newcolumntype{C}{>{\centering\arraybackslash}X}
\newcolumntype{L}{>{\raggedright\arraybackslash}X}
\newcolumntype{R}{>{\raggedleft\arraybackslash}X}
\usepackage{array}
\usepackage{multicol}
\usepackage{siunitx}
% parse-numbers=false : accepte aussi \SI{2,5 \times 10^{-3}}{...}
\sisetup{locale=FR,output-decimal-marker={,},group-minimum-digits=4,parse-numbers=false}
\DeclareSIUnit\ohm{\ensuremath{\symup{\Omega}}} % U+2126 absent de la police
\DeclareSIUnit\division{div} % réglages d'oscilloscope (ms/div, V/div)
\DeclareSIUnit\tr{tr} % tours (vitesse de rotation, ex. \SI{1500}{\tr\per\minute})
\DeclareSIUnit\molar{M} % concentration molaire (ex. \SI{3}{\molar}, \SI{1}{\micro\molar})
\DeclareSIUnit\an{an} % année (le modèle confond parfois avec \year, un compteur TeX) — ex. \SI{5}{\kilo\meter\per\an}
\DeclareSIUnit\FCFA{FCFA} % franc CFA (ex. \SI{500}{\FCFA\per\kilo\gram})
\DeclareSIUnit\degreeBrix{\textdegree Brix} % teneur en sucre d'un sirop/jus (réfractométrie)
\DeclareSIUnit\month{mois} % le modèle utilise parfois \month, un compteur TeX, comme unité de durée
\DeclareSIUnit\microliter{\micro\liter} % le modèle écrit parfois \microliter en un seul token
\DeclareSIUnit\million{millions} % dénombrement (ex. \SI{15}{\million\per\mL} spermatozoïdes)
\usepackage{qrcode}
% \degree : commande fréquemment utilisée par les modèles pour un angle ou une
% température en dehors de \SI{}{} (non fournie par les paquets chargés).
\providecommand{\degree}{^{\circ}}
% \qed (fin de démonstration), utilisé par les modèles sans amsthm
\providecommand{\qed}{\hfill$\square$}
% \barre : double barre des valeurs interdites dans un tableau de variations (tkz-tab)
\providecommand{\barre}{\ensuremath{\|}}
% environnement proof (amsthm non chargé) utilisé par les modèles
\ifdefined\proof\else\newenvironment{proof}[1][Démonstration]{\par\noindent\textit{#1.}\ }{\qed\par}\fi
\usepackage{lastpage}
\usepackage{hyperref}
\hypersetup{hidelinks}

% ============================================================
% Palette « Pop » OnBuch+ — mêmes hex que la classe P (plus_ui.dart)
% ============================================================
\definecolor{poporange}{HTML}{F59321}
\definecolor{popdark}{HTML}{E07A0C}
\definecolor{popcream}{HTML}{FFF6E8}
\definecolor{popink}{HTML}{1C1714}
\definecolor{poppurple}{HTML}{7A5AE0}
\definecolor{popgreen}{HTML}{2E9E6A}
\definecolor{poppink}{HTML}{E0457B}
\definecolor{popblue}{HTML}{2F7DE1}
\definecolor{popgold}{HTML}{C98A12}
\definecolor{popmuted}{HTML}{8A7F76}
\definecolor{popline}{HTML}{EADFD2}
\definecolor{popgray}{HTML}{8A7F76}
\colorlet{popgrey}{popgray}
% Alias tolérants (le modèle nomme parfois les couleurs différemment)
\colorlet{popwhite}{white}
\colorlet{popblack}{popink}
\colorlet{popred}{poppink}
\colorlet{popyellow}{popgold}
% Teintes claires (fonds de tableaux/encadrés) — le modèle nomme parfois
% les couleurs avec un suffixe L (ex. popblueL) sans les avoir définies.
\colorlet{poporangeL}{poporange!20}
\colorlet{popdarkL}{popdark!20}
\colorlet{poppurpleL}{poppurple!20}
\colorlet{popgreenL}{popgreen!20}
\colorlet{poppinkL}{poppink!20}
\colorlet{popblueL}{popblue!20}
\colorlet{popgoldL}{popgold!20}
\colorlet{popredL}{popred!20}
\colorlet{popgrayL}{popgray!20}
% Teintes claires pour fonds de tableaux / zones
\colorlet{popblueL}{popblue!10!white}
\colorlet{poporangeL}{poporange!14!white}
\colorlet{popgreenL}{popgreen!12!white}
\colorlet{poppurpleL}{poppurple!12!white}
\colorlet{poppinkL}{poppink!12!white}
\colorlet{popgoldL}{popgold!14!white}

\pagecolor{popcream}

% ============================================================
% Typographie
% ============================================================
\defaultfontfeatures{Ligatures=TeX}
\setmainfont{PlusJakartaSans-Medium.ttf}[Path=../../../fonts/,BoldFont=PlusJakartaSans-Bold.ttf]
% Maths en sans-serif (Fira Math) avec chiffres et lettres droites de Plus
% Jakarta, pour que les formules se fondent dans le texte.
\usepackage[math-style=ISO,bold-style=ISO]{unicode-math}
\setmathfont{FiraMath-Regular.otf}[Path=../../../fonts/]
\setmathfont{PlusJakartaSans-Medium.ttf}[Path=../../../fonts/,range={up/num,up/{Latin,latin},\mathup}]
\usepackage{icomma} % virgule décimale sans espace en mode math
% Caractères mathématiques Unicode absents des polices de texte (Plus Jakarta,
% Archivo Black) : rendus par leur équivalent mathématique, en texte comme en
% formule — sinon ils s'affichent en carré vide (ex. « Arithmétique dans ℤ »).
\usepackage{newunicodechar}
\newunicodechar{ℝ}{\ensuremath{\mathbb{R}}}
\newunicodechar{ℤ}{\ensuremath{\mathbb{Z}}}
\newunicodechar{ℂ}{\ensuremath{\mathbb{C}}}
\newunicodechar{ℕ}{\ensuremath{\mathbb{N}}}
\newunicodechar{ℚ}{\ensuremath{\mathbb{Q}}}
\newunicodechar{𝔻}{\ensuremath{\mathbb{D}}}
\newunicodechar{α}{\ensuremath{\alpha}}
\newunicodechar{β}{\ensuremath{\beta}}
\newunicodechar{γ}{\ensuremath{\gamma}}
\newunicodechar{δ}{\ensuremath{\delta}}
\newunicodechar{ε}{\ensuremath{\varepsilon}}
\newunicodechar{θ}{\ensuremath{\theta}}
\newunicodechar{λ}{\ensuremath{\lambda}}
\newunicodechar{μ}{\ensuremath{\mu}}
\newunicodechar{σ}{\ensuremath{\sigma}}
\newunicodechar{τ}{\ensuremath{\tau}}
\newunicodechar{φ}{\ensuremath{\varphi}}
\newunicodechar{ψ}{\ensuremath{\psi}}
\newunicodechar{ω}{\ensuremath{\omega}}
\newunicodechar{Σ}{\ensuremath{\Sigma}}
% majuscules produites par \MakeUppercase dans les titres (α-aminés -> Α)
\newunicodechar{Α}{\ensuremath{\alpha}}
\newunicodechar{Β}{\ensuremath{\beta}}
\newunicodechar{Γ}{\ensuremath{\Gamma}}
\newunicodechar{Δ}{\ensuremath{\Delta}}
\newunicodechar{Ε}{\ensuremath{\varepsilon}}
\newunicodechar{Θ}{\ensuremath{\Theta}}
\newunicodechar{Λ}{\ensuremath{\Lambda}}
\newunicodechar{Μ}{\ensuremath{\mu}}
\newunicodechar{Τ}{\ensuremath{\tau}}
\newunicodechar{Φ}{\ensuremath{\Phi}}
\newunicodechar{Ψ}{\ensuremath{\Psi}}
\newunicodechar{Ω}{\ensuremath{\Omega}}
\newunicodechar{↦}{\ensuremath{\mapsto}}
\newunicodechar{∈}{\ensuremath{\in}}
\newunicodechar{∉}{\ensuremath{\notin}}
\newunicodechar{∧}{\ensuremath{\wedge}}
\newunicodechar{⇔}{\ensuremath{\Leftrightarrow}}
\newunicodechar{⟺}{\ensuremath{\Longleftrightarrow}}
\newunicodechar{⇒}{\ensuremath{\Rightarrow}}
\newunicodechar{⟹}{\ensuremath{\Longrightarrow}}
\newunicodechar{∘}{\ensuremath{\circ}}
\newunicodechar{∪}{\ensuremath{\cup}}
\newunicodechar{∩}{\ensuremath{\cap}}
\newunicodechar{≡}{\ensuremath{\equiv}}
\newunicodechar{⊂}{\ensuremath{\subset}}
\newunicodechar{⊥}{\ensuremath{\perp}}
\newunicodechar{∃}{\ensuremath{\exists}}
\newunicodechar{∀}{\ensuremath{\forall}}
\newunicodechar{∛}{\ensuremath{\sqrt[3]{\,}}}
\newunicodechar{∜}{\ensuremath{\sqrt[4]{\,}}}
\newunicodechar{⃗}{\ensuremath{{}^{\to}}}
\newunicodechar{ⁿ}{\ifmmode^{n}\else\textsuperscript{\ensuremath{n}}\fi}
\newunicodechar{ˣ}{\ifmmode^{x}\else\textsuperscript{\ensuremath{x}}\fi}
\newunicodechar{ᵇ}{\ifmmode^{b}\else\textsuperscript{\ensuremath{b}}\fi}
\newunicodechar{ʳ}{\ifmmode^{r}\else\textsuperscript{\ensuremath{r}}\fi}
\newunicodechar{ᵘ}{\ifmmode^{u}\else\textsuperscript{\ensuremath{u}}\fi}
\newunicodechar{ʸ}{\ifmmode^{y}\else\textsuperscript{\ensuremath{y}}\fi}
\newunicodechar{ᵃ}{\ifmmode^{a}\else\textsuperscript{\ensuremath{a}}\fi}
\newunicodechar{ᵉ}{\ifmmode^{e}\else\textsuperscript{\ensuremath{e}}\fi}
\newunicodechar{ᵏ}{\ifmmode^{k}\else\textsuperscript{\ensuremath{k}}\fi}
\newunicodechar{ᵖ}{\ifmmode^{p}\else\textsuperscript{\ensuremath{p}}\fi}
\newunicodechar{ᴺ}{\ifmmode^{N}\else\textsuperscript{\ensuremath{N}}\fi}
\newunicodechar{ᶜ}{\ifmmode^{c}\else\textsuperscript{\ensuremath{c}}\fi}
\newunicodechar{ʰ}{\ifmmode^{h}\else\textsuperscript{\ensuremath{h}}\fi}
\newunicodechar{ᵠ}{\ifmmode^{q}\else\textsuperscript{\ensuremath{q}}\fi}
\newunicodechar{ₙ}{\ifmmode_{n}\else\textsubscript{\ensuremath{n}}\fi}
\newunicodechar{ₐ}{\ifmmode_{a}\else\textsubscript{\ensuremath{a}}\fi}
\newunicodechar{ᵢ}{\ifmmode_{i}\else\textsubscript{\ensuremath{i}}\fi}
\newunicodechar{ᵣ}{\ifmmode_{r}\else\textsubscript{\ensuremath{r}}\fi}
\newunicodechar{ₖ}{\ifmmode_{k}\else\textsubscript{\ensuremath{k}}\fi}
\newunicodechar{ᵦ}{\ifmmode_{\beta}\else\textsubscript{\ensuremath{\beta}}\fi}
\newunicodechar{ₓ}{\ifmmode_{x}\else\textsubscript{\ensuremath{x}}\fi}
\newfontfamily\popdisplay{ArchivoBlack-Regular.ttf}[Path=../../../fonts/]
\newfontfamily\poplogo{SpaceGrotesk-Bold.ttf}[Path=../../../fonts/]
\newfontfamily\popsemi{PlusJakartaSans-SemiBold.ttf}[Path=../../../fonts/]
\newfontfamily\popxbold{PlusJakartaSans-ExtraBold.ttf}[Path=../../../fonts/]
\newfontfamily\popxboldls{PlusJakartaSans-ExtraBold.ttf}[Path=../../../fonts/,LetterSpace=3.0]

\setlist[enumerate]{leftmargin=1.7em,itemsep=5pt,topsep=4pt}
\setlist[itemize]{leftmargin=1.7em,itemsep=4pt,topsep=4pt,label={\color{poporange}\textbullet}}
\setlength{\parindent}{0pt}
\setlength{\parskip}{5pt}
\renewcommand{\arraystretch}{1.25}

% pgfplots : style Pop par défaut
\pgfplotsset{
  every axis/.append style={axis line style={popink,line width=1pt},tick style={popink},
    grid style={popline,line width=0.5pt},label style={font=\small},tick label style={font=\footnotesize}},
  popcurve/.style={poporange,line width=1.8pt,no markers,smooth},
}
% Même style disponible dans un \draw TikZ simple (hors pgfplots)
\tikzset{popcurve/.style={poporange,line width=1.8pt}}
\tikzset{popgrid/.style={popline,very thin}}
\colorlet{popcurve}{poporange} % le nom du style est parfois utilisé comme couleur

% ============================================================
% Pill / squiggle
% ============================================================
\newcommand{\poppill}[3]{%
  \tikz[baseline=-0.55ex]\node[draw=popink,line width=1.2pt,rounded corners=2.6pt,%
    fill=#2,inner xsep=6.5pt,inner ysep=3pt]{%
    \popxboldls\fontsize{7}{7}\selectfont\color{#3}\MakeUppercase{#1}};%
}
\newcommand{\popsquiggle}[1]{%
  \tikz[baseline=-0.4ex]{
    \draw[#1,line width=2.6pt,line cap=round]
      plot [smooth] coordinates {(0,0.09) (0.34,0.01) (0.68,0.21) (1.02,0.01) (1.36,0.10)};
  }%
}
% Mot-clé surligné (marqueur orange sous le texte)
\newcommand{\cle}[1]{{\popxbold\color{popdark}#1}}

% ============================================================
% En-tête / pied
% ============================================================
\pagestyle{fancy}
\fancyhf{}
\renewcommand{\headrulewidth}{0pt}
\renewcommand{\footrulewidth}{0pt}
\lhead{\raisebox{-0.15ex}{\poplogo\fontsize{13}{13}\selectfont\color{popink}OnBuch\color{poporange}+}}
\rhead{\poppill{\DOCMATIERE\ \textbullet\ \DOCNIVEAU\ \textbullet\ Leçon \DOCLECON}{white}{popink}}
\lfoot{\popxbold\fontsize{7.6}{7.6}\selectfont\color{popmuted}\MakeUppercase{Cours signé OnBuch+}\ \textbullet\ {\popsemi\DOCTITRE}}
\rfoot{\tikz[baseline=-0.6ex]\node[rounded corners=8.5pt,fill=popink,minimum height=17pt,minimum width=17pt,inner xsep=5pt,inner sep=0]{\popxbold\fontsize{8}{8}\selectfont\color{popcream}\thepage\,/\,\pageref*{LastPage}};}

% ============================================================
% Titres
% ============================================================
\newcommand{\chaptitre}[2]{%
  \begin{tcolorbox}[enhanced,colback=poporange,colframe=popink,boxrule=2.6pt,arc=11pt,
      drop shadow=popink,left=15pt,right=15pt,top=12pt,bottom=11pt,before skip=3pt,after skip=18pt]
    \poppill{#2}{popink}{popcream}\par\vspace{9pt}
    {\raggedright\hyphenpenalty=10000\exhyphenpenalty=10000\popdisplay\fontsize{22}{24}\selectfont\color{popcream}\MakeUppercase{#1}\par}
  \end{tcolorbox}
}
% Section numérotée : pastille orange avec le numéro + titre Archivo
\newcounter{popsec}
\newcommand{\coursec}[1]{%
  \stepcounter{popsec}\setcounter{popsub}{0}%
  \par\vspace{16pt}\noindent
  \begin{minipage}{\linewidth}
  \tikz[baseline=-0.7ex]\node[draw=popink,line width=1.6pt,fill=poporange,rounded corners=4pt,minimum size=22pt,inner sep=0]{\popdisplay\fontsize{12}{12}\selectfont\color{popcream}\thepopsec};%
  \hspace{8pt}{\popdisplay\fontsize{15}{17}\selectfont\color{popink}\MakeUppercase{#1}}\par
  \vspace{4pt}\noindent{\color{poporange}\rule{36pt}{3.6pt}}
  \end{minipage}\par\nopagebreak\vspace{8pt}
}
\newcounter{popsub}
\newcommand{\courssub}[1]{%
  \stepcounter{popsub}%
  \par\vspace{9pt}\noindent{\popxbold\fontsize{11.5}{13}\selectfont\color{popink}{\color{poporange}\thepopsec.\thepopsub}\hspace{6pt}#1}\par\nopagebreak\vspace{4pt}
}

% ============================================================
% Boîtes pédagogiques — même recette : fond blanc, bordure épaisse colorée,
% ombre dure encre, badge plein. Titre optionnel [#1].
% ============================================================
\tcbset{popbase/.style={breakable,enhanced,colback=white,boxrule=2.3pt,arc=9pt,
  shadow={3pt}{-3pt}{0pt}{popink},left=14pt,right=14pt,top=11pt,bottom=11pt,before skip=11pt,after skip=15pt,
  fontupper=\popsemi\fontsize{10.6}{14.5}\selectfont\color{popink},
  fontlower=\popsemi\fontsize{10.6}{14.5}\selectfont\color{popink}}}
% Coche dessinée (le glyphe ✓ n'existe pas dans les polices du document)
\newcommand{\popcheck}[1]{\tikz[baseline=-0.5ex]\draw[#1,line width=1.6pt,line cap=round,line join=round](-0.13,0.0)--(-0.04,-0.09)--(0.13,0.1);}
\providecommand{\coche}{\popcheck{popgreen}}
\providecommand{\resultat}[1]{\par\smallskip\noindent\fcolorbox{popgreen}{popgreenL}{\parbox{\dimexpr\linewidth-2\fboxsep-2\fboxrule\relax}{\textbf{Résultat.} #1}}\par\smallskip}
\newcommand{\popboxhead}[3]{\poppill{#1}{#2}{white}\ifx\relax#3\relax\else\hspace{6pt}{\popxbold\fontsize{10.6}{12}\selectfont\color{popink}#3}\fi\par\vspace{7pt}}

\newtcolorbox{aretenir}[1][]{popbase,colframe=popgreen,before upper={\popboxhead{À retenir}{popgreen}{#1}}}
\newtcolorbox{exemplebox}[1][]{popbase,colframe=poporange,before upper={\popboxhead{Exemple}{poporange}{#1}}}
\newtcolorbox{definition}[1][]{popbase,colframe=popblue,before upper={\popboxhead{Définition}{popblue}{#1}}}
\newtcolorbox{propriete}[1][]{popbase,colframe=poppurple,before upper={\popboxhead{Propriété}{poppurple}{#1}}}
\newtcolorbox{methode}[1][]{popbase,colframe=popgold,before upper={\popboxhead{Méthode}{popgold}{#1}}}
\newtcolorbox{attention}[1][]{popbase,colframe=poppink,before upper={\popboxhead{Attention}{poppink}{#1}}}
\newtcolorbox{experience}[1][]{popbase,colframe=popblue,colback=popblueL,before upper={\popboxhead{Expérience}{popblue}{#1}}}
% « Le savais-tu ? » : carte violette pleine (contexte, culture, Cameroun)
\newtcolorbox{savaistu}[1][]{popbase,colback=poppurple,colframe=popink,
  fontupper=\popsemi\fontsize{10.4}{14.2}\selectfont\color{white},
  before upper={\poppill{Le savais-tu\,?}{white}{poppurple}\ifx\relax#1\relax\else\hspace{6pt}{\popxbold\color{white}#1}\fi\par\vspace{7pt}}}
% Exercice résolu : énoncé en haut, \tcblower puis la solution sur fond crème
\newcounter{popexo}\newcounter{popexr}
\newtcolorbox{exoresolu}[1][]{popbase,colframe=poporange,colbacklower=popcream,
  segmentation style={popink,line width=1.2pt,dashed},
  before upper={\stepcounter{popexr}\popboxhead{Exercice résolu \thepopexr}{poporange}{#1}},
  before lower={\poppill{Solution}{popink}{popcream}\par\vspace{6pt}}}
% Exercice (non résolu en place) — la correction est dans la partie Corrigés
\newtcolorbox{exercice}[1][]{popbase,colframe=popink,
  before upper={\stepcounter{popexo}\popboxhead{Exercice \thepopexo}{popink}{#1}}}
% Corrigé
\newtcolorbox{corrige}[1][]{popbase,colframe=popgreen,colback=popgreenL,
  before upper={\popboxhead{Corrigé}{popgreen}{#1}}}
% Objectifs / prérequis / bilan
\newtcolorbox{objectifs}{popbase,colframe=popink,colback=poporangeL,
  before upper={\popboxhead{Objectifs de la leçon}{popink}{}}}
\newtcolorbox{prerequis}{popbase,colframe=popink,colback=popblueL,
  before upper={\popboxhead{Prérequis}{popblue}{}}}
\newtcolorbox{bilan}{popbase,colframe=popink,colback=white,boxrule=2.8pt,
  before upper={\popboxhead{Fiche bilan}{poporange}{L'essentiel à connaître pour l'examen}}}
% Figure encadrée avec légende
\newtcolorbox{popfigure}[1][]{enhanced,breakable=false,colback=white,colframe=popline,boxrule=1.4pt,arc=8pt,
  left=10pt,right=10pt,top=10pt,bottom=8pt,before skip=10pt,after skip=12pt,halign=center,
  lower separated=false,halign lower=center,
  fontlower=\popsemi\fontsize{9}{11.5}\selectfont\color{popmuted},
  #1}
\newcommand{\legende}[1]{\par\vspace{6pt}{\popsemi\fontsize{9}{11.5}\selectfont\color{popmuted}#1}}

% ============================================================
% Signature OnBuch+
% ============================================================
\newcommand{\poplogoinline}[1]{{\poplogo\fontsize{#1}{#1}\selectfont\color{popink}OnBuch\color{poporange}+}}
\newcommand{\signatureonbuch}{%
  \begin{tcolorbox}[enhanced,colback=popink,colframe=popink,boxrule=2.2pt,arc=10pt,
      drop shadow=poporange,left=14pt,right=14pt,top=10pt,bottom=10pt,before skip=0pt,after skip=0pt]
    \tikz[baseline=-0.7ex]\node[circle,fill=popgreen,draw=popcream,line width=1.3pt,minimum size=22pt,inner sep=0]{\popcheck{white}};%
    \hspace{9pt}\begin{minipage}[c]{0.62\linewidth}
      {\popxbold\fontsize{12}{13}\selectfont\color{popcream}Signé\ \textbullet\ {\poplogo OnBuch\color{poporange}+}}\par\vspace{2pt}
      {\raggedright\popsemi\fontsize{8}{10.5}\selectfont\color{popline}\MakeUppercase{Équipe pédagogique OnBuch+}\\\MakeUppercase{Conforme au programme officiel MINESEC}\par}
    \end{minipage}\hfill
    {\poplogo\fontsize{22}{22}\selectfont\color{popcream}OnBuch\color{poporange}+}
  \end{tcolorbox}%
}

% ============================================================
% Page de garde
% #1 titre · #2 matière · #3 niveau · #4 module · #5 n° leçon · #6 durée ·
% #7 description / objectifs (texte ou itemize)
% ============================================================
\newcommand{\pagegarde}[7]{%
  \thispagestyle{empty}
  \newgeometry{a4paper,margin=1.7cm,top=1.6cm,bottom=1.4cm}%
  \begin{tikzpicture}[remember picture,overlay]
    \fill[poporange!18] ([xshift=-3.2cm,yshift=-3.6cm]current page.north east) circle (4.6cm);
    \fill[poppurple!14] ([xshift=2.4cm,yshift=7.6cm]current page.south west) circle (3.8cm);
  \end{tikzpicture}%
  {\poplogo\fontsize{30}{30}\selectfont\color{popink}OnBuch\color{poporange}+}\hfill
  \raisebox{0.6ex}{\poppill{Cours signé \textbullet\ Premium}{popink}{popcream}}\par
  \vspace{1pt}\popsquiggle{poppurple}\par
  \vspace{8pt}
  \begin{tcolorbox}[enhanced,colback=poporange,colframe=popink,boxrule=2.8pt,arc=13pt,
      drop shadow=popink,left=18pt,right=18pt,top=12pt,bottom=13pt,after skip=10pt]
    \poppill{#2 \textbullet\ #3}{popink}{popcream}\hspace{5pt}\poppill{#4}{white}{popink}\par\vspace{8pt}
    {\popdisplay\fontsize{14}{14}\selectfont\color{popink}LEÇON #5}\par\vspace{4pt}
    {\raggedright\hyphenpenalty=10000\exhyphenpenalty=10000\popdisplay\fontsize{25}{27}\selectfont\color{popcream}\MakeUppercase{#1}\par}
    \vspace{8pt}
    {\popxbold\fontsize{9.5}{9.5}\selectfont\color{popink}\MakeUppercase{Durée conseillée : #6}}
  \end{tcolorbox}
  \begin{tcolorbox}[enhanced,colback=white,colframe=popink,boxrule=2.2pt,arc=10pt,
      drop shadow=popink,left=16pt,right=16pt,top=10pt,bottom=10pt,after skip=8pt]
    \poppill{Dans cette leçon}{poppurple}{white}\par\vspace{5pt}
    \popsemi\fontsize{9.3}{12.6}\selectfont\color{popink}#7
  \end{tcolorbox}
  \noindent\begin{minipage}[t]{0.487\linewidth}
    \begin{tcolorbox}[enhanced,colback=popgreen,colframe=popink,boxrule=2.2pt,arc=10pt,
        drop shadow=popink,left=11pt,right=11pt,top=10pt,bottom=10pt,height=3.1cm,valign=center]
      \tcbox[colback=white,colframe=popink,boxrule=1.3pt,arc=5pt,boxsep=0pt,nobeforeafter,
        left=3pt,right=3pt,top=3pt,bottom=3pt,valign=center]{\qrcode[height=1.9cm]{https://play.google.com/store/apps/details?id=com.luvvix.onbuch&hl=fr}}%
      \hspace{8pt}\begin{minipage}[c]{0.5\linewidth}
        {\popdisplay\fontsize{11}{12.5}\selectfont\color{white}\MakeUppercase{Télécharge OnBuch}}\par\vspace{4pt}
        {\raggedright\popsemi\fontsize{8.4}{11}\selectfont\color{white}Gratuit \textbullet\ Google Play\\Tuteur IA, annales, concours, résultats\par}
      \end{minipage}
    \end{tcolorbox}
  \end{minipage}\hfill
  \begin{minipage}[t]{0.487\linewidth}
    \begin{tcolorbox}[enhanced,colback=poppurple,colframe=popink,boxrule=2.2pt,arc=10pt,
        drop shadow=popink,left=11pt,right=11pt,top=10pt,bottom=10pt,height=3.1cm,valign=center]
      \tikz[baseline=-0.7ex]\node[circle,fill=white,draw=popink,line width=1.3pt,minimum size=1.6cm,inner sep=0]{\popdisplay\fontsize{24}{24}\selectfont\color{poppurple}+};%
      \hspace{8pt}\begin{minipage}[c]{0.6\linewidth}
        {\popdisplay\fontsize{11}{12.5}\selectfont\color{white}\MakeUppercase{Passe à OnBuch+}}\par\vspace{4pt}
        {\raggedright\popsemi\fontsize{8.4}{11}\selectfont\color{white}Tous les cours signés, Tuteur IA illimité, fiches PDF — onglet OnBuch+ de l'app\par}
      \end{minipage}
    \end{tcolorbox}
  \end{minipage}
  \vspace{8pt}
  \popcontenu
  \vfill
  \signatureonbuch
  \restoregeometry
  \newpage
}

% Rangée « Ce cours contient » (page de garde)
\newcommand{\poptile}[3]{% #1 couleur #2 gros texte #3 légende
  \begin{tcolorbox}[enhanced,colback=#1,colframe=popink,boxrule=2pt,arc=9pt,shadow={2.5pt}{-2.5pt}{0pt}{popink},
      left=6pt,right=6pt,top=9pt,bottom=9pt,halign=center,valign=center,height=1.75cm,width=\linewidth,nobeforeafter]
    {\popdisplay\fontsize{12.5}{13}\selectfont\color{white}\MakeUppercase{#2}}\par\vspace{3pt}
    {\centering\popsemi\fontsize{7.8}{9.5}\selectfont\color{white}#3\par}
  \end{tcolorbox}}
\newcommand{\popcontenu}{%
  \par\noindent{\popxboldls\fontsize{7.5}{7.5}\selectfont\color{popmuted}CE COURS SIGNÉ CONTIENT}\par\vspace{7pt}
  \noindent\begin{minipage}[t]{0.235\linewidth}\poptile{popblue}{Cours}{complet, illustré, pas à pas}\end{minipage}\hfill
  \begin{minipage}[t]{0.235\linewidth}\poptile{poporange}{Méthodes}{et exercices résolus}\end{minipage}\hfill
  \begin{minipage}[t]{0.235\linewidth}\poptile{poppink}{Exercices}{type examen + corrigés}\end{minipage}\hfill
  \begin{minipage}[t]{0.235\linewidth}\poptile{popgreen}{Fiche bilan}{l'essentiel à retenir}\end{minipage}\par}

% Sommaire
\newtcolorbox{sommaire}{enhanced,colback=white,colframe=popink,boxrule=2.2pt,arc=10pt,
  drop shadow=popink,left=18pt,right=18pt,top=13pt,bottom=13pt,before skip=6pt,after skip=20pt,
  before upper={\poppill{Sommaire}{poppurple}{white}\par\vspace{9pt}\popsemi\fontsize{10.6}{14.5}\selectfont\color{popink}}}

% Signature de fin de cours — poussée tout en bas de la dernière page
\newcommand{\findecours}{%
  \par\vfill\nopagebreak
  \begin{center}\popsquiggle{poporange}\end{center}\vspace{4pt}
  \signatureonbuch
}
\AtBeginDocument{\def\llbracket{\mathopen{[\mkern-3.5mu[}}\def\rrbracket{\mathclose{]\mkern-3.5mu]}}} % intervalles d'entiers (glyphe ⟦ absent de la police)
% Glyphes absents de Fira Math : redéfinis à partir de glyphes disponibles
\AtBeginDocument{%
  \def\setminus{\mathbin{\backslash}}\let\smallsetminus\setminus
  \def\vdots{\mathord{\vbox{\baselineskip3.2pt\lineskiplimit0pt\kern4pt\hbox{.}\hbox{.}\hbox{.}}}}%
  \def\ddots{\mathinner{\mkern1mu\raise6.5pt\hbox{.}\mkern2mu\raise3.6pt\hbox{.}\mkern2mu\raise0.7pt\hbox{.}\mkern1mu}}%
  \def\triangle{\mathord{\scalebox{0.9}{$\symup{\Delta}$}}}%
  \def\nabla{\mathord{\raisebox{\depth}{\scalebox{1}[-1]{$\symup{\Delta}$}}}}%
  \def\checkmark{\popcheck{popdark}}%
  \def\star{\mathbin{\ast}}\def\bigstar{\mathord{\ast}}%
  \def\diamond{\mathbin{\scalebox{0.6}{$\Diamond$}}}%
  \def\bigcup{\mathop{\mathchoice{\raisebox{-0.3em}{\scalebox{1.7}{$\cup$}}}{\scalebox{1.25}{$\cup$}}{\cup}{\cup}}\displaylimits}%
  \def\bigcap{\mathop{\mathchoice{\raisebox{-0.3em}{\scalebox{1.7}{$\cap$}}}{\scalebox{1.25}{$\cap$}}{\cap}{\cap}}\displaylimits}%
  \def\lesseqqgtr{\mathrel{\lessgtr}}\def\bigoplus{\mathop{\oplus}\displaylimits}%
}
\providecommand{\ssi}{si et seulement si} % abréviation fréquente
\AtBeginDocument{\providecommand{\wideparen}{\overparen}} % arc de cercle

%% FIGFIX-BEGIN v3 — correctifs globaux des figures (docs/onbuch-generation/tools/figfix)
\usepackage{adjustbox}\usepackage{etoolbox}
% toute figure TikZ/pgfplots plus large que la ligne est réduite à la largeur disponible
\BeforeBeginEnvironment{tikzpicture}{\begin{adjustbox}{max width=\linewidth}}
\AfterEndEnvironment{tikzpicture}{\end{adjustbox}}
% légendes pgfplots lisibles par-dessus les courbes
\pgfplotsset{every axis/.append style={legend style={fill=white,fill opacity=0.92,text opacity=1,draw=black!15,inner sep=3pt}}}
% étiquettes d'arêtes (syntaxe edge["..."]) avec fond blanc
\tikzset{every edge quotes/.append style={fill=white,inner sep=1.2pt}}
% bulles de cartes mentales : texte à la ligne dans la bulle, bulle un peu plus grande
\tikzset{every concept/.append style={text width=2.0cm,align=center,minimum size=2.3cm,inner sep=2pt}}
%% FIGFIX-END
\begin{document}

\pagegarde{Le Droit et la Justice}{Philosophie}{Tle Tech}{Philosophie – classe de Terminale technique}{N06}{3 séances (fiche de progression harmonisée)}{Cette leçon analyse les concepts de droit et de justice, leurs fondements théoriques (droit positif, droit naturel) et leurs déclinaisons pratiques (justice commutative, distributive, réparatrice). Elle outille l'élève pour argumenter philosophiquement sur l'écart entre légalité et légitimité à travers des cas concrets du système judiciaire camerounais.
\vspace{4pt}
\begin{multicols}{2}\small
\begin{itemize}
\item À la fin de cette leçon, je sais définir rigoureusement les concepts de droit, justice, loi, légalité, légitimité et équité.
\item Je sais distinguer le droit positif du droit naturel et expliquer le débat entre légalisme et légitimité (Kelsen, Hart, Finnis).
\item Je sais identifier et caractériser les trois formes classiques de justice (commutative, distributive, réparatrice) et leurs formes modernes (transitionnelle, sociale, environnementale).
\item Je sais analyser une situation concrète camerounaise (conflit foncier, corruption, justice coutumière) en mobilisant les concepts de droit et de justice.
\item Je sais structurer une dissertation philosophique : analyser le sujet, proposer une problématique, construire un plan dialectique et rédiger une introduction/conclusion conformes aux attendus du Baccalauréat.
\item Je sais expliquer un texte philosophique (extrait d'Aristote, Kant, Rawls ou auteur africain) en identifiant la thèse, les arguments et en le critiquant.
\end{itemize}\end{multicols}}

\begin{sommaire}
\begin{multicols}{2}
\begin{enumerate}
\item 1. Définitions et distinctions conceptuelles : Droit, Justice, Loi, Morale
\item 2. Les fondements du Droit : Droit positif vs Justice (Droit naturel)
\item 3. Des formes de la Justice à la Justice comme valeur : Typologie aristotélicienne et moderne
\item 4. La Justice au Cameroun : Institutions, Effectivité et Défis majeurs
\item 5. Méthodologie : Argumenter, Dissertation et Explication de texte (Entraînement Bac)
\item Activité d'intégration
\item Méthodes et exercices résolus
\item Exercices
\item Corrigés des exercices
\item Fiche bilan
\end{enumerate}
\end{multicols}
\end{sommaire}

\begin{objectifs}
\begin{itemize}
\item À la fin de cette leçon, je sais définir rigoureusement les concepts de droit, justice, loi, légalité, légitimité et équité.
\item Je sais distinguer le droit positif du droit naturel et expliquer le débat entre légalisme et légitimité (Kelsen, Hart, Finnis).
\item Je sais identifier et caractériser les trois formes classiques de justice (commutative, distributive, réparatrice) et leurs formes modernes (transitionnelle, sociale, environnementale).
\item Je sais analyser une situation concrète camerounaise (conflit foncier, corruption, justice coutumière) en mobilisant les concepts de droit et de justice.
\item Je sais structurer une dissertation philosophique : analyser le sujet, proposer une problématique, construire un plan dialectique et rédiger une introduction/conclusion conformes aux attendus du Baccalauréat.
\item Je sais expliquer un texte philosophique (extrait d'Aristote, Kant, Rawls ou auteur africain) en identifiant la thèse, les arguments et en le critiquant.
\item Je sais évaluer l'effectivité de la justice au Cameroun (indépendance, accès au droit, délais, corruption) à l'aide de données chiffrées et de documents officiels.
\item Je sais argumenter à l'oral et à l'écrit sur la tension entre droit coutumier et droit moderne dans la résolution des conflits au Cameroun.
\end{itemize}
\end{objectifs}

\begin{prerequis}
\begin{itemize}
\item Notion d'État de droit et de séparation des pouvoirs (programme de 1ère / Éducation à la citoyenneté).
\item Distinction de base entre morale, religion et droit (notion de norme).
\item Connaissance des institutions judiciaires camerounaises (Tribunaux, Cour Suprême, Conseil Constitutionnel, Chefferies traditionnelles).
\item Méthode de base de la dissertation et de l'explication de texte (acquise en 1ère).
\item Repères historiques : Code Napoléon, Déclaration 1789, Charte africaine des droits de l'homme.
\end{itemize}
\end{prerequis}

\newpage

\coursec{Définitions et distinctions conceptuelles : Droit, Justice, Loi, Morale}

\courssub{Situation d'accroche et problématique}

Dans un quartier populaire de Douala, M. Ngando voit sa maison démolie sur ordre du maire pour \og utilité publique \fg{}, sans indemnisation préalable ni procès équitable ; il saisit le tribunal administratif mais attend l'audience depuis trois ans. Parallèlement, un notable reconnu coupable de détournement de fonds publics bénéficie d'une liberté conditionnelle après deux mois de détention. Face à ces deux situations, les élèves s'interrogent : \textbf{la loi appliquée strictement suffit-elle à rendre justice ?} Qu'est-ce qui distingue le \og droit \fg{} de la \og justice \fg{} dans notre réalité camerounaise ? Cette tension entre la \cle{légalité} (conformité à la loi écrite) et la \cle{légitimité} (acceptabilité morale) structure toute la philosophie du droit. Pour y répondre, il faut d'abord disséquer les concepts avec rigueur.

\courssub{Étymologie : les racines latines du vocabulaire juridique}

\begin{definition}[Étymologie fondamentale]
\begin{itemize}
    \item \textbf{Jus, Juris} : Le lien, l'obligation, ce qui lie les hommes entre eux. C'est la racine de \emph{Justitia} (justice) et de \emph{Jurisprudence} (science du droit).
    \item \textbf{Justitia} : La conformité au \emph{Jus}, l'art du bon et de l'équitable (\emph{ars boni et aequi} chez le juriste romain Ulpien).
    \item \textbf{Lex, Legis} : La loi écrite, la règle posée par l'autorité législative (comices, puis Parlement). Elle est \emph{posée} (d'où \emph{droit positif}).
    \item \textbf{Directum} (bas latin) : Ce qui est droit, ce qui ne se courbe pas. Donne \og Droit \fg{} (français), \emph{Derecho} (espagnol), \emph{Right} (anglais). Il exprime la rectitude, la norme qui redresse la conduite.
\end{itemize}
\end{definition}

\begin{aretenir}[Distinction capitale]
\emph{Jus} = le droit subjectif (le pouvoir reconnu à une personne : \og J'ai le droit de...\fg{}). \emph{Lex} = la loi objective (la règle générale : \og La loi dit que...\fg{}). La justice vise à accorder le \emph{Jus} à chacun selon l'équité ; la loi édicte le \emph{Lex} pour tous.
\end{aretenir}

\courssub{Le Droit : définition, caractères et structure}

\begin{definition}[Le Droit (sens objectif)]
Ensemble de \textbf{règles de conduite} obligatoires, \textbf{générales} et \textbf{abstraites}, sanctionnées par la \textbf{force publique} (État), qui régissent les rapports sociaux. (Inspiré de Hans Kelsen, \emph{Théorie pure du droit}).
\end{definition}

\begin{propriete}[Les quatre caractères de la règle de droit]
\begin{enumerate}
    \item \textbf{Généralité} : S'adresse à une catégorie indéterminée de personnes (ex: \og Tout citoyen a le droit de voter \fg{}), non à un individu nommé.
    \item \textbf{Abstraction} : S'applique à des situations types, répétitives (ex: \og Le contrat se forme par l'échange de consentements \fg{}), non à un cas unique.
    \item \textbf{Impersonnalité} : Ne vise personne en particulier ; elle est anonyme.
    \item \textbf{Coercibilité} : Sanctionnée par la contrainte étatique (police, prison, amende, exécution forcée). C'est la marque distinctive par rapport à la morale ou la politesse.
\end{enumerate}
\end{propriete}

\begin{exemplebox}[Droit subjectif vs Droit objectif]
Mme Mvondo signe un contrat de bail. Le \textbf{Droit objectif} est l'ensemble des articles du Code civil camerounais sur le bail (règles générales). Le \textbf{Droit subjectif} de Mme Mvondo est son pouvoir exclusif d'exiger la jouissance paisible du logement (prérogative personnelle née du contrat).
\end{exemplebox}

\begin{methode}[Définir un concept philosophique (Genre + Différence spécifique)]
Pour définir \og Droit \fg{} au Bac :
1. \textbf{Genre prochain} : Ensemble de règles sociales / normes de conduite.
2. \textbf{Différence spécifique} : Sanctionnées par la force publique (coercibilité étatique), générales et abstraites.
3. \textbf{Exemple concret} : Code pénal camerounais (Art. 277 : le vol est puni de 5 à 10 ans de prison).
4. \textbf{Contre-exemple/Limite} : La règle morale \og Tu ne mentiras pas \fg{} (sanction = remords, pas prison) ou la règle de politesse \og Tenir la porte \fg{} (pas de sanction étatique).
\end{methode}

\courssub{La Justice : vertu, institution et idéal}

\begin{definition}[Justice (Aristote, \emph{Éthique à Nicomaque}, Livre V)]
Vertu complète consistant à \textbf{donner à chacun ce qui lui est dû} (\emph{suum cuique tribuere}). C'est à la fois une \textbf{vertu particulière} (respect des lois, équité dans les échanges) et la \textbf{vertu générale} (somme de toutes les vertus sociales).
\end{definition}

\begin{aretenir}[Double acception du mot \og Justice \fg{}]
\begin{itemize}
    \item \textbf{Justice comme Idéal / Valeur} : L'équité, la juste mesure, le but vers lequel tend le droit. \og La justice est la fin du droit \fg{} (Kant).
    \item \textbf{Justice comme Institution / Appareil} : L'ensemble des organes juridictionnels (Tribunaux, Cours, Parquets) et des auxiliaires (avocats, huissiers) chargés de \emph{dire le droit} (\emph{jus dicere}).
\end{itemize}
\end{aretenir}

\begin{attention}[Piège majeur : Confondre l'idéal et l'institution]
\begin{itemize}
    \item \textbf{Erreur} : \og Le juge a rendu justice \fg{} = \og Le juge a appliqué la loi \fg{}.
    \item \textbf{Correction} : Le juge \emph{dit le droit} (il applique la loi positive). Il ne \emph{fait} pas toujours la justice idéale (équité parfaite). Une décision peut être \textbf{légale} (conforme à la loi) mais \textbf{injuste} (ex: loi ségrégationniste appliquée strictement).
\end{itemize}
\end{attention}

\begin{definition}[Justice commutative (Aristote)]
\textbf{Égalité arithmétique} dans les échanges entre particuliers (contrats, vente, réparation de dommage). \emph{Exemple :} Si j'achète 10 kg de riz à 500 F/kg, je dois 5000 F. Ni plus, ni moins. Le juge rétablit l'égalité \emph{stricte} (proportion arithmétique).
\end{definition}

\begin{exemplebox}[Justice distributive vs Justice commutative]
\begin{itemize}
    \item \textbf{Distributive (Aristote)} : Répartition des honneurs, richesses, charges selon le \textbf{mérite} ou le \textbf{besoin} (proportion géométrique). \emph{Ex:} Bourses d'études sur critères sociaux ; impôt progressif.
    \item \textbf{Commutative} : Échanges contractuels, responsabilité délictuelle. \emph{Ex:} Indemnisation d'un accident de la route (réparation intégrale du préjudice).
\end{itemize}
\end{exemplebox}

\courssub{La Loi : source formelle, hiérarchie des normes (Kelsen)}

\begin{definition}[La Loi (sens formel)]
Règle générale, abstraite et impersonnelle, votée par le \textbf{Parlement} (Assemblée nationale + Sénat au Cameroun), promulguée par le Président de la République. Elle émane du \textbf{pouvoir législatif}.
\end{definition}

\begin{propriete}[Hiérarchie des normes (Hans Kelsen, Pyramide de Kelsen)]
Une norme tire sa validité de la norme supérieure. Au sommet, la \textbf{Grundnorm} (norme fondamentale : \og Il faut se conformer à la Constitution \fg{}).
\end{propriete}

\begin{popfigure}
\begin{tikzpicture}[scale=0.85, every node/.style={font=\small}]
    % Couleurs
    \colorlet{constcol}{poppurple}
    \colorlet{orgacol}{popblue}
    \colorlet{ordcol}{popgreen}
    \colorlet{decrcol}{poporange}
    \colorlet{arrc}{poppink}
    \colorlet{custcol}{popgold}
    % Niveaux
    \node[draw=constcol, thick, fill=constcol!15, rounded corners, minimum width=9cm, align=center] (const) at (0,4.5) {\textbf{CONSTITUTION DE 1996 (rév. 2008)}\\\textit{Norme fondamentale (Grundnorm camerounaise)}\\\textit{\textit{+ Préambule + Textes référencés (Bloc de constitutionnalité)}}};
    \node[draw=orgacol, thick, fill=orgacol!15, rounded corners, minimum width=8cm, align=center] (orga) at (0,3.1) {\textbf{LOIS ORGANIQUES}\\\textit{(Organisation des pouvoirs publics, statut de la magistrature, etc.)}\\\textit{Contrôle de constitutionnalité obligatoire}};
    \node[draw=ordcol, thick, fill=ordcol!15, rounded corners, minimum width=7cm, align=center] (ord) at (0,1.8) {\textbf{LOIS ORDINAIRES}\\\textit{(Code civil, Code pénal, Code du travail, Loi de finances, etc.)}};
    \node[draw=decrcol, thick, fill=decrcol!15, rounded corners, minimum width=6cm, align=center] (decr) at (0,0.5) {\textbf{DÉCRETS PRÉSIDENTIELS / DÉCRETS EN CONSEIL DES MINISTRES}\\\textit{(Application des lois, nominations, organisation administrative)}};
    \node[draw=popline, thick, fill=popcream, rounded corners, minimum width=5.5cm, align=center] (arrete) at (0,-0.7) {\textbf{ARRÊTÉS MINISTÉRIELS / PRÉFECTORAUX / MUNICIPAUX}\\\textit{(Mesures d'exécution locale, police administrative)}};
    \node[draw=custcol, thick, fill=custcol!20, rounded corners, minimum width=5cm, align=center] (coutume) at (0,-1.9) {\textbf{COUTUME / USAGE (Droit coutumier)}\\\textit{Art. 372 Code Civil : applicable si non contraire à l'ordre public}};
    % Flèches de validité
    \draw[->, thick, popdark] (const.south) -- (orga.north) node[midway, right, font=\tiny] {Fondement de validité};
    \draw[->, thick, popdark] (orga.south) -- (ord.north);
    \draw[->, thick, popdark] (ord.south) -- (decr.north);
    \draw[->, thick, popdark] (decr.south) -- (arrete.north);
    \draw[->, thick, popdark] (arrete.south) -- (coutume.north) node[midway, right, font=\tiny] {Subsidiarité / Contrôle constitutionnalité};
    % Annotation bloc de constitutionnalité
    \node[draw=popred, dashed, thick, fit=(const), inner sep=5pt, label={[popred, font=\tiny]above right:Bloc de constitutionnalité (Conseil Constitutionnel 2018)}] {};
\end{tikzpicture}
\legende{Pyramide de Kelsen adaptée au droit camerounais. La Constitution 1996 (révisée 2008) et son Préambule forment le Bloc de constitutionnalité au sommet. Les lois organiques lui sont soumises par contrôle obligatoire. La coutume (Art. 372 C. civ.) est une source autonome reconnue sous réserve de conformité à l'ordre public constitutionnel.}
\end{popfigure}

\courssub{Distinction fondamentale : Droit vs Morale}

\begin{center}
\begin{tabularx}{\linewidth}{|l|X|X|}
\hline
\rowcolor{popblueL}
\textbf{Critère} & \textbf{Droit} & \textbf{Morale} \\
\hline
\textbf{Origine (Source)} & Hétéronomie : Règle imposée de l'extérieur (État, Législateur). & Autonomie : Règle que le sujet se donne à lui-même (conscience, raison pratique). \\
\hline
\textbf{Portée (Intériorité)} & \textbf{Extériorité} : Juge les \emph{actes} extérieurs, observables. & \textbf{Intériorité} : Juge les \emph{intentions}, les motifs, la volonté. \\
\hline
\textbf{Sanction} & \textbf{Coercitive / Juridique} : Peine (prison, amende), contrainte physique, nullité. Exécutée par la force publique. & \textbf{Non-coercitive / Intérieure} : Remords, culpabilité, estime de soi, blâme social. Aucune force publique. \\
\hline
\textbf{Domaine} & Conduites sociales à incidence collective (sécurité, propriété, contrats). & Conduites totales de la vie (vertus, vices, devoirs envers soi-même). \\
\hline
\textbf{Obligation} & Juridique (obligation de \emph{faire} ou \emph{ne pas faire} sous peine de contrainte). & Morale (devoir-être, impératif catégorique kantien). \\
\hline
\end{tabularx}
\end{center}

\begin{savaistu}[Zone de chevauchement : Droits de l'homme et Bioéthique]
La frontière n'est pas étanche. Les \textbf{Droits de l'homme} (Déclaration 1789, Charte africaine 1981, Préambule Constitution Cameroun 1996) sont des \emph{revendications morales} érigées en \emph{normes juridiques suprêmes}.
\textbf{Exemple camerounais :} L'interdiction de la torture (Art. 5 Déclaration 1789, Art. 19 Constitution) est une exigence morale absolue devenue règle de droit intangible. En \textbf{bioéthique} (Loi 2002/003 sur la procréation médicalement assistée), le législateur traduit des consensus moraux (dignité de l'embryon, consentement libre) en règles juridiques sanctionnées.
\end{savaistu}

\courssub{Légalité vs Légitimité : le critère weberien}

\begin{definition}[Légalité]
Conformité formelle d'un acte ou d'une décision à la \textbf{loi positive} en vigueur. C'est un critère \textbf{juridique}, binaire (légal / illégal), vérifiable par le juge administratif ou constitutionnel.
\end{definition}

\begin{definition}[Légitimité (Max Weber, \emph{Économie et Société})]
Acceptabilité \textbf{morale}, \textbf{politique} ou \textbf{sociologique} du pouvoir ou de la norme. Elle repose sur la \textbf{croyance} en la validité de l'ordre. Trois types purs :
\begin{enumerate}
    \item \textbf{Rationnelle-légale} : Obéissance à la loi parce qu'elle a été édictée selon une procédure valide (État de droit moderne). \emph{Ex:} Président élu au suffrage universel.
    \item \textbf{Traditionnelle} : Obéissance aux \og coutumes éternelles \fg{} et aux détenteurs d'autorité héréditaire. \emph{Ex:} Chefferies traditionnelles au Cameroun (Lamidats, Fon).
    \item \textbf{Charismatique} : Obéissance à un chef porteur de \og grâce \fg{} personnelle, héros, prophète. Instable, routinisée ensuite en traditionnelle ou légale.
\end{enumerate}
\end{definition}

\begin{exoresolu}[Analyse : Démolition de la maison de M. Ngando (Situation d'accroche)]
\textbf{Énoncé :} Le maire signe un arrêté de démolition pour \og utilité publique \fg{} (base légale : Loi 85/09 du 4 juillet 1985 sur l'expropriation). Aucune indemnisation préalable. M. Ngando saisit le tribunal administratif (délai 3 ans).
\textbf{Analyse :}
\begin{itemize}
    \item \textbf{Légalité formelle :} L'arrêté municipal est \textbf{légal} (compétence du maire, motif d'utilité publique prévu par la loi, forme régulière).
    \item \textbf{Légitimité / Justice :} L'absence d'indemnisation \emph{préalable et juste} (Art. 16 Déclaration 1789, Art. 17 Constitution) rend l'acte \textbf{illégitime} et \textbf{injuste}, bien que légal. Le délai de 3 ans au tribunal viole le droit à un procès équitable dans un délai raisonnable (Art. 14 PIDCP ratifié par le Cameroun).
    \item \textbf{Conclusion :} Un acte peut être \textbf{parfaitement légal} (conforme à la procédure) et \textbf{radicalement injuste / illégitime} (atteinte disproportionnée au droit de propriété, déni de justice par lenteur).
\end{itemize}
\tcblower
\textbf{Solution détaillée :} La distinction permet de comprendre la revendication de M. Ngando : il ne conteste pas forcément le pouvoir d'exproprier (légalité), mais les \emph{modalités} (indemnisation, délai) qui violent la justice commutative (réparation) et la justice procédurale (accès au juge effectif). C'est le cœur du \emph{constitutionnalisme} : la légalité doit être soumise au contrôle de légitimité constitutionnelle (Droits fondamentaux).
\end{exoresolu}

\courssub{Le cas camerounais : Dualisme juridique (Bijuridisme)}

\begin{propriete}[Article 372 du Code Civil Camerounais (Loi 87-01 du 16 juillet 1987)]
\og Les coutumes sont applicables dans les matières où la loi n'a pas prévu de dispositions, pourvu qu'elles ne soient pas contraires à l'ordre public, aux bonnes mœurs et aux lois écrites. \fg{}
\end{propriete}

\begin{aretenir}[Le Cameroun, pays bijuridique unique en Afrique]
\begin{itemize}
    \item \textbf{8 Régions (Ex-Français) :} \textbf{Civil Law} (Droit écrit, codifié, juge d'instruction, procédure inquisitoire). Sources : Codes (Civil, Pénal, Travail, OHADA), Constitution.
    \item \textbf{2 Régions (Nord-Ouest / Sud-Ouest - Ex-Britannique) :} \textbf{Common Law} (Droit jurisprudentiel, \emph{stare decisis}, procédure accusatoire, \emph{Habeas Corpus}, Jury). Sources : \emph{Case Law}, Statutes, Equity.
    \item \textbf{Point de rencontre :} Cour Suprême (Chambre Judiciaire, Chambre Administrative, Chambre des Comptes) + \textbf{Conseil Constitutionnel} (opérationnel depuis 2018). Droit OHADA (droit des affaires uniforme) applicable sur tout le territoire.
\end{itemize}
\end{aretenir}

\begin{savaistu}[Disparités procédurales concrètes : \emph{Habeas Corpus} vs \emph{Référé-liberté}]
\begin{itemize}
    \item \textbf{Régions Anglophones (Common Law) :} \textbf{\emph{Habeas Corpus}} (Art. 404 Criminal Procedure Code). Recours \emph{ultra-rapide} (24h-48h) devant le \emph{High Court} pour contester une détention arbitraire. Le juge ordonne la production du corps du détenu. Charge de la preuve sur l'État.
    \item \textbf{Régions Francophones (Civil Law) :} \textbf{Référé-liberté} (Art. 809-1 CPC modifié, Loi 2011/027). Recours devant le Président du Tribunal de Grande Instance. Délai plus long, condition d'urgence + atteinte manifestement illégale à une liberté fondamentale. Charge de la preuve partagée.
    \item \textbf{Conséquence :} Inégalité devant la justice selon le lieu de résidence. Un justiciable à Bamenda a un remède plus efficace contre l'arbitraire qu'un justiciable à Yaoundé ou Douala pour le même fait.
    \item \textbf{Perspective :} La Commission d'Harmonisation du Droit (CHR) œuvre à réduire ces écarts, mais l'effectivité reste inégale.
\end{itemize}
\end{savaistu}

\begin{popfigure}
\begin{tikzpicture}[scale=0.9]
    % Diagramme de Venn 3 ensembles : Droit, Morale, Religion
    \def\R{3.5cm}
    \def\offset{2.2cm}
    \colorlet{droitc}{popblue}
    \colorlet{moralec}{popgreen}
    \colorlet{religc}{poppurple}
    % Cercles
    \begin{scope}[blend mode=multiply, opacity=0.7]
        \fill[droitc] (-\offset, 0) circle (\R);
        \fill[moralec] (\offset, 0) circle (\R);
        \fill[religc] (0, {-\offset*0.866}) circle (\R); % Hauteur triangle équilatéral
    \end{scope}
    % Contours
    \draw[thick, droitc] (-\offset, 0) circle (\R) node[above left=0.2cm and 0.2cm, font=\bfseries, droitc] {DROIT};
    \draw[thick, moralec] (\offset, 0) circle (\R) node[above right=0.2cm and 0.2cm, font=\bfseries, moralec] {MORALE};
    \draw[thick, religc] (0, {-\offset*0.866}) circle (\R) node[below=0.2cm, font=\bfseries, religc] {RELIGION};
    % Étiquettes intersections (positionnées manuellement pour lisibilité)
    % Centre triple intersection
    \node[align=center, font=\small, text=popdark] at (0, -0.2) {\textbf{Intersection Triple}\\\textit{Droits de l'homme, Bioéthique,}\\\textit{Interdit meurtre/vol, Mariage}};
    % Droit \cap Morale (haut centre)
    \node[align=center, font=\small, text=popdark] at (0, 1.8) {\textbf{Droit $\cap$ Morale}\\\textit{Respect contrats, Protection famille,}\\\textit{Responsabilité civile}};
    % Droit \cap Religion (gauche bas)
    \node[align=center, font=\small, text=popdark] at (-2.2, -1.5) {\textbf{Droit $\cap$ Religion}\\\textit{Statut personnel (mariage,}\\\textit{héritage) droit coutumier,}\\\textit{Fêtes légales (Noël, Aid)}};
    % Morale \cap Religion (droite bas)
    \node[align=center, font=\small, text=popdark] at (2.2, -1.5) {\textbf{Morale $\cap$ Religion}\\\textit{Charité, Vertus, Péchés non}\\\textit{pénalisés (orgueil, paresse)}};
    % Zones propres
    \node[align=center, font=\small, text=droitc] at (-3.8, 1.2) {\textbf{Propre au Droit}\\\textit{Code de la route,}\\\textit{Procédure civile,}\\\textit{Droit des sociétés OHADA}};
    \node[align=center, font=\small, text=moralec] at (3.8, 1.2) {\textbf{Propre à la Morale}\\\textit{Devoirs envers soi-même,}\\\textit{Vertus (courage, tempérance),}\\\textit{Promesse gratuite}};
    \node[align=center, font=\small, text=religc] at (0, -3.8) {\textbf{Propre à la Religion}\\\textit{Rites funéraires,}\\\textit{Dogmes, Sacrements,}\\\textit{Diète alimentaire (porc, jeûne)}};
\end{tikzpicture}
\legende{Diagramme de Venn : Relations entre Droit, Morale et Religion. L'intersection centrale (Droits fondamentaux, bioéthique) montre la juridicisation de l'éthique. Les zones propres illustrent l'autonomie relative de chaque ordre normatif.}
\end{popfigure}

\begin{popfigure}
\begin{tikzpicture}
    \begin{axis}[
        ybar,
        bar width=18pt,
        width=\linewidth,
        height=8cm,
        enlarge x limits=0.25,
        symbolic x coords={Tribunaux Étatiques, Chefferies/Notables, Médiation Familiale/AMI},
        xtick=data,
        x tick label style={rotate=30, anchor=east, font=\small},
        ylabel={\% de justiciables (estimation)},
        ymin=0, ymax=100,
        ymajorgrids=true,
        grid style=dashed,
        legend style={at={(0.5,-0.25)}, anchor=north, legend columns=-1, font=\small},
        nodes near coords={\pgfmathprintnumber\pgfplotspointmeta\%},
        nodes near coords align={vertical},
        every node near coord/.append style={font=\small, popdark},
        popbar1/.style={fill=popblue!80},
        popbar2/.style={fill=poporange!80},
        popbar3/.style={fill=popgreen!80},
    ]
    \addplot[popbar1] coordinates {(Tribunaux Étatiques, 35) (Chefferies/Notables, 10) (Médiation Familiale/AMI, 5)};
    \addplot[popbar2] coordinates {(Tribunaux Étatiques, 15) (Chefferies/Notables, 45) (Médiation Familiale/AMI, 10)};
    \addplot[popbar3] coordinates {(Tribunaux Étatiques, 5) (Chefferies/Notables, 5) (Médiation Familiale/AMI, 60)};
    \legend{Contentieux Pénal, Contentieux Foncier/Coutumier, Contentieux Familial/Quotidien}
    \end{axis}
    % Source note
    \node[anchor=north west, font=\tiny, text=popmuted] at (0.02\linewidth, -0.15\linewidth) {Source : Inspiré World Bank \og Justice for the Poor \fg{} Cameroun \& Études MINJUSTICE 2020-2023. Données illustratives.};
\end{tikzpicture}
\legende{Sources du droit effectivement mobilisé par les justiciables au Cameroun. Le pluralisme juridique est une réalité vécue : l'État n'a pas le monopole effectif du règlement des conflits. La médiation familiale/AMI (Mécanismes Alternatifs de Règlement des Différends) domine le contentieux familial ; la chefferie domine le foncier coutumier.}
\end{popfigure}

\courssub{Synthèse et ouverture vers la Section 2}

\begin{aretenir}[Bilan conceptuel pour la dissertation]
\begin{enumerate}
    \item \textbf{Droit} = Norme coercitive, étatique, générale (Légalité).
    \item \textbf{Justice} = Idéal d'équité (valeur) + Institution de jugement (Double sens). But du droit.
    \item \textbf{Loi} = Source formelle majeure, votée par le Parlement. Hiérarchisée (Kelsen).
    \item \textbf{Morale} = Intériorité, autonomie, sanction intérieure. Source d'inspiration du droit (Droits de l'homme).
    \item \textbf{Légalité $\neq$ Légitimité} : Un régime peut être légal (Hitler élu, lois votées) et illégitime (injuste). Le constitutionnalisme vise à soumettre la légalité à la légitimité constitutionnelle (Droits fondamentaux).
    \item \textbf{Cameroun} : Bijuridisme (Civil Law / Common Law) + Pluralisme vivant (Coutume Art. 372, Chefferies, Médiation). Défis : Effectivité des droits, délais, accès au juge, harmonisation, \textbf{Conseil Constitutionnel}.
\end{enumerate}
\end{aretenir}

\begin{methode}[Transition vers Section 2 : Les Fondements]
Maintenant que les concepts sont posés, nous devons demander : \textbf{D'où vient la force obligatoire du droit ?} Pourquoi obéit-on ? Par peur de la sanction (Positivisme / Kelsen / Hart) ou par conviction morale que la loi est juste (Jusnaturalisme / Droit naturel / Finnis) ? C'est l'objet de la Section 2 : \emph{Fondements du Droit : Droit positif vs Justice (Droit naturel)}.
\end{methode}

\begin{exercice}[Entraînement Bac - Type Dissertation]
\textbf{Sujet :} \og La loi est-elle la justice ? \fg{}
\textbf{Consignes :} 1. Analyser les termes (définitions, distinctions). 2. Problématiser (opposition légalité/légitimité, droit positif/Justice idéale). 3. Construire un plan dialectique (Thèse : La loi incarne la justice / Antithèse : La loi peut être injuste / Synthèse : Le contrôle de constitutionnalité et l'équité comme correctifs). 4. Illustrer par le droit camerounais (Art. 372, Bijuridisme, Conseil Constitutionnel).
\end{exercice}

\begin{corrige}[Exercice n°1 - Éléments de corrigé]
\textbf{Introduction :} Accroche sur M. Ngando (légal mais injuste). Définitions : Loi = règle générale votée ; Justice = idéal d'équité (suum cuique). Problématique : La conformité à la loi (légalité) suffit-elle à réaliser l'équité (justice) ? Plan : I. La loi comme tentative d'objectivation de la justice (sécurité juridique, généralité). II. Les limites : loi injuste, rigueur de la loi, lacunes (cas camerounais : coutume vs ordre public, bijuridisme). III. Les correctifs : Contrôle de constitutionnalité (\textbf{Conseil Constitutionnel} 2018), Équité (pouvoir du juge), Droit naturel comme critère de validité morale.
\end{corrige}

\coursec{Les fondements du Droit : Droit positif vs Justice (Droit naturel)}

Dans la section précédente, nous avons distingué le \cle{Droit} (ensemble de règles sanctionnées par l'autorité publique), la \cle{Justice} (idéal de rendre à chacun son dû), la \cle{Loi} (règle écrite votée par le législateur) et la \cle{Morale} (règles intérieures de la conscience). Mais une question demeure en suspens : d'où le Droit tire-t-il son autorité ? Qu'est-ce qui fait qu'une règle est \emph{valide} et s'impose à nous ? Deux grandes réponses s'affrontent depuis des siècles : pour les uns, le Droit se fonde sur lui-même, sur sa propre organisation formelle (c'est le \cle{positivisme juridique}) ; pour les autres, le Droit ne vaut que s'il est conforme à des principes supérieurs de justice (c'est le \cle{jusnaturalisme}). Ce débat n'est pas théorique : il décide de savoir si l'on doit obéir à une loi injuste.

\courssub{Le positivisme juridique : le Droit se fonde sur lui-même}

\textbf{Intuition.} Imaginez une partie de football à Yaoundé. Ce qui rend une règle du jeu valable, ce n'est pas qu'elle soit « morale », c'est qu'elle figure dans le règlement officiel adopté par la fédération et appliqué par l'arbitre. De même, dit le positiviste, une loi est du droit non pas parce qu'elle est juste, mais parce qu'elle a été \emph{produite selon la procédure correcte} par l'autorité compétente.

\begin{definition}[Positivisme juridique]
Le \cle{positivisme juridique} est la théorie selon laquelle le droit est un ensemble de normes \emph{posées} (en latin \emph{positum}) par les autorités humaines compétentes, dont la validité dépend uniquement de leur conformité aux procédures de production juridique, et non de leur contenu moral. Principe fondamental : \textbf{séparation stricte du Droit et de la Morale}. Une loi injuste reste une loi valide tant qu'elle n'est pas abrogée.
\end{definition}

\textbf{Hans Kelsen (1881-1973)} est le grand représentant du positivisme. Dans sa \emph{Théorie pure du droit} (1934), il veut faire du droit une science « pure », débarrassée de la morale, de la politique et de la religion. Pour lui, les normes juridiques forment une \emph{pyramide hiérarchique} : chaque norme tire sa validité d'une norme supérieure (le règlement du décret, le décret de la loi, la loi de la Constitution). Mais alors, d'où la Constitution elle-même tire-t-elle sa validité ?

\begin{definition}[La \emph{Grundnorm} ou norme fondamentale]
Chez Kelsen, la \cle{Grundnorm} (norme fondamentale) est une \emph{hypothèse transcendantale logique}, c'est-à-dire un présupposé nécessaire de la pensée juridique : « il faut se conduire comme la Constitution le prescrit ». Elle valide la Constitution sans être elle-même une norme juridique positive. Elle est le point d'arrêt ultime de la chaîne de validité, comme le premier maillon imaginaire qui tient toute la chaîne.
\end{definition}

\textbf{H.L.A. Hart (1907-1992)}, philosophe anglais, propose dans \emph{Le Concept de droit} (1961) une version sociologique du positivisme. Pour lui, un système juridique repose sur une \cle{règle de reconnaissance} : une pratique sociale partagée par les juges et les fonctionnaires, qui permet d'identifier ce qui compte comme loi (au Cameroun : « ce qui est voté par l'Assemblée nationale, promulgué et publié au Journal officiel »). Là encore, la \textbf{validité n'est pas la moralité} : une norme peut être juridiquement valide et moralement condamnable.

\begin{attention}[Piège fréquent]
Ne pas écrire : « le positivisme affirme que toutes les lois sont justes ». C'est l'inverse ! Le positiviste dit que la question de la justice d'une loi est une question \emph{morale}, extérieure à la question de sa \emph{validité juridique}. Kelsen lui-même reconnaît qu'une loi peut être valide et injuste. Confondre validité et moralité est l'erreur classique de la dissertation.
\end{attention}

\courssub{Le jusnaturalisme : un droit supérieur au droit des hommes}

\textbf{Intuition.} Lors des procès de Nuremberg (1945-1946), des criminels nazis se défendaient en disant : « nous avons obéi aux lois de notre pays ». Les juges ont répondu qu'il existe des crimes qu'aucune loi ne peut légitimer. Cela suppose qu'au-dessus des lois positives, il existe un \emph{droit naturel}, inscrit dans la raison humaine, auquel toute loi doit se conformer.

\begin{definition}[Jusnaturalisme ou droit naturel]
Le \cle{jusnaturalisme} est la théorie selon laquelle il existe des principes de justice \emph{universels, immuables et supérieurs} au droit positif, fondés dans la nature humaine, la raison ou Dieu. Une loi positive qui les contredit gravement perd sa qualité de loi : \emph{Lex injusta non est lex} (« une loi injuste n'est pas une loi », formule reprise de saint Augustin par Thomas d'Aquin).
\end{definition}

\textbf{Thomas d'Aquin (1225-1274)} donne au droit naturel sa formulation classique dans la \emph{Somme théologique} : la loi naturelle est la participation de la créature raisonnable à la loi éternelle de Dieu ; elle commande de « faire le bien et éviter le mal ». À l'époque moderne, \textbf{John Finnis} (\emph{Droit naturel et droits naturels}, 1980) fonde le droit naturel sur des biens humains fondamentaux évidents (la vie, la connaissance, l'amitié...) que toute loi juste doit respecter. \textbf{Lon Fuller}, quant à lui, défend une « morale interne du droit » : pour être du droit, des règles doivent être générales, publiées, non rétroactives, compréhensibles, stables. Un système de lois secrètes et rétroactives (comme sous le nazisme) n'est pas un vrai système juridique.

\begin{aretenir}[Citation fondatrice]
« Une loi injuste est une espèce de violence. » — \textbf{Saint Augustin}, \emph{Du libre arbitre} (\emph{De libero arbitrio}). Cette formule est la base du jusnaturalisme classique : la force ne fait pas le droit, et une norme qui opprime ne mérite pas le nom de loi.
\end{aretenir}

\courssub{Le débat Hart-Fuller (1958) : faut-il obéir à une loi injuste ?}

En 1958, la \emph{Harvard Law Review} publie un échange célèbre entre Hart et Fuller à propos des « lois iniques » : les lois nazies, les lois de l'apartheid sud-africain, les décrets d'exception. \textbf{Hart} (positiviste) soutient que ces lois étaient juridiquement valides, quoique moralement monstrueuses ; dire « c'est du droit, mais un droit mauvais » permet de garder une conscience morale claire. \textbf{Fuller} répond qu'un tel système, par son injustice radicale et ses procédures corrompues, cesse d'être du droit.

Face à la loi injuste, une troisième voie existe : la \cle{désobéissance civile}. \textbf{Henry David Thoreau} (\emph{La Désobéissance civile}, 1849) refuse de payer un impôt finançant l'esclavage et la guerre : « sous un gouvernement qui emprisonne injustement, la place de l'homme juste est aussi en prison ». \textbf{Gandhi} en fait une arme de lutte non violente contre la colonisation britannique en Inde. \textbf{Martin Luther King} (\emph{Lettre de la prison de Birmingham}, 1963) la justifie ainsi : « une loi injuste est un code qui n'est pas en harmonie avec la loi morale » ; on la transgresse ouvertement, en acceptant la sanction, pour réveiller la conscience de la communauté.

\begin{exemplebox}[Trois attitudes face à la loi injuste]
\begin{enumerate}
\item \textbf{Obéissance positiviste} : j'applique la loi car elle est valide, tout en la jugeant immorale (position de Hart).
\item \textbf{Désobéissance civile} : je transgresse publiquement et pacifiquement la loi injuste, en acceptant la sanction (Thoreau, Gandhi, King).
\item \textbf{Dénonciation jusnaturaliste} : je déclare que cette norme n'est pas du droit du tout (Fuller, Finnis).
\end{enumerate}
\end{exemplebox}

\courssub{La théorie du « Droit vivant » (Ehrlich)}

\textbf{Eugen Ehrlich} (1862-1922), juriste autrichien, ouvre une perspective sociologique : le droit réellement efficace n'est pas seulement celui des codes et des tribunaux, c'est le \cle{droit vivant} (\emph{living law}), c'est-à-dire l'ensemble des règles que les gens \emph{pratiquent effectivement} dans leur vie sociale : coutumes villageoises, usages des marchands, règles des associations, conventions familiales. Le droit de l'État n'est qu'une petite partie du droit réel.

\begin{exemplebox}[Le droit vivant au quotidien camerounais]
À Mokolo ou à Foumban, la plupart des ventes de terrains, des mariages et des successions se règlent selon la coutume, devant le chef ou les notables, sans jamais passer devant un tribunal d'État. Pour Ehrlich, ce droit coutumier \emph{est} du droit à part entière, car il organise réellement la vie sociale. Le droit écrit qui ignore ces pratiques risque de rester lettre morte.
\end{exemplebox}

\begin{savaistu}[Le droit vivant et la sociologie juridique]
Ehrlich, fondateur de la \emph{sociologie juridique}, invite les juristes à sortir des codes pour observer le terrain. Au Cameroun, les travaux de l'anthropologue \textbf{Francis Nyamnjoh} ou du juriste \textbf{Emmanuel Bindzi} montrent comment le « droit vivant » (coutumes, arrangements à l'amiable, médiation des chefs) résout 80\% des conflits sans saisir les tribunaux. Cette approche éclaire l'effectivité du droit bien mieux que la seule lecture des textes.
\end{savaistu}

\courssub{Application au Cameroun : Constitution, contrôle des lois et pluralisme juridique}

\textbf{La Constitution du 18 janvier 1996} est la norme suprême camerounaise. Son \textbf{Préambule} — qui a valeur constitutionnelle — intègre solennellement les droits de l'homme (Déclaration universelle de 1948, Charte africaine de 1981) et les devoirs du citoyen : « le peuple camerounais affirme son attachement aux libertés fondamentales inscrites dans la Déclaration universelle des droits de l'homme ». Le \textbf{Conseil constitutionnel} contrôle la conformité des lois à la Constitution. Toutefois, l'effectivité de la justice reste fragile : le \textbf{Conseil supérieur de la magistrature} (CSM), qui gère la carrière des magistrats, est présidé par le Président de la République, ce qui traduit un poids important de l'exécutif sur le judiciaire et nourrit le débat sur l'indépendance de la justice.

\begin{propriete}[Monisme et dualisme en droit international]
\textbf{Monisme} (défendu par Kelsen) : le droit international et le droit interne forment \emph{un seul ordre juridique}, avec primauté du droit international ; les traités s'appliquent directement. \textbf{Dualisme} : les deux ordres sont distincts ; un traité ne vaut en droit interne qu'après une loi de transposition. \textbf{Le Cameroun est moniste modéré} : selon l'article 45 de la Constitution de 1996, « les traités ou accords internationaux régulièrement approuvés ou ratifiés ont, dès leur publication, une autorité supérieure à celle des lois ».
\end{propriete}

Le Cameroun connaît en outre un \cle{pluralisme juridique} : coexistent le droit étatique (Constitution, codes), le droit coutumier (successions, mariages, foncier), le droit religieux (droit musulman appliqué par les tribunaux \emph{alkali} dans le Nord) et le droit international. Cette coexistence est théoriquement ordonnée par la hiérarchie des normes, mais pratiquement conflictuelle.

\begin{methode}[Analyser un conflit de normes]
\begin{enumerate}
\item \textbf{Identifier} les normes en conflit (ex. : coutume successorale vs Code civil).
\item \textbf{Vérifier la hiérarchie} : Constitution $>$ Traité ratifié (art. 45) $>$ Loi $>$ Règlement $>$ Coutume.
\item \textbf{Chercher la clause de sauvegarde} : une coutume n'est recevable que si elle n'est pas contraire à l'ordre public, à la loi et aux droits fondamentaux.
\item \textbf{Conclure} sur la norme applicable et justifier la réponse.
\end{enumerate}
\end{methode}

\begin{exoresolu}[L'héritage de la fille : coutume contre Constitution]
\textbf{Énoncé.} Dans certaines communautés de l'Ouest et du Littoral, la coutume exclut les filles de la succession foncière du père, au profit des seuls héritiers mâles. Une jeune fille de Bafoussam, privée de sa part d'héritage, saisit le juge. Quelle norme doit s'appliquer ? Justifiez votre réponse en mobilisant les notions du cours.
\tcblower
\textbf{Solution détaillée.}
\begin{enumerate}
\item \emph{Normes en conflit} : la coutume successorale (droit vivant, au sens d'Ehrlich) contre le Code civil (art. 745 et suivants, qui ne distingue pas les héritiers selon le sexe) et l'article 1er de la Constitution (égalité de tous devant la loi).
\item \emph{Hiérarchie} : la Constitution est suprême ; la coutume est une source subordonnée, admise seulement à titre supplétif.
\item \emph{Clause de sauvegarde} : la coutume ne peut s'appliquer si elle viole l'égalité, droit fondamental. Dans son \textbf{avis consultatif n° 001/2017}, la Cour suprême du Cameroun a jugé discriminatoires les coutumes excluant les filles de la succession.
\item \emph{Conclusion} : la norme applicable est le droit étatique conforme à la Constitution ; la fille hérite au même titre que le fils. Ce cas illustre la tension entre l'article 1er de la Constitution (égalité) et le respect des traditions, et montre que la coutume, source du droit, reste soumise au contrôle de justice.
\end{enumerate}
\end{exoresolu}

\begin{attention}[Ne pas confondre droit naturel et droit coutumier]
Le \textbf{droit naturel} relève de la \emph{philosophie} : il prétend à l'\emph{universalité} (valable pour tous les hommes, partout). Le \textbf{droit coutumier} relève de la \emph{sociologie} : il est \emph{particulier} à une communauté locale. Surtout, la coutume n'est pas automatiquement juste : l'excision ou l'exclusion successorale des filles sont des coutumes, et elles violent la dignité humaine. Le droit naturel peut servir à \emph{critiquer} une coutume injuste.
\end{attention}

\begin{exemplebox}[Un chiffre qui fait réfléchir]
Selon le rapport 2023 du MINJUSTICE, environ \textbf{65 \%} des litiges fonciers ruraux sont réglés par les chefferies (droit coutumier) sans saisine des tribunaux, et seuls \textbf{12 \%} aboutissent à un titre foncier sécurisé. Le droit vivant domine donc la pratique, mais au prix d'une insécurité juridique pour les populations.
\end{exemplebox}

\begin{popfigure}
\begin{center}
\begin{tikzpicture}[
  font=\small,
  racine/.style={draw=popdark, fill=popcream, rounded corners=2pt, text width=2.6cm, align=center, minimum height=0.9cm},
  tronc/.style={draw=popblue, fill=popblueL, rounded corners=2pt, text width=3.2cm, align=center, font=\bfseries},
  branche/.style={draw=poppurple, fill=poppurpleL, rounded corners=2pt, text width=3.4cm, align=center},
  fruit/.style={draw=popgreen, fill=popgreenL, rounded corners=2pt, text width=3.4cm, align=center},
  lien/.style={-{Stealth[length=2mm]}, thick, popmuted}]
% Racines
\node[racine] (r1) at (-4,0) {Nature, Dieu, Raison humaine};
\node[racine] (r2) at (4,0) {Convention, Force, Efficacité sociale};
% Troncs
\node[tronc] (t1) at (-4,1.6) {Jusnaturalisme};
\node[tronc] (t2) at (4,1.6) {Positivisme};
% Branches
\node[branche] (b1) at (-4,3.2) {Augustin, Thomas d'Aquin, Finnis, Fuller};
\node[branche] (b2) at (4,3.2) {Kelsen (\emph{Grundnorm}), Hart (règle de reconnaissance)};
% Fruits
\node[fruit] (f1) at (-4,4.8) {On peut désobéir à la loi injuste ; contrôle moral des lois};
\node[fruit] (f2) at (4,4.8) {Sécurité juridique ; la loi valide s'applique même si elle est critiquable};
\draw[lien] (r1) -- (t1);
\draw[lien] (r2) -- (t2);
\draw[lien] (t1) -- (b1);
\draw[lien] (t2) -- (b2);
\draw[lien] (b1) -- (f1);
\draw[lien] (b2) -- (f2);
\end{tikzpicture}
\end{center}
\legende{Arbre des fondements du droit : des racines (fondements ultimes) aux troncs (les deux grandes théories), aux branches (auteurs clés) et aux fruits (conséquences concrètes pour le justiciable).}
\end{popfigure}

\begin{popfigure}
\begin{center}
\begin{tikzpicture}[font=\small]
\draw[-{Stealth}, very thick, popdark] (0,0) -- (13.5,0);
\foreach \x/\an/\txt in {0.5/1960/{Indépendance, 1ère Constitution}, 3.2/1961/{Réunification, République fédérale}, 5.9/1972/{État unitaire (référendum du 20 mai)}, 8.6/1996/{Constitution : droits de l'homme, décentralisation}, 11.6/2008/{Révision : suppression de la limite des mandats}}{
  \fill[poporange] (\x,0) circle (2.5pt);
  \draw[poporange, thick] (\x,0) -- (\x,0.5);
  \node[above, align=center, text width=2.3cm] at (\x,0.6) {\textbf{\an}\\ \txt};
}
\end{tikzpicture}
\end{center}
\legende{Grandes dates du constitutionnalisme camerounais : de l'indépendance (1960) à la révision de 2008, en passant par la réunification (1961), l'État unitaire (1972) et la Constitution de 1996.}
\end{popfigure}

\begin{popfigure}
\begin{center}
\begin{tikzpicture}[font=\small]
% Contour très simplifié du Cameroun
\draw[thick, popdark, fill=popcream] (2.2,6.4) -- (3.4,5.6) -- (4.4,4.9) -- (5.6,4.6) -- (6.6,3.9) -- (6.4,2.6) -- (5.4,1.6) -- (4.2,0.4) -- (2.6,0.2) -- (1.4,0.9) -- (0.8,2.2) -- (1.0,3.6) -- (1.6,4.6) -- (1.2,5.4) -- cycle;
% Zones
\fill[popgoldL, opacity=0.75] (2.2,6.4) -- (3.4,5.6) -- (4.4,4.9) -- (3.6,4.4) -- (1.6,4.6) -- (1.2,5.4) -- cycle;
\fill[popgreenL, opacity=0.75] (1.6,4.6) -- (3.6,4.4) -- (4.4,4.9) -- (5.6,4.6) -- (6.6,3.9) -- (6.4,2.6) -- (5.4,1.6) -- (4.2,0.4) -- (2.6,0.2) -- (1.4,0.9) -- (0.8,2.2) -- (1.0,3.6) -- cycle;
% Villes
\fill[popink] (1.6,1.3) circle (2.2pt) node[below left=1pt, popink]{\textbf{Douala}};
\fill[popink] (2.7,1.7) circle (2.2pt) node[below right=1pt, popink]{\textbf{Yaoundé}};
\fill[popink] (2.9,5.6) circle (2.2pt) node[above right=1pt, popink]{\textbf{Garoua}};
\fill[popink] (1.5,3.4) circle (2pt) node[left=1pt, popink]{\textbf{Bafoussam}};
% Nord
\draw[-{Stealth}, thick, popdark] (5.6,6.9) -- (5.6,6.2) node[midway, right]{N};
% Légende
\node[anchor=west, draw=popmuted, fill=white, align=left, text width=5.2cm] at (7.2,4.6) {%
\textcolor{popgold}{\rule{0.35cm}{0.35cm}}~Nord : droit musulman et coutumier fort (tribunaux \emph{alkali})\\[2pt]
\textcolor{popgreen}{\rule{0.35cm}{0.35cm}}~Centre-Sud-Est-Ouest : droit coutumier foncier et successoral\\[2pt]
\textcolor{popink}{$\bullet$}~Grandes villes : droit étatique moderne dominant};
\end{tikzpicture}
\end{center}
\legende{Carte schématique (approximative) des zones d'influence des droits au Cameroun : droit musulman et coutumier au Nord, droit coutumier foncier au Centre-Sud, droit moderne dominant dans les grandes villes.}
\end{popfigure}

\begin{aretenir}[L'essentiel de la section]
\begin{itemize}
\item Le \textbf{positivisme} (Kelsen, Hart) : la validité d'une norme dépend de sa production conforme, non de sa moralité ; sommet de la pyramide : la \emph{Grundnorm} ou la règle de reconnaissance.
\item Le \textbf{jusnaturalisme} (Augustin, Thomas d'Aquin, Finnis, Fuller) : une loi radicalement injuste n'est pas une loi ; il existe un droit supérieur fondé sur la raison et la dignité humaine.
\item Face à la loi inique : obéir (Hart), désobéir civilement (Thoreau, Gandhi, King) ou nier sa qualité de droit (Fuller).
\item Le \textbf{droit vivant} (Ehrlich) : le droit réel inclut les coutumes et usages effectivement pratiqués.
\item Au Cameroun : Constitution de 1996, monisme modéré (art. 45), pluralisme juridique ordonné par la hiérarchie des normes et le contrôle de conventionalité et de constitutionnalité.
\end{itemize}
\end{aretenir}

\begin{exercice}[Pour s'entraîner]
\textbf{Exercice 1.} Définissez la \emph{Grundnorm} et expliquez pourquoi Kelsen la qualifie d'« hypothèse » et non de norme positive.

\textbf{Exercice 2.} Un décret pris lors d'un état d'exception prive un groupe de citoyens de leurs droits sans procès. Analysez cette situation successivement du point de vue de Hart, puis de Fuller, puis de Martin Luther King.

\textbf{Exercice 3 (sujet de dissertation).} « Faut-il obéir à une loi injuste ? » Rédigez l'introduction et le plan détaillé de la dissertation, en mobilisant le débat Hart-Fuller et le cas camerounais du pluralisme juridique.
\end{exercice}

\begin{corrige}[Exercice 1]
\textbf{La \emph{Grundnorm} chez Kelsen.} La \emph{Grundnorm} (norme fondamentale) est le présupposé logique ultime qui valide la Constitution sans être elle-même une norme positive. Elle n'est pas écrite, pas votée, pas promulguée : c'est une \emph{hypothèse transcendantale} (« il faut se conduire comme la Constitution le prescrit ») qui permet d'arrêter la régression à l'infini de la chaîne de validité (loi $\leftarrow$ Constitution $\leftarrow$ ?). Kelsen la qualifie d'hypothèse car elle relève de la \emph{théorie} du droit (science), non de la pratique juridique positive ; elle est le point de vue du juriste qui décrit le système, non une norme produite par le législateur.
\end{corrige}

\begin{corrige}[Exercice 2]
\textbf{Analyse croisée d'un décret d'exception liberticide.}
\begin{itemize}
\item \textbf{Hart (positiviste) :} Le décret est \emph{juridiquement valide} s'il a été pris selon la procédure constitutionnelle (compétence, forme, publication). Sa moralité est une question distincte : on peut dire « c'est la loi, mais c'est une loi mauvaise ». Le juge doit l'appliquer ; le citoyen a une obligation légale d'obéissance, mais pas nécessairement une obligation morale.
\item \textbf{Fuller (jusnaturaliste procédural) :} Un décret qui prive de droits sans procès viole la \emph{morale interne du droit} (généralité, publicité, non-rétroactivité, congruence entre loi officielle et action administrative). Un tel système, par son injustice radicale, \emph{cesse d'être du droit}. Le juge ne devrait pas l'appliquer ; le citoyen n'a pas d'obligation juridique de l'obéir.
\item \textbf{Martin Luther King (désobéissance civile) :} Ce décret est une \emph{loi injuste} car il n'est pas en harmonie avec la loi morale (dignité, égalité). Le citoyen doit le transgresser \emph{ouvertement, pacifiquement, en acceptant la sanction} (prison), pour dramatiser l'injustice et forcer la majorité à réformer la loi. La désobéissance n'est pas anarchie : c'est une fidélité supérieure à la justice.
\end{itemize}
\end{corrige}

\begin{corrige}[Exercice 3]
\textbf{Sujet : « Faut-il obéir à une loi injuste ? »}
\textbf{Introduction.}
\begin{itemize}
\item \textbf{Accroche :} « Une loi injuste n'est pas une loi » (saint Augustin) — mais les procès de Nuremberg ont montré qu'on jugeait des hommes pour avoir \emph{obéi} à la loi.
\item \textbf{Définition des termes :} « Obéir » = se conformer à une norme contraignante ; « loi injuste » = norme valide formellement mais contraire aux principes de justice (dignité, égalité, droits fondamentaux).
\item \textbf{Problématique :} L'obéissance à la loi est-elle une obligation absolue (positivisme) ou conditionnée par sa conformité à la justice (jusnaturalisme) ? Comment trancher dans un État de droit pluraliste comme le Cameroun ?
\item \textbf{Annonce du plan :} I. L'obéissance comme exigence de l'ordre juridique (thèse positiviste). II. Les limites de l'obéissance : la justice comme critère de validité (thèse jusnaturaliste). III. La voie médiane : désobéissance civile et contrôle de constitutionnalité dans le pluralisme camerounais.
\end{itemize}
\textbf{Plan détaillé.}
\textbf{I. L'obéissance à la loi, condition de l'ordre et de la sécurité juridique (thèse positiviste)}
\begin{enumerate}
\item A. Kelsen : la validité ne dépend que de la procédure (pyramide des normes, \emph{Grundnorm}). La séparation Droit/Morale garantit la prévisibilité.
\item B. Hart : la règle de reconnaissance identifie la loi ; dire « c'est du droit, mais c'est mauvais » préserve la clarté morale et évite l'arbitraire subjectif.
\item C. Application camerounaise : l'article 45 de la Constitution (primauté des traités) et le contrôle du Conseil constitutionnel supposent qu'on identifie la loi \emph{avant} de la juger.
\end{enumerate}
\textbf{II. La loi injuste perd sa force obligatoire : le droit naturel comme limite (thèse jusnaturaliste)}
\begin{enumerate}
\item A. Augustin / Thomas d'Aquin : \emph{Lex injusta non est lex}. La loi tire sa force de sa conformité à la raison droite / loi éternelle.
\item B. Fuller : la « morale interne du droit » (8 principes) ; un système qui viole gravement la procédure (lois secrètes, rétroactives, impossibles à respecter) n'est pas un système juridique.
\item C. Finnis : les biens fondamentaux (vie, connaissance, sociabilité) sont la mesure de la justesse des lois.
\item D. Exemple : Nuremberg, apartheid, lois raciales — la communauté internationale a jugé que l'obéissance n'excusait pas.
\end{enumerate}
\textbf{III. Dépasser l'alternative : désobéissance civile, contrôle juridictionnel et pluralisme au Cameroun}
\begin{enumerate}
\item A. La désobéissance civile (Thoreau, Gandhi, King) : transgression publique, non-violente, assumée, pour réveiller la conscience collective. Ni anarchie, ni soumission.
\item B. Le contrôle de constitutionnalité et de conventionalité : le Conseil constitutionnel (a priori / a posteriori via QPC) et les juges ordinaires (contrôle de conventionnalité) permettent d'\emph{écarter} la loi injuste sans la violer.
\item C. Le pluralisme juridique camerounais comme laboratoire : le conflit coutume/Code civil (ex. succession des filles) montre que la hiérarchie des normes (Constitution $>$ Loi $>$ Coutume) et la clause de sauvegarde (non-contrariété aux droits fondamentaux) opèrent une \emph{résolution juridique} du conflit, évitant le face-à-face brutal obéissance/révolte.
\item D. Limite : l'effectivité de la justice (indépendance du CSM, accès au juge, coût des procédures) conditionne la réalité de ces garanties.
\end{enumerate}
\textbf{Conclusion.} Obéir à la loi injuste n'est ni un devoir absolu (positivisme naïf) ni un refus systématique (anarchie). L'État de droit moderne, au Cameroun comme ailleurs, institue des \emph{procédures} (contrôle de constitutionnalité, conventionnalité, désobéissance civile encadrée) pour que la justice ne reste pas un idéal abstrait mais devienne une norme effective. La philosophie du droit éclaire ainsi l'architecture même de nos institutions.
\end{corrige}

\coursec{Des formes de la Justice à la Justice comme valeur : Typologie aristotélicienne et moderne}

Dans la section précédente, nous avons vu que le droit positif (l'ensemble des lois écrites d'un État) ne coïncide pas toujours avec la justice : une loi peut être légale sans être juste. Mais alors, qu'est-ce que la justice elle-même ? Y a-t-il une seule justice, ou plusieurs formes de justice ? C'est la question qu'Aristote, il y a plus de 2300 ans, a été le premier à traiter de manière systématique dans le livre V de l'\emph{Éthique à Nicomaque}. Sa typologie reste aujourd'hui le socle de toute réflexion juridique et politique. Nous verrons ensuite comment les philosophes modernes et contemporains (Rawls, les théoriciens de la justice transitionnelle et environnementale) ont prolongé et dépassé cette grille, jusqu'à faire de la justice une \cle{valeur} fondatrice de nos sociétés.

\courssub{Justice générale et justice particulière : la typologie d'Aristote}

\textbf{Intuition de départ.} Imaginez deux situations. (a) Un chauffeur de taxi-moto à Douala respecte le code de la route, paie ses impôts, ne vole personne : on dit de lui qu'il est un homme \emph{juste}, au sens large. (b) Un proviseur doit répartir 10 bourses d'excellence entre 500 élèves : il doit être juste d'une autre manière, en \emph{répartissant} correctement. Aristote distingue précisément ces deux sens du mot « justice ».

\begin{definition}[Justice générale et justice particulière]
La \cle{justice générale} (ou justice légale) est la vertu complète : elle consiste à obéir aux lois de la cité et à respecter autrui dans toutes ses actions. Être juste, c'est alors être un bon citoyen. La \cle{justice particulière} est une vertu spécifique qui concerne la \emph{répartition} des biens et la \emph{régulation} des échanges entre les hommes. Elle se divise en deux espèces : la justice \emph{distributive} et la justice \emph{commutative}.
\end{definition}

Aristote formule cette distinction dans l'\emph{Éthique à Nicomaque}, livre V, chapitres 2 à 5. Son raisonnement part d'une observation simple : le mot « injuste » a deux sens. Est injuste celui qui \emph{viole la loi} (le délinquant), mais aussi celui qui \emph{prend plus que sa part} (le cupide, le profiteur). Au premier sens correspond la justice générale ; au second, la justice particulière, qui est le cœur du droit au sens strict.

\begin{attention}[Piège fréquent]
Ne confondez pas justice générale et justice particulière dans une dissertation. Dire « la justice c'est obéir aux lois » n'est qu'une moitié de la réponse (justice générale). Le cœur du chapitre au programme est la justice particulière : distributive et commutative.
\end{attention}

\begin{popfigure}
\begin{center}
\begin{tikzpicture}[scale=1]
% Triangle
\coordinate (A) at (0,4.2);
\coordinate (B) at (-4.5,0);
\coordinate (C) at (4.5,0);
\draw[very thick, popblue, fill=popblueL] (A) -- (B) -- (C) -- cycle;
% Sommet : justice générale
\node[draw=popdark, fill=popgoldL, rounded corners, align=center, font=\small\bfseries] at (0,3.1) {Justice générale\\ (obéissance aux lois)};
% Base gauche : distributive
\node[draw=popdark, fill=popgreenL, rounded corners, align=center, font=\small\bfseries] at (-2.6,0.9) {Justice\\ distributive\\ (géométrique)};
% Base droite : commutative
\node[draw=popdark, fill=poppinkL, rounded corners, align=center, font=\small\bfseries] at (2.6,0.9) {Justice\\ commutative\\ (arithmétique)};
% Flèche distributive : État vers citoyens
\node[circle, draw=popdark, fill=white, font=\footnotesize, inner sep=1.5pt] (E) at (-6.8,2.6) {État};
\node[circle, draw=popdark, fill=white, font=\footnotesize, inner sep=1.5pt] (c1) at (-6.8,0.4) {C1};
\node[circle, draw=popdark, fill=white, font=\footnotesize, inner sep=1.5pt] (c2) at (-5.9,0.4) {C2};
\node[circle, draw=popdark, fill=white, font=\footnotesize, inner sep=1.5pt] (c3) at (-7.7,0.4) {C3};
\draw[-{Stealth}, thick, popgreen] (E) -- (c1);
\draw[-{Stealth}, thick, popgreen] (E) -- (c2);
\draw[-{Stealth}, thick, popgreen] (E) -- (c3);
% Double flèche commutative
\node[circle, draw=popdark, fill=white, font=\footnotesize, inner sep=1.5pt] (d1) at (6.2,1.6) {A};
\node[circle, draw=popdark, fill=white, font=\footnotesize, inner sep=1.5pt] (d2) at (7.8,1.6) {B};
\draw[{Stealth}-{Stealth}, very thick, poppink] (d1) -- (d2);
\node[font=\footnotesize, popmuted, align=center] at (7.0,2.2) {échange\\ contrat, délit};
\end{tikzpicture}
\end{center}
\legende{Les trois justices d'Aristote. Au sommet, la justice générale (vertu du citoyen obéissant aux lois). À la base, les deux espèces de justice particulière : à gauche, la justice distributive (l'État répartit vers les citoyens C1, C2, C3 selon une proportion) ; à droite, la justice commutative (échange bilatéral d'égal à égal entre deux particuliers A et B).}
\end{popfigure}

\courssub{La justice distributive : la proportion géométrique}

\textbf{Intuition.} Trois ouvriers construisent une maison à Bafoussam. Le premier a travaillé 10 jours, le deuxième 5 jours, le troisième 5 jours. Le maître d'ouvrage verse \SI{200000}{} FCFA. Serait-il juste de donner la même somme à chacun ? Non, dit Aristote : la justice distributive exige que chacun reçoive \emph{proportionnellement} à sa contribution : \num{100000}, \num{50000} et \num{50000} FCFA.

\begin{definition}[Justice distributive]
La \cle{justice distributive} est la répartition proportionnelle des honneurs, des richesses et des charges entre les membres d'une communauté, selon un critère déterminé (mérite, besoin, travail). Elle suit une proportion \emph{géométrique} : Part de A / Mérite de A = Part de B / Mérite de B. Autrement dit : à mérites égaux, parts égales ; à mérites inégaux, parts inégales dans la même proportion.
\end{definition}

La formule aristotélicienne s'écrit :
\[ \frac{\text{Part de A}}{\text{Mérite de A}} = \frac{\text{Part de B}}{\text{Mérite de B}} \]

\begin{exemplebox}[Application numérique]
Deux candidats passent un concours de la fonction publique. A obtient 15/20, B obtient 10/20. Il y a 3 postes à pourvoir par tranche de mérite. La justice distributive impose que la probabilité d'accès aux postes respecte le rapport 15/10, c'est-à-dire 3/2 : A doit être prioritairement servi. De même, au Cameroun, les allocations familiales sont proportionnelles au nombre d'enfants à charge : critère du \emph{besoin}, non du mérite.
\end{exemplebox}

\textbf{Le grand problème : quel critère de répartition ?} Aristote lui-même reconnaît que le choix du critère dépend du régime politique. C'est un débat toujours actuel :

\begin{center}
\begin{tabularx}{\linewidth}{|l|X|X|}
\hline
\rowcolor{popblueL} \textbf{Critère} & \textbf{Doctrine politique} & \textbf{Exemple d'application} \\
\hline
Le mérite (naissance, excellence) & Aristocratie & Concours, grades, distinctions honorifiques \\
\hline
Le besoin & Socialisme & Allocations sociales, gratuité des soins pour les indigents \\
\hline
L'égalité stricte (chacun compte pour un) & Démocratie libérale & Suffrage universel : une personne = une voix \\
\hline
La liberté (droit d'acquérir et d'échanger) & Libéralisme (Robert Nozick) & Propriété privée, marché libre \\
\hline
\end{tabularx}
\end{center}

Nozick, dans \emph{Anarchie, État et Utopie} (1974), soutient qu'une distribution est juste si elle résulte d'acquisitions et d'échanges libres, quelle que soit l'inégalité finale. À l'inverse, les doctrines sociales estiment que l'État doit corriger le marché pour répondre aux besoins. Le débat sur les impôts progressifs au Cameroun illustre cette tension : faut-il taxer davantage les riches (critère du besoin social) ou laisser chacun jouir du fruit de son travail (critère libéral) ?

\courssub{La justice commutative : l'égalité arithmétique}

\textbf{Intuition.} Vous achetez un téléphone \SI{50000}{} FCFA au marché Mokolo à Yaoundé. Le vendeur vous remet un téléphone défectueux. La justice exige qu'il vous rembourse exactement \SI{50000}{} FCFA — ni plus, ni moins. Peu importe que vous soyez riche ou pauvre, ministre ou étudiant : dans l'échange, les personnes sont traitées comme \emph{égales}.

\begin{definition}[Justice commutative]
La \cle{justice commutative} (du latin \emph{commutare}, échanger) règle les rapports bilatéraux entre particuliers : contrats, ventes, prêts, mais aussi réparation des délits et dommages. Elle suit une égalité \emph{arithmétique} stricte, résumée par la formule latine \emph{do ut des} (« je donne pour que tu donnes ») : Valeur perdue = Valeur réparée. Elle ignore le mérite des personnes : elle ne regarde que l'égalité des prestations.
\end{definition}

Le rôle du juge, dans cette justice, est de \emph{rétablir l'égalité rompue}. Aristote compare le juge à un géomètre qui, devant une ligne coupée inégalement, retranche du segment trop long ce qui a été pris et l'ajoute au segment trop court. C'est exactement la logique des \cle{dommages-intérêts} en droit civil camerounais : si un entrepreneur détruit la clôture de son voisin, le tribunal condamne à réparer l'équivalent exact du préjudice.

\begin{exoresolu}[Distinguer distributive et commutative]
Énoncé. Le gouvernement camerounais décide (a) d'augmenter les bourses des étudiants les plus pauvres et (b) de condamner une entreprise de BTP à indemniser les riverains dont les maisons ont été fissurées par les travaux d'une route à Bafoussam. Qualifiez chaque mesure.
\tcblower
Solution. (a) Il s'agit d'une décision de l'État répartissant une ressource (les bourses) entre des citoyens selon un critère (le besoin, ici la pauvreté) : c'est la \textbf{justice distributive}, de type proportionnel au besoin. (b) Il s'agit d'un rapport entre deux parties privées (entreprise et riverains) où l'on répare un préjudice pour rétablir l'égalité rompue : c'est la \textbf{justice commutative}, sous sa forme corrective (dommages-intérêts). Le critère décisif est la structure du rapport : vertical (État $\to$ citoyens) pour la distributive, horizontal (particulier $\leftrightarrow$ particulier) pour la commutative.
\end{exoresolu}

\begin{methode}[Qualifier un conflit de justice]
Pour analyser une situation concrète dans une copie :
\begin{enumerate}
\item \textbf{Identifier les parties} : État/Citoyens (rapport vertical) $\Rightarrow$ justice distributive ou droit public ; Citoyen/Citoyen (rapport horizontal) $\Rightarrow$ justice commutative ou droit privé.
\item \textbf{Identifier l'enjeu} : s'agit-il de répartir des biens, honneurs ou charges (distributive) ? Ou de régler un échange, un contrat, un dommage (commutative) ?
\item \textbf{Choisir le principe pertinent} : proportionnalité au mérite ou au besoin (distributive) ; égalité stricte et réparation intégrale (commutative).
\end{enumerate}
\end{methode}

\courssub{Les prolongements modernes : justice réparatrice, sociale, transitionnelle, environnementale}

\textbf{3.4.1. La justice réparatrice (restaurative).} La justice pénale classique se contente de punir le coupable (prison, amende). Mais une fois le voleur emprisonné, la victime reste avec son préjudice, et le lien social reste rompu. La \cle{justice réparatrice} dépasse la sanction : elle vise la \emph{réparation du préjudice} et la \emph{restauration du lien social} par la médiation, le dialogue entre victime et auteur, la reconnaissance des torts. L'exemple le plus célèbre est celui des juridictions \emph{Gacaca} au Rwanda après le génocide de 1994 : des tribunaux populaires de proximité où les accusés avouaient, demandaient pardon et réparaient (travaux d'intérêt général), permettant de juger des centaines de milliers de dossiers impossibles à traiter par la justice classique. Au Cameroun, le projet de Commission Vérité, Justice et Réconciliation (CVJR), évoqué notamment dans le contexte de la crise anglophone, s'inscrit dans cette logique.

\textbf{3.4.2. La justice sociale selon John Rawls.} En 1971, le philosophe américain John Rawls publie \emph{Théorie de la justice}, l'ouvrage le plus important de philosophie politique du XX\textsuperscript{e} siècle.

\begin{aretenir}[Citation fondamentale]
« La justice est la première vertu des institutions sociales, comme la vérité l'est des systèmes de pensée. » (John Rawls, \emph{Théorie de la justice}, §1). Autrement dit : une institution injuste doit être réformée ou abolie, de même qu'une théorie fausse doit être rejetée, quelle que soit son efficacité.
\end{aretenir}

Rawls propose une expérience de pensée : la \cle{position originelle}. Imaginez que des individus doivent choisir les règles de leur société, mais placés derrière un \cle{voile d'ignorance} : ils ignorent s'ils seront riches ou pauvres, hommes ou femmes, anglophones ou francophones, en bonne santé ou handicapés. Que choisiraient-ils rationnellement ? Rawls répond : deux principes.

\begin{propriete}[Les deux principes de justice de Rawls]
\begin{enumerate}
\item \textbf{Principe des libertés égales} : chacun a droit au maximum de libertés fondamentales (expression, conscience, vote, propriété) compatibles avec les mêmes libertés pour tous.
\item \textbf{Principe de différence} : les inégalités sociales et économiques ne sont acceptables que si (a) elles profitent d'abord aux plus défavorisés, et (b) les positions avantageuses sont ouvertes à tous selon une équité réelle des chances.
\end{enumerate}
\end{propriete}

Le raisonnement est celui du \emph{maximin} : derrière le voile, un individu rationnel choisit la société où la \emph{pire} situation possible est la meilleure possible, car il pourrait l'occuper. Mieux vaut une société où le plus pauvre vit décemment qu'une société très riche en moyenne mais où le plus pauvre meurt de faim.

\begin{popfigure}
\begin{center}
\begin{tikzpicture}
\begin{axis}[
  width=0.85\linewidth, height=6cm,
  xlabel={Rang social (du plus favoris\'e au plus d\'efavoris\'e)},
  ylabel={Niveau de bien-\^etre},
  xmin=0, xmax=10, ymin=0, ymax=12,
  xtick={2,8}, xticklabels={Favoris\'es, D\'efavoris\'es},
  ytick=\empty,
  legend style={at={(0.02,0.98)}, anchor=north west, font=\footnotesize},
  axis lines=left]
% Société A : inégalitaire, défavorisés très bas
\addplot[popcurve, poppink, domain=0:10, samples=100]{11 - 1.05*x};
\addlegendentry{Soci\'et\'e A : in\'egalitaire (moyenne haute, minimum tr\`es bas)}
% Société B : rawlsienne, plancher relevé
\addplot[popcurve, popblue, domain=0:10, samples=100]{8.5 - 0.45*x};
\addlegendentry{Soci\'et\'e B : rawlsienne (maximin : plancher relev\'e)}
% Ligne du plancher
\draw[dashed, popgreen, thick] (axis cs:0,4.9) -- (axis cs:10,4.9);
\node[font=\footnotesize, popgreen, anchor=south west] at (axis cs:0.2,4.9) {Plancher minimal garanti (principe de diff\'erence)};
\end{axis}
\end{tikzpicture}
\end{center}
\legende{Illustration schématique du raisonnement de Rawls derrière le voile d'ignorance. La société A a une richesse moyenne plus élevée mais ses plus défavorisés sont très bas. La société B, conforme au principe de différence, relève le plancher : derrière le voile d'ignorance, on choisit rationnellement B (stratégie du maximin).}
\end{popfigure}

\begin{attention}[La justice sociale n'est pas la charité]
Erreur classique de dissertation : assimiler justice sociale et générosité. L'aumône dépend du bon vouloir du donneur ; la justice sociale est une \emph{exigence structurelle} (Rawls) et \emph{légale} : le Code du travail camerounais, le salaire minimum (SMIG), l'assurance maladie sont des \emph{droits} exigibles, pas des faveurs. Confondre aumône et droit est une erreur politique : le bénéficiaire de la charité doit remercier ; le titulaire d'un droit peut réclamer.
\end{attention}

\textbf{3.4.3. La justice transitionnelle.} Comment une société sort-elle d'une dictature, d'une guerre civile ou d'une crise grave ? Ni l'oubli pur (l'injustice reste vivace) ni la vengeance (le cycle recommence) ne fonctionnent. La \cle{justice transitionnelle} combine quatre mécanismes, appelés « piliers ».

\begin{popfigure}
\begin{center}
\begin{tikzpicture}[mindmap, grow cyclic,
  every node/.style={concept, font=\small},
  concept color=popblueL,
  level 1/.append style={level distance=3.4cm, sibling angle=90},
  level 2/.append style={level distance=2.2cm, sibling angle=60, font=\footnotesize}]
\node[concept, fill=popgoldL, font=\small\bfseries] {Justice transitionnelle}
  child[concept color=poppinkL] { node {Poursuites judiciaires}
    child { node {Tribunal militaire vs civil (Cameroun)} } }
  child[concept color=popgreenL] { node {Commissions V\'erit\'e}
    child { node {Projet CVJR ; CNPBM (2017)} } }
  child[concept color=poporangeL] { node {R\'eparations}
    child { node {Indemnisations, m\'emoriaux} } }
  child[concept color=poppurpleL] { node {R\'eformes institutionnelles}
    child { node {Centres DDR (ex-combattants)} } };
\end{tikzpicture}
\end{center}
\legende{Carte mentale : les quatre piliers de la justice transitionnelle et leurs illustrations camerounaises dans le contexte de la crise anglophone (depuis 2016) : Commission nationale pour la promotion du bilinguisme et du multiculturalisme (CNPBM), centres de Désarmement, Démobilisation et Réintégration (DDR) pour les ex-combattants, débat sur la juridiction compétente (tribunal militaire ou civil) pour juger les crimes commis.}
\end{popfigure}

Le cas camerounais est directement au programme des débats contemporains : face à la crise dans les régions du Nord-Ouest et du Sud-Ouest depuis 2016, l'État a créé la CNPBM (2017) et des centres DDR à Bamenda, Buea et Mora, tandis que la société civile réclame une véritable commission Vérité-Justice-Réconciliation. La question philosophique posée : peut-on réconcilier sans vérité ? Peut-on punir sans diviser davantage ?

\textbf{3.4.4. La justice environnementale et climatique.} Nouveau front de la justice, elle repose sur trois principes : le principe \cle{pollueur-payeur} (qui pollue doit réparer — extension de la justice commutative à l'environnement) ; l'\cle{équité intergénérationnelle} (nous devons transmettre aux générations futures une planète habitable — extension de la justice distributive dans le temps) ; l'accès à l'information et la participation des populations (consacrés par la loi camerounaise n°96/12 du 5 août 1996 relative à la gestion de l'environnement, inspirée des standards internationaux comme la Convention d'Aarhus de 1998). Au Cameroun, l'affaire Herakles Farms (projet de palmeraie industrielle dans le Sud-Ouest, contesté pour déforestation et spoliation des terres villageoises) et l'exploitation minière dans la région de l'Est illustrent ces conflits : qui supporte les dommages écologiques ? Souvent les populations locales les plus pauvres, tandis que les profits partent ailleurs — une injustice distributive à l'échelle écologique.

\begin{savaistu}[Le Cameroun et la justice internationale]
Le Cameroun a signé le Statut de Rome (instituant la Cour pénale internationale) dès 1998 et l'a ratifié en 2006 ; son Code pénal révisé de 2016 (Loi n°2016/007) n'a toutefois intégré que partiellement les définitions des crimes internationaux (génocide, crimes contre l'humanité). Conséquence : certains crimes de masse commis sur le territoire restent difficiles à poursuivre devant les juridictions nationales, ce qui ravive le débat sur l'effectivité de la justice.
\end{savaistu}

\courssub{De la typologie à la valeur : la justice comme exigence absolue}

Parcourons le chemin : Aristote décrit des \emph{formes} de justice (répartir, échanger, réparer) ; Rawls en fait le \emph{premier} critère de légitimité des institutions ; les justices transitionnelle et environnementale l'étendent à l'histoire et à la planète. La justice est ainsi passée du statut de vertu parmi d'autres à celui de \cle{valeur fondatrice} : on ne demande plus seulement « cette répartition est-elle proportionnée ? » mais « cette société est-elle juste ? ». C'est cette exigence qui fonde le préambule des constitutions modernes, dont celle du Cameroun qui proclame l'attachement du peuple aux droits de l'homme.

\begin{exercice}[Entraînement à la dissertation]
Sujet : « Y a-t-il plusieurs justices ou une seule justice ? »
Consignes : (1) Définissez justice générale et justice particulière ; (2) Montrez, avec Aristote, la pluralité des formes (distributive, commutative) ; (3) Interrogez l'unité profonde de la justice comme valeur (Rawls) ; (4) Illustrez par un exemple camerounais (bourses, dommages-intérêts, crise anglophone ou justice environnementale). Rédigez l'introduction et le plan détaillé.
\end{exercice}

\begin{corrige}[Exercice — éléments de corrigé]
Introduction : accroche possible par l'expérience quotidienne (un proviseur répartit des bourses, un juge condamne un voleur : dit-on « juste » dans le même sens ?). Problématique : la justice est-elle une vertu unique ou un ensemble de formes hétérogènes ? Plan dialectique possible : I. La pluralité des formes (Aristote : générale vs particulière ; distributive proportionnelle vs commutative arithmétique — exemples chiffrés). II. L'extension moderne du concept (sociale chez Rawls, réparatrice, transitionnelle, environnementale) qui semble éclater encore la notion. III. L'unité retrouvée : toutes ces formes expriment une même exigence — rendre à chacun son dû (\emph{suum cuique}) et traiter les égaux égaux — qui fait de la justice la valeur première des institutions (Rawls). Conclusion : pluralité des formes, unité de la valeur.
\end{corrige}

\coursec{La Justice au Cameroun : Institutions, Effectivité et Défis majeurs}

Après avoir distingué les formes de la justice (distributive, commutative, corrective) et montré que la justice est aussi une valeur, il nous faut descendre de l'idéal vers le réel : comment la justice est-elle \emph{instituée} et \emph{effective} au Cameroun ? Une justice qui n'existerait que dans les textes serait une justice formelle sans effectivité. Cette section examine l'architecture judiciaire camerounaise, le pluralisme juridique (droit moderne et droit coutumier), les obstacles à l'accès au droit, la question de l'indépendance des juges, les délais, et enfin les réformes en cours.

\courssub{L'architecture judiciaire camerounaise}

\begin{definition}[Ordre juridictionnel]
Un \cle{ordre juridictionnel} est un ensemble de tribunaux hiérarchisés compétents pour une catégorie de litiges. Le Cameroun connaît deux ordres principaux : l'\cle{ordre judiciaire} (litiges entre particuliers, matière pénale) et l'\cle{ordre administratif} (litiges impliquant l'État et les personnes publiques).
\end{definition}

\textbf{L'ordre judiciaire} comprend, de bas en haut : les \cle{Tribunaux d'instance} (petits litiges civils et délits mineurs), les \cle{Tribunaux de grande instance} (affaires civiles importantes, crimes et délits), les \cle{Cours d'appel} (réexamen des jugements), et enfin la \cle{Cour Suprême} (Chambre judiciaire), juge de cassation qui ne juge pas les faits mais vérifie l'application correcte du droit.

\textbf{L'ordre administratif} comprend les \cle{Tribunaux administratifs}, la \cle{Cour administrative d'appel} et la \cle{Cour Suprême} (Chambre administrative). Il juge par exemple les litiges entre un citoyen et une administration (expropriation, licenciement d'un fonctionnaire, concours publics).

S'y ajoutent des \cle{juridictions spécialisées} : tribunaux militaires, tribunaux de commerce (créés par la loi de 2023), juridictions comptables. Le contentieux constitutionnel est dévolu au \cle{Conseil Constitutionnel}, la Cour Constitutionnelle pleinement distincte n'étant pas encore opérationnelle.

\begin{popfigure}
\begin{center}
\begin{tikzpicture}[
  box/.style={draw=popblue, thick, rounded corners, fill=popblueL, text width=3.4cm, align=center, minimum height=0.8cm, font=\small},
  boxa/.style={draw=poporange, thick, rounded corners, fill=poporangeL, text width=3.4cm, align=center, minimum height=0.8cm, font=\small},
  top/.style={draw=popdark, very thick, rounded corners, fill=popgoldL, text width=5.5cm, align=center, minimum height=0.9cm, font=\small\bfseries},
  spec/.style={draw=poppurple, thick, rounded corners, fill=poppurpleL, text width=4.2cm, align=center, minimum height=0.8cm, font=\small},
  arr/.style={-{Stealth[length=2.5mm]}, thick, popmuted}]
  \node[top, align=center] (cs) at (0,0){Cour Suprême\\ (Ch. judiciaire / Ch. administrative)};
  \node[box] (ca) at (-3.2,-1.6) {Cours d'appel};
  \node[boxa] (caa) at (3.2,-1.6) {Cour administrative d'appel};
  \node[box] (tgi) at (-3.2,-3.2) {Tribunaux de grande instance};
  \node[boxa] (ta) at (3.2,-3.2) {Tribunaux administratifs};
  \node[box] (ti) at (-3.2,-4.8) {Tribunaux d'instance};
  \node[spec] (sp) at (0,-6.4) {Juridictions spécialisées : militaire, commerce, comptable};
  \node[spec, fill=popgreenL, draw=popgreen] (cc) at (0,-7.6) {Conseil Constitutionnel};
  \draw[arr] (ti) -- (tgi); \draw[arr] (tgi) -- (ca); \draw[arr] (ca) -- (cs);
  \draw[arr] (ta) -- (caa); \draw[arr] (caa) -- (cs);
\end{tikzpicture}
\end{center}
\legende{Organisation judiciaire du Cameroun (loi de 2023) : deux ordres convergent vers la Cour Suprême ; les flèches indiquent les voies de recours (appel puis cassation).}
\end{popfigure}

\begin{propriete}[Principe du double degré de juridiction]
Tout justiciable a le droit de faire rejuger sa cause par une juridiction supérieure : c'est le \cle{droit à l'appel}, garantie essentielle du procès équitable. Exceptions : les matières gracieuses et les petites créances, jugées en premier et dernier ressort.
\end{propriete}

\courssub{La justice coutumière : un pluralisme juridique}

Le Cameroun connaît un \cle{pluralisme juridique} : à côté du droit moderne coexistent les \cle{chefferies traditionnelles} (de 1er, 2e et 3e degré), compétentes en matière de droit personnel, familial, successoral et foncier coutumier. Leurs décisions peuvent faire l'objet d'un appel devant le Tribunal d'instance, juge de droit commun.

\begin{attention}[Les limites de la justice coutumière]
La justice coutumière souffre de trois défauts majeurs : l'\cle{absence d'écrits} (pas de procès-verbaux fiables, donc contrôle difficile en appel), le risque de \cle{partialité} (le chef est parfois partie au conflit), et le \cle{non-respect des droits des femmes et des enfants} (successions excluant les filles, mariages forcés). Or le droit coutumier ne peut déroger ni à la Constitution ni aux conventions internationales ratifiées par le Cameroun.
\end{attention}

\begin{attention}[La « justice populaire » n'est pas de la justice]
Le lynchage ou \emph{jungle justice} n'est pas une forme de justice : c'est une violation du \cle{monopole de la violence légitime} de l'État (Max Weber) et du droit à un procès équitable. Punir sans procès, sans juge impartial et sans droit de défense, c'est détruire la justice au nom de la justice.
\end{attention}

\courssub{L'accès au droit : entre droits proclamés et obstacles réels}

\begin{definition}[Aide juridictionnelle]
L'\cle{aide juridictionnelle} est la prise en charge par l'État des frais de justice (avocat, huissier, expertise) pour les justiciables dont les ressources sont insuffisantes. Elle est instituée au Cameroun par la loi de 2005. S'y ajoutent les \cle{cliniques juridiques} (universités, ONG) et les \cle{Maisons de justice}.
\end{definition}

L'accès effectif au juge se heurte cependant à quatre obstacles : le \cle{coût} (timbre fiscal, honoraires d'avocat), la \cle{distance} (les « déserts judiciaires » : certaines sous-préfectures sont à des centaines de kilomètres du tribunal), la \cle{langue} (la justice se rend en français ou en anglais, rarement dans les langues locales), et l'\cle{analphabétisme juridique} (ignorance de ses droits et des procédures).

\begin{exemplebox}[Un budget en deçà des standards]
Le budget de la Justice en 2024 est d'environ 65 milliards de FCFA, soit moins de 1\% du budget de l'État, alors que le standard international (ONU) recommande 2 à 3\%. Conséquence : environ 2 000 magistrats pour 28 millions d'habitants ; il manquerait près de 1 500 magistrats pour atteindre le ratio d'un magistrat pour 10 000 habitants.
\end{exemplebox}

\begin{popfigure}
\begin{center}
\begin{tikzpicture}[scale=0.9]
  \draw[thick, popdark, fill=popcream] plot coordinates {(0,0) (1.2,1.4) (2.6,2.0) (3.4,3.4) (2.8,4.6) (3.6,5.8) (3.0,7.0) (1.4,7.6) (0.2,6.4) (-0.8,5.0) (-1.4,3.4) (-1.0,1.6)} -- cycle;
  % régions (schématiques)
  \fill[popgreen!60] (0.4,0.9) circle (0.55); % Centre
  \fill[popgreen!60] (-0.5,1.7) circle (0.45); % Littoral
  \fill[popgold!70] (1.6,2.6) circle (0.5); % Ouest
  \fill[popgold!70] (0.2,2.9) circle (0.45); % Nord-Ouest/Sud-Ouest
  \fill[poppink!55] (2.2,4.4) circle (0.6); % Adamaoua/Est
  \fill[poppink!55] (1.6,6.2) circle (0.55); % Nord
  \fill[poppink!70] (1.0,7.1) circle (0.45); % Extrême-Nord
  \node[font=\scriptsize] at (0.4,0.9) {Centre};
  \node[font=\scriptsize] at (-0.5,1.7) {Litt.};
  \node[font=\scriptsize] at (1.6,2.6) {Ouest};
  \node[font=\scriptsize] at (2.2,4.4) {Est/Ad.};
  \node[font=\scriptsize] at (1.6,6.2) {Nord};
  \node[font=\scriptsize] at (1.0,7.1) {E.-N.};
  \draw[-{Stealth}, thick, popdark] (-1.8,6.6) -- (-1.8,7.4) node[above, font=\scriptsize]{N};
  % légende
  \fill[poppink!60] (-1.9,-0.9) rectangle (-1.5,-0.6); \node[right, font=\scriptsize] at (-1.4,-0.75) {< 50 magistrats (faible dotation)};
  \fill[popgold!70] (-1.9,-1.4) rectangle (-1.5,-1.1); \node[right, font=\scriptsize] at (-1.4,-1.25) {50 à 150 magistrats};
  \fill[popgreen!60] (-1.9,-1.9) rectangle (-1.5,-1.6); \node[right, font=\scriptsize] at (-1.4,-1.75) {> 150 magistrats (bonne dotation)};
\end{tikzpicture}
\end{center}
\legende{Carte schématique (croquis simplifié, données simulées réalistes, 2023) : densité de magistrats par région. L'Extrême-Nord, l'Est et l'Adamaoua sont sous-dotés ; le Centre et le Littoral concentrent l'essentiel des juridictions.}
\end{popfigure}

\courssub{Indépendance de la justice et délais : les défis majeurs}

\begin{definition}[Indépendance de la justice]
L'\cle{indépendance de la justice} signifie que le juge ne statue qu'en fonction du droit, sans recevoir d'ordre du pouvoir exécutif ni subir de pression. Elle est la condition de la justice comme valeur : un juge dépendant est un juge injuste.
\end{definition}

Au Cameroun, la carrière des magistrats est gérée par le \cle{Conseil Supérieur de la Magistrature} (CSM), présidé par le Président de la République, le Garde des Sceaux en étant vice-président. Les magistrats du \cle{parquet} (procureurs) sont hiérarchisés et reçoivent des instructions écrites. Cette organisation nourrit le débat : le pouvoir judiciaire est-il une simple \emph{autorité} (terme de la Constitution) ou un véritable \emph{pouvoir} séparé ? La perception de la corruption reste forte : selon Transparency International (CPI 2023), le Cameroun se classe 140\textsuperscript{e} sur 180.

Autre défi : les \cle{délais}. Un procès civil au tribunal de première instance dure en moyenne 3 à 5 ans. Causes : pénurie de magistrats, infrastructures vétustes, lourdeur de la procédure écrite, grèves récurrentes.

\begin{aretenir}[« La justice retardée est une justice niée »]
Cette formule, attribuée à William Gladstone, s'applique aux délais camerounais : la détention préventive y dure souvent bien au-delà des délais légaux. Un jugement juste mais rendu trop tard cesse d'être juste : la \emph{temporalité} fait partie de la justice elle-même.
\end{aretenir}

\begin{popfigure}
\begin{center}
\begin{tikzpicture}
\begin{axis}[
  ybar, bar width=0.35cm, width=0.95\linewidth, height=6cm,
  symbolic x coords={Civil, Pénal, Administratif, Coutumier},
  xtick=data, ylabel={Durée moyenne (années)}, ymin=0, ymax=6.5,
  legend style={at={(0.5,1.02)}, anchor=south, legend columns=2, font=\small},
  nodes near coords, every node near coord/.append style={font=\scriptsize},
  ylabel style={font=\small}]
\addplot[fill=popblue] coordinates {(Civil,4.5) (Pénal,3.0) (Administratif,5.0) (Coutumier,1.0)};
\addplot[fill=poporange] coordinates {(Civil,3.5) (Pénal,2.5) (Administratif,4.0) (Coutumier,0.8)};
\legend{2015, 2023}
\end{axis}
\end{tikzpicture}
\end{center}
\legende{Durée moyenne des procédures par matière (données simulées réalistes) : légère amélioration entre 2015 et 2023, mais les contentieux administratif et civil restent très longs.}
\end{popfigure}

\courssub{Réformes récentes et étude de cas}

Les réformes récentes témoignent d'une volonté de modernisation : la \cle{loi de 2023} sur l'organisation judiciaire (création des tribunaux de commerce et de cours d'appel spécialisées), la \cle{numérisation} (e-justice, e-filing à Douala et Yaoundé) et la formation continue des magistrats à l'ENAM.

\begin{exoresolu}[Étude de cas : l'affaire de la chefferie de Bafut (Nord-Ouest)]
Énoncé : À la mort du fon (chef) de Bafut, deux prétendants revendiquent la succession selon la coutume. Le ministère de l'Administration territoriale reconnaît l'un d'eux par arrêté. L'autre saisit le tribunal administratif pour excès de pouvoir. Identifier les ordres de juridiction en présence et ce que ce cas illustre philosophiquement.
\tcblower
Solution : (1) Le litige de succession relève d'abord de la \emph{justice coutumière} (droit personnel et successoral), avec appel possible devant le tribunal d'instance. (2) L'arrêté du ministre est un \emph{acte administratif} : son annulation relève du \emph{tribunal administratif}, puis de la Cour administrative d'appel, puis de la Cour Suprême (chambre administrative). (3) Ce cas illustre le \cle{pluralisme juridictionnel} camerounais (droit coutumier, droit moderne, droit administratif) et le \emph{rôle politique de la justice} : en reconnaissant un chef, l'État fait de la tradition un enjeu de pouvoir, et le juge administratif devient arbitre entre légitimité coutumière et légalité étatique.
\end{exoresolu}

\begin{methode}[Lire un arrêt de la Cour Suprême]
\begin{enumerate}
\item \textbf{Le visa} : relever les textes cités (Constitution, lois, conventions).
\item \textbf{Faits et procédure} : résumer le litige et les décisions antérieures.
\item \textbf{Les moyens} : identifier les arguments des parties.
\item \textbf{Les motifs} : analyser le raisonnement juridique du juge.
\item \textbf{Le dispositif} : dégager la décision finale (rejet, cassation, annulation).
\end{enumerate}
Repérer la \emph{ratio decidendi} : la règle de droit qui fonde la décision et qui fera jurisprudence.
\end{methode}

\begin{savaistu}[Le saviez-vous ?]
L'e-filing (dépôt électronique des actes de procédure), expérimenté à Douala et Yaoundé, permet à un avocat d'introduire une requête sans se déplacer au greffe : un outil concret contre les « déserts judiciaires » et les délais, qui rapproche la justice du justiciable.
\end{savaistu}

\begin{aretenir}[L'essentiel de la section]
La justice camerounaise repose sur deux ordres juridictionnels coiffés par la Cour Suprême, complétés par une justice coutumière aux limites réelles. Son effectivité est entravée par le coût, la distance, la pénurie de magistrats, les délais et les doutes sur l'indépendance. Les réformes (loi de 2023, numérisation) montrent que la justice comme \emph{valeur} exige un effort constant pour devenir justice \emph{effective} : c'est là un enjeu philosophique autant que politique.
\end{aretenir}

\coursec{Méthodologie : Argumenter, Dissertation et Explication de texte (Entraînement Bac)}

Après avoir étudié les institutions camerounaises de la justice et leurs défis concrets, il nous reste à transformer ce savoir en \cle{compétence d'examen}. Au Probatoire et au Baccalauréat technique, l'épreuve de philosophie ne demande pas de réciter : elle demande de \emph{penser par écrit}, c'est-à-dire de problématiser un sujet, d'argumenter avec rigueur et de mobiliser les auteurs. Cette dernière séquence vous donne les outils pour réussir les deux exercices possibles : la \cle{dissertation} et l'\cle{explication de texte}.

\courssub{Analyse d'un sujet type Bac : « Le droit positif suffit-il à garantir la justice ? »}

\textbf{Premier temps : dégager les mots-clés.} Tout sujet de dissertation repose sur trois ou quatre termes qu'il faut définir avant d'écrire la moindre phrase. Ici :

\begin{itemize}
\item \cle{Droit positif} : l'ensemble des règles juridiques effectivement en vigueur dans un État, à un moment donné (Constitution, lois votées, règlements, jurisprudence, coutumes reconnues). Il est « positif » parce qu'il est \emph{posé} par les hommes, et non découvert dans la nature.
\item \cle{Garantir} : assurer, rendre certain, protéger contre tout échec. Le mot est fort : il ne s'agit pas de « favoriser » ou de « tendre vers », mais de \emph{suffire à assurer}.
\item \cle{Justice} : l'idéal d'équité qui consiste à rendre à chacun ce qui lui est dû, tant sur le plan distributif (répartition des biens et des charges) que correctif (réparation des torts).
\end{itemize}

\textbf{Deuxième temps : repérer l'opérateur logique.} Le verbe interrogatif « suffit-il » est un mot-opérateur qui appelle une réponse nuancée : il invite à distinguer la \emph{condition nécessaire} (sans laquelle rien n'est possible) de la \emph{condition suffisante} (qui suffit à elle seule). Toute la dissertation se jouera sur cette distinction.

\begin{attention}[Erreur fatale : le réductionnisme juridique]
Écrire « le droit, c'est la loi » est une faute qui plafonne la copie. Le droit positif inclut aussi la \cle{jurisprudence} (les décisions des juges), la \cle{coutume}, les principes généraux du droit et les traités internationaux ratifiés par le Cameroun. Autre piège : une copie sans aucune référence à un auteur (Aristote, Kant, Rawls...) est une copie plafonnée. Citez brièvement, mais citez.
\end{attention}

\courssub{Problématisation : la tension entre sécurité juridique et équité matérielle}

L'intuition de départ est simple : sans lois écrites, la justice serait livrée à l'arbitraire du plus fort. Mais l'observation de la réalité camerounaise montre aussitôt la difficulté : une loi peut être égale pour tous dans son texte et inégale dans ses effets. Par exemple, la législation foncière camerounaise (ordonnances de 1974) fait de l'immatriculation le mode principal de constatation de la propriété, formellement identique pour tous ; matériellement, seuls les citoyens informés et disposant de ressources peuvent accomplir la procédure, si bien que des communautés rurales risquent de perdre des terres qu'elles occupent depuis des générations.

La \cle{problématique} naît de cette tension : le droit positif assure la \emph{sécurité juridique} (prévisibilité des règles, égalité formelle devant la loi, limitation de l'arbitraire), mais l'équité matérielle exige parfois d'\emph{adapter} la règle au cas particulier et de \emph{corriger} les inégalités structurelles. On peut donc formuler la problématique ainsi : \emph{la règle générale et égale suffit-elle à réaliser l'équité concrète, ou la justice exige-t-elle un supplément qui dépasse le droit positif ?}

\begin{methode}[L'introduction « en entonnoir »]
\begin{enumerate}
\item \textbf{Accroche} : citation, fait d'actualité, étymologie ou paradoxe. Exemple : « "Il n'y a point de pire injustice que de traiter également des choses inégales", écrit Aristote. »
\item \textbf{Définitions rigoureuses} des termes du sujet (droit positif, garantir, justice).
\item \textbf{Reformulation} du sujet en une phrase claire.
\item \textbf{Problématique} : la question-tension, souvent introduite par « Mais alors... » ou « Cependant... ».
\item \textbf{Annonce de plan} : « Nous verrons d'abord que..., puis que..., enfin que... ».
\end{enumerate}
\end{methode}

\courssub{Construction du plan dialectique : thèse, antithèse, synthèse}

Le plan dialectique est le plan-roi du sujet interrogatif. Il suit le mouvement même de la pensée qui doute.

\textbf{I. Thèse — Le droit positif comme condition nécessaire de la justice.} Sans lois écrites, publiques et stables, chacun jugerait selon ses passions : la vengeance privée remplacerait le procès. Le droit positif assure la sécurité juridique, l'égalité formelle (la même loi s'applique à tous) et limite l'arbitraire du juge comme du puissant. Référence : Kant, pour qui le droit est « l'ensemble des conditions sous lesquelles l'arbitre de l'un peut s'accorder avec l'arbitre de l'autre selon une loi universelle de la liberté ». Exemple : le Code pénal camerounais définit les infractions et les peines, empêchant les sanctions arbitraires.

\textbf{II. Antithèse — Mais condition insuffisante.} La loi est générale, donc rigide : appliquée mécaniquement, elle peut produire de l'injustice. Surtout, une loi positive peut être \emph{injuste} : l'apartheid sud-africain ou les lois ségrégationnistes étaient du droit positif formellement valide. Enfin, les inégalités de fait (richesse, instruction, accès aux avocats) rendent l'égalité formelle parfois fictive. Référence : Aristote distingue déjà la justice distributive, qui répartit \emph{proportionnellement} aux mérites, de la simple égalité arithmétique.

\textbf{III. Synthèse — La justice comme horizon critique et finalité du droit.} Le dépassement ne consiste pas à abolir le droit positif, mais à le soumettre à un idéal de justice qui le juge et le corrige : l'équité (le juge adapte la règle au cas particulier), les droits fondamentaux, le contrôle de conformité des lois à la Constitution et aux traités internationaux ratifiés (comme la CEDAW, Convention sur l'élimination de toutes les formes de discrimination à l'égard des femmes). Rawls montre, avec le « voile d'ignorance », que des lois justes sont celles que des personnes choisiraient sans savoir quelle place elles occuperont dans la société. Le droit positif est ainsi l'\emph{instrument} nécessaire de la justice, mais la justice demeure sa \emph{finalité} et son juge.

\begin{popfigure}
\begin{center}
\begin{tikzpicture}[font=\small]
% Entonnoir introduction
\fill[popblueL] (-3.2,4.2) -- (3.2,4.2) -- (1.1,3.2) -- (-1.1,3.2) -- cycle;
\node at (0,3.75) {\textbf{Introduction} : accroche $\rightarrow$ problématique};
% Thèse / Antithèse
\draw[fill=popgreenL, rounded corners] (-3.6,0.4) rectangle (-0.4,2.9);
\node[align=center] at (-2,1.65) {\textbf{I. Thèse}\\ Le droit positif,\\ condition nécessaire\\ (sécurité, égalité formelle)};
\draw[fill=poporangeL, rounded corners] (0.4,0.4) rectangle (3.6,2.9);
\node[align=center] at (2,1.65) {\textbf{II. Antithèse}\\ Condition insuffisante\\ (rigidité, lois injustes,\\ inégalités de fait)};
% Synthèse
\draw[fill=poppurpleL, rounded corners] (-2.2,-1.3) rectangle (2.2,0.1);
\node[align=center] at (0,-0.6) {\textbf{III. Synthèse} : la justice, horizon critique\\ et finalité du droit};
% Entonnoir inversé conclusion
\fill[popgoldL] (-1.1,-2.2) -- (1.1,-2.2) -- (3.2,-3.2) -- (-3.2,-3.2) -- cycle;
\node at (0,-2.65) {\textbf{Conclusion} : bilan $\rightarrow$ réponse $\rightarrow$ ouverture};
% Flèches
\draw[-{Stealth}, thick, popdark] (0,3.15) -- (0,2.95);
\draw[-{Stealth}, thick, popdark] (-2,0.35) -- (-1,-0.05);
\draw[-{Stealth}, thick, popdark] (2,0.35) -- (1,-0.05);
\draw[-{Stealth}, thick, popdark] (0,-1.35) -- (0,-2.15);
\end{tikzpicture}
\end{center}
\legende{Structure type d'une dissertation dialectique de philosophie au Baccalauréat technique : entonnoir d'introduction, confrontation thèse/antithèse, synthèse, entonnoir inversé de conclusion.}
\end{popfigure}

\begin{savaistu}[Les sujets réels du Bac technique]
Au Baccalauréat technique, les sujets portent souvent sur l'articulation Droit / Morale / Politique. Exemples : « La loi est-elle la seule limite de la liberté ? » ; « Le juge doit-il se borner à appliquer la loi ? ». Vous remarquez que ces sujets mobilisent exactement les notions de cette leçon : droit positif, justice, liberté, équité.
\end{savaistu}

\courssub{La rédaction : paragraphes, transitions, conclusion}

\textbf{Le développement.} Chaque paragraphe suit la structure : \emph{Idée force} (une phrase) $\rightarrow$ \emph{Argument} (le raisonnement) $\rightarrow$ \emph{Exemple ou référence} (auteur ou fait camerounais) $\rightarrow$ \emph{Transition} (qui annonce la limite de l'idée et prépare le paragraphe suivant). Une partie compte deux ou trois paragraphes, jamais un seul.

\begin{exemplebox}[Astuce Bac : l'ancrage territorial]
Illustrez toujours par un exemple camerounais : la loi de 2004 sur la décentralisation face à la lenteur réelle des transferts de compétences ; le Code pénal face aux débats bioéthiques ; les coutumes successorales face à la CEDAW ratifiée par le Cameroun. Cet ancrage montre à l'examinateur que vous savez \emph{appliquer les notions à des situations concrètes}, savoir-faire explicitement exigé par le programme.
\end{exemplebox}

\textbf{La conclusion.} Elle comporte trois temps : le \emph{bilan} (résumé du chemin parcouru), la \emph{réponse} claire à la question posée (ici : le droit positif est nécessaire mais non suffisant ; il doit être orienté et corrigé par l'idéal de justice), et l'\emph{ouverture} (une question nouvelle : « Mais qui, alors, doit incarner cet idéal : le juge, le législateur ou le citoyen ? »).

\courssub{L'explication de texte : la méthode en quatre temps}

L'explication de texte est l'autre exercice possible. Elle n'est ni un résumé, ni une dissertation déguisée : il s'agit de \emph{faire parler le texte}, de dégager sa thèse et de la discuter.

\begin{methode}[Les quatre temps de l'explication de texte]
\begin{enumerate}
\item \textbf{Présentation} : auteur, œuvre, date, contexte, thèse générale du texte, place du passage dans l'argumentation.
\item \textbf{Analyse structurée} : suivre le texte \emph{dans l'ordre}, dégager les idées directrices, expliquer chaque argument et définir les concepts clés.
\item \textbf{Discussion critique} : portée et limites de la thèse, objections, exemples contraires, confrontation avec d'autres auteurs.
\item \textbf{Conclusion} : intérêt philosophique du texte et actualité de la question posée.
\end{enumerate}
\end{methode}

\begin{definition}[Ratio decidendi et obiter dictum]
Dans une décision de justice, la \cle{ratio decidendi} est le motif juridique essentiel qui fonde la décision : c'est lui qui a force obligatoire. L'\cle{obiter dictum} est un propos accessoire du juge, qui n'a pas d'autorité de chose jugée. Savoir distinguer l'un de l'autre, c'est savoir lire un arrêt — compétence utile pour analyser un texte juridique au Bac.
\end{definition}

\begin{popfigure}
\begin{center}
\begin{tikzpicture}[font=\scriptsize, x=0.92cm]
\draw[fill=popblueL] (0,0) rectangle (15.2,0.75);
\node[popdark] at (1.15,0.37) {\textbf{Paragraphe}};
\node[popdark] at (3.6,0.37) {\textbf{Idée principale}};
\node[popdark] at (6.1,0.37) {\textbf{Concept clé}};
\node[popdark] at (8.6,0.37) {\textbf{Argument}};
\node[popdark] at (11.1,0.37) {\textbf{Exemple cité}};
\node[popdark] at (13.6,0.37) {\textbf{Mon commentaire}};
\draw (0,-0.75) rectangle (15.2,0);
\node[popdark] at (1.15,-0.37) {lignes 1--5};
\node[popdark] at (3.6,-0.37) {\ldots};
\node[popdark] at (6.1,-0.37) {\ldots};
\node[popdark] at (8.6,-0.37) {\ldots};
\node[popdark] at (11.1,-0.37) {\ldots};
\node[popdark] at (13.6,-0.37) {\ldots};
\draw (0,-1.5) rectangle (15.2,-0.75);
\node[popdark] at (1.15,-1.12) {lignes 6--12};
\node[popdark] at (3.6,-1.12) {\ldots};
\node[popdark] at (6.1,-1.12) {\ldots};
\node[popdark] at (8.6,-1.12) {\ldots};
\node[popdark] at (11.1,-1.12) {\ldots};
\node[popdark] at (13.6,-1.12) {\ldots};
\foreach \x in {2.3,4.9,7.3,9.9,12.3} {\draw (\x,0.75) -- (\x,-1.5);}
\end{tikzpicture}
\end{center}
\legende{Grille d'analyse d'un texte philosophique : à remplir au brouillon avant toute rédaction, en suivant le texte ligne par ligne.}
\end{popfigure}

\courssub{Entraînement sur corpus de textes}

Voici un corpus représentatif des textes qui peuvent tomber à l'examen. Pour chacun, appliquez la grille d'analyse ci-dessus.

\begin{itemize}
\item \textbf{Aristote} (\emph{Éthique à Nicomaque}, livre V) : la justice distributive répartit honneurs et richesses \emph{proportionnellement} au mérite ; la justice corrective rétablit l'égalité brisée par la faute. Question : l'égalité proportionnelle est-elle compatible avec l'égalité des droits ?
\item \textbf{Kant} (\emph{Métaphysique des mœurs}, \emph{Doctrine du droit}, 1797) : le droit est la liberté de chacun soumise à la condition de la loi universelle. Question : la loi limite-t-elle la liberté ou la rend-elle possible ?
\item \textbf{Rawls} (\emph{Théorie de la justice}, 1971) : derrière le « voile d'ignorance », des personnes qui ignorent leur future place sociale choisiraient des principes justes. Question : cette fiction est-elle un bon critère de justice ?
\item \textbf{Eboussi Boulaga} : en Afrique, le droit hérité de la colonisation coexiste avec les coutumes ; la question est celle de leur articulation légitime. Question : le droit positif étatique peut-il ignorer le droit coutumier ?
\item \textbf{Jurisprudence camerounaise} (femme et héritage) : des décisions de justice camerounaises ont écarté des coutumes successorales discriminatoires au nom des droits fondamentaux et des engagements internationaux du Cameroun. Question : le juge peut-il corriger la coutume au nom de la justice ?
\end{itemize}

\begin{popfigure}
\begin{center}
\begin{tikzpicture}[mindmap, grow cyclic, every node/.style={concept}, root concept/.append style={concept color=popblue, font=\small\bfseries},
level 1/.append style={level distance=3.4cm, sibling angle=90},
level 2/.append style={level distance=2.4cm, sibling angle=50, font=\scriptsize}]
\node[concept color=popblue, text=white]{Mots-opérateurs du sujet}
  child[concept color=popgreenL] { node {\textbf{Suffit-il ?}}
      child { node {Oui : condition suffisante} }
      child { node {Non : nécessaire mais insuffisant} } }
  child[concept color=poporangeL] { node {\textbf{Faut-il ?}}
      child { node {Obligation morale} }
      child { node {Obligation juridique} } }
  child[concept color=poppurpleL] { node {\textbf{Peut-on ?}}
      child { node {Possibilité de fait} }
      child { node {Légitimité de droit} } }
  child[concept color=popgoldL] { node {\textbf{Est-ce que ?}}
      child { node {Question de vérité} }
      child { node {Question de définition} } };
\end{tikzpicture}
\end{center}
\legende{Carte mentale des mots-opérateurs : identifier l'opérateur du sujet, c'est savoir quel type de réponse l'examinateur attend.}
\end{popfigure}

\courssub{Évaluation formative : le barème de l'examinateur}

\begin{exemplebox}[Barème type au Baccalauréat technique (coefficient 2 à 4 selon la série)]
\textbf{Dissertation /20} : Introduction (3 pts) + Développement (12 pts) + Conclusion (2 pts) + Présentation et style (3 pts). \textbf{Explication de texte /20} : Présentation (3 pts) + Analyse (10 pts) + Discussion critique (5 pts) + Conclusion (2 pts). Grille de critères : Compréhension du sujet ou du texte (4 pts), Analyse et problématisation (6 pts), Argumentation et références (6 pts), Expression et orthographe (4 pts).
\end{exemplebox}

\begin{exoresolu}[Sujet d'entraînement : rédiger une introduction complète]
Énoncé : rédigez l'introduction de la dissertation « Le droit positif suffit-il à garantir la justice ? » en suivant la méthode de l'entonnoir.
\tcblower
\textbf{Solution détaillée.} \emph{Accroche} : « "Une loi injuste ne semble pas être une loi", écrivait saint Augustin, formulant une intuition que toute l'histoire du droit n'a cessé de confronter à la réalité des États. » \emph{Définitions} : le droit positif désigne l'ensemble des règles juridiques effectivement en vigueur dans un État ; garantir signifie assurer avec certitude ; la justice est l'idéal d'équité qui rend à chacun ce qui lui est dû. \emph{Reformulation} : la question est de savoir si la seule existence de lois valides suffit à rendre une société juste. \emph{Problématique} : mais alors, si le droit positif assure la sécurité et l'égalité formelle, comment expliquer que des lois formellement valides aient pu instituer l'injustice, et que des règles égales produisent des effets inégaux ? \emph{Annonce de plan} : nous montrerons d'abord que le droit positif est la condition nécessaire de toute justice effective ; nous verrons ensuite qu'il ne saurait en être la condition suffisante ; enfin, nous soutiendrons que la justice demeure l'horizon critique et la finalité du droit.
\end{exoresolu}

\begin{exercice}[Entraînement à l'explication de texte]
Prenez l'extrait de Rawls sur le voile d'ignorance fourni par votre professeur. 1) Remplissez la grille d'analyse (idée principale, concept clé, argument, exemple). 2) Rédigez la présentation du texte en cinq lignes. 3) Formulez une objection : le voile d'ignorance est-il réaliste ? 4) Échangez votre copie avec un camarade et évaluez-la selon la grille critériée (compréhension /4, analyse /6, argumentation /6, expression /4) : c'est l'évaluation par les pairs.
\end{exercice}

\begin{corrige}[Exercice — Éléments de correction attendus]
1) Idée principale : des principes de justice sont justes s'ils sont choisis par des personnes ignorant leur future position sociale. Concept clé : voile d'ignorance (situation originelle). Argument : l'ignorance de son propre intérêt garantit l'impartialité du choix. Exemple : nul ne choisirait l'esclavage s'il risquait d'être l'esclave. 2) La présentation doit mentionner : John Rawls, \emph{Théorie de la justice} (1971), philosophe américain, contexte des luttes pour les droits civiques, thèse : la justice comme équité. 3) Objection possible : on ne peut pas se dépouiller de son identité ; mais on peut répondre que le voile est une \emph{fiction méthodologique}, non une expérience réelle. 4) L'évaluation par les pairs doit justifier chaque point attribué par une observation précise (ex. : « problématique absente : -2 en analyse »).
\end{corrige}

\begin{aretenir}[L'essentiel de la méthode]
Un sujet de philosophie se traite en quatre gestes : \textbf{définir} les mots-clés, \textbf{problématiser} (dégager la tension), \textbf{argumenter} en trois moments dialectiques avec auteurs et exemples camerounais, \textbf{conclure} en répondant clairement. En explication de texte : présenter, analyser dans l'ordre du texte, discuter, conclure. La copie gagnante est celle qui pense, non celle qui récite.
\end{aretenir}

\newpage

\coursec{Activité d'intégration}

\coursec{Le Droit et la Justice — Activité d'intégration : Simulation d'audience}

\courssub{Objectifs de l'activité}

Cette activité d'intégration mobilise l'ensemble des savoirs de la leçon \emph{Le Droit et la Justice} dans une mise en situation collaborative et réaliste : une simulation d'audience judiciaire portant sur un conflit successoral entre droit coutumier et droit positif moderne. À l'issue de l'activité, l'élève doit être capable de :

\begin{itemize}
\item définir et distinguer les concepts de \cle{droit}, de \cle{justice}, de \cle{loi} et de \cle{morale} ;
\item identifier les fondements du droit (droit positif, droit naturel, justice comme valeur) et repérer un conflit de normes ;
\item appliquer la typologie des formes de justice (distributive, commutative, corrective) à un cas concret ;
\item rédiger et justifier une décision argumentée selon la structure du \cle{syllogisme judiciaire} ;
\item discuter philosophiquement la question : \emph{la loi suffit-elle à rendre justice ?}
\end{itemize}

\begin{attention}[Durée et organisation]
L'activité se déroule sur 2 heures en 4 phases : analyse juridique en groupes (30 min), simulation d'audience (45 min), délibération écrite individuelle (30 min), mise en commun et débat philosophique (15 min). Chaque élève doit produire une décision écrite individuelle, même si la simulation est collective.
\end{attention}

\courssub{Situation}

\textbf{Simulation d'audience : « L'héritage de M. Essomba »}

M. Essomba Martin, ressortissant de la communauté Bassa, installé à Yaoundé depuis trente ans, décède brutalement sans avoir rédigé de testament. Il laisse derrière lui :

\begin{itemize}
\item son épouse, Mme Essomba Charlotte, mariée sous le régime du \textbf{mariage civil} à la mairie de Yaoundé 5\textsuperscript{e} ;
\item deux filles, Aïcha (24 ans, étudiante) et Sandrine (19 ans, élève en Terminale) ;
\item un fils cadet, Junior (15 ans) ;
\item son neveu aîné, Nkoulou André (38 ans), désigné par la famille élargie comme successeur coutumier.
\end{itemize}

Le patrimoine comprend : une maison à Yaoundé évaluée à \num{45000000} FCFA, un terrain familial de 2 hectares au village (siège de la chefferie), un compte bancaire de \num{8000000} FCFA et un commerce.

\textbf{Le conflit.} Selon la coutume Bassa invoquée par le neveu et le chef de quartier, seul un héritier mâle de la lignée paternelle peut succéder à la tête de la concession familiale et aux terres ; les filles et l'épouse sont exclues de la succession foncière. Selon le droit positif camerounais, le mariage civil soumet la succession au Code civil : les enfants — filles et garçons sans distinction — et le conjoint survivant sont héritiers réservataires.

\textbf{Extraits du dossier (fourni aux élèves) :}
\begin{itemize}
\item Article 745 du Code civil : les enfants succèdent à parts égales, sans distinction de sexe ni de primogéniture ;
\item Article 1\textsuperscript{er} de la Constitution de 1996 : « la loi assure à tous le droit à l'égalité » ;
\item Convention CEDAW (élimination des discriminations à l'égard des femmes), ratifiée par le Cameroun ;
\item Avis consultatif de la Cour suprême (2017) : les coutumes contraires à l'ordre public et à l'égalité ne peuvent être appliquées ;
\item Témoignages contradictoires de l'épouse, du neveu, du chef de quartier et de l'avocat ;
\item Extrait de John Rawls, \emph{Théorie de la justice} : le « voile d'ignorance ».
\end{itemize}

\textbf{Problématique de l'activité :} quelle décision le juge doit-il rendre pour que le droit soit dit \emph{et} que justice soit faite ?

\courssub{Consignes}

\textbf{Phase 1 — Analyse juridique (groupes de 4, 30 min).}
\begin{enumerate}
\item Identifiez toutes les normes en présence (coutume, Code civil, Constitution, convention internationale) et classez-les selon la hiérarchie des normes.
\item Formulez précisément le conflit de normes : quelles règles s'opposent, et sur quel point ?
\item Ce litige relève-t-il de la justice distributive, de la justice commutative ou de la justice corrective ? Justifiez.
\end{enumerate}

\textbf{Phase 2 — Simulation d'audience (45 min).} Attribution des rôles : Juge, Médiateur, Avocat de l'épouse (droit moderne et droits des femmes), Avocat du neveu (coutume et tradition), Observateur (grille d'évaluation).
\begin{enumerate}
\item Chaque avocat rédige puis plaide un argumentaire de 5 minutes fondé sur un fondement du droit (droit positif pour l'un, droit naturel et coutume pour l'autre).
\item Le médiateur propose une solution d'équité : que suggère-t-il et pourquoi ?
\end{enumerate}

\textbf{Phase 3 — Délibération écrite individuelle (30 min).}
\begin{enumerate}
\item Rédigez la décision motivée du Juge (2 pages maximum) selon la structure du syllogisme judiciaire : la majeure (la règle de droit applicable), la mineure (les faits qualifiés), la conclusion (le dispositif : qui hérite de quoi).
\end{enumerate}

\textbf{Phase 4 — Débat philosophique (15 min).}
\begin{enumerate}
\item La loi (Code civil) a-t-elle suffi à « rendre justice » dans cette affaire ?
\item La médiation coutumière aurait-elle été plus « juste » ?
\item En vous plaçant sous le « voile d'ignorance » de Rawls (vous ignorez si vous naîtrez fille, garçon, épouse ou neveu), quelle règle successorale choisiriez-vous ? Concluez.
\end{enumerate}

\courssub{Ressources à mobiliser}

\begin{aretenir}[Notions utiles à réactiver]
\begin{itemize}
\item \cle{Droit positif} : ensemble des règles juridiques en vigueur dans un État à un moment donné (Constitution, lois, Code civil).
\item \cle{Droit naturel} : ensemble de principes de justice supposés universels et antérieurs aux lois humaines (égalité, dignité).
\item \cle{Hiérarchie des normes} : Constitution $>$ traités ratifiés $>$ lois $>$ règlements $>$ coutumes (recevables seulement si conformes à la loi et à l'ordre public).
\item \cle{Justice distributive} (Aristote) : à chacun selon son mérite ou ses droits, proportionnellement ; \cle{justice commutative} : égalité arithmétique dans les échanges ; \cle{justice corrective} : réparation du tort par le juge.
\item \cle{Syllogisme judiciaire} : majeure (la loi) + mineure (les faits) = conclusion (la décision).
\item \cle{Voile d'ignorance} (Rawls) : les principes justes sont ceux que choisiraient des personnes ignorant leur place future dans la société.
\end{itemize}
\end{aretenir}

\courssub{Démarche et corrigé commenté}

\begin{exoresolu}[Corrigé commenté de la simulation]
\textbf{Étape 1 — Identification et hiérarchie des normes.} Quatre normes sont en présence : la coutume Bassa, le Code civil (art. 745), la Constitution (art. 1\textsuperscript{er}, égalité) et la CEDAW ratifiée. Selon la hiérarchie des normes, la Constitution prime les lois, lesquelles priment les coutumes. Une coutume n'est applicable que si elle n'est pas contraire à la loi, à l'ordre public et aux engagements internationaux du Cameroun. Or la coutume Bassa exclut les filles et l'épouse : elle contredit l'article 745 du Code civil et le principe constitutionnel d'égalité. \emph{Elle est donc irrecevable devant le juge étatique.}

\textbf{Étape 2 — Qualification du type de justice.} Le litige porte sur le partage d'un patrimoine entre plusieurs ayants droit : il relève principalement de la \emph{justice distributive} (répartir les biens selon les droits de chacun). Le jugement qui tranche le conflit relève de la \emph{justice corrective} (le juge répare l'atteinte aux droits de l'épouse et des filles). La justice commutative n'est qu'accessoire (aucun échange contractuel en cause).

\textbf{Étape 3 — Plaidoiries.} L'avocat de l'épouse fonde son argumentation sur le droit positif : mariage civil, article 745, égalité constitutionnelle, CEDAW. Il peut ajouter un argument de droit naturel : l'égalité de dignité entre filles et garçons est un principe universel de justice. L'avocat du neveu invoque la tradition, la continuité de la lignée et la fonction sociale de la chefferie : argument respectable sur le plan culturel, mais juridiquement subordonné à la loi, et philosophiquement critiquable car il sacrifie l'égalité des personnes à la coutume du groupe.

\textbf{Étape 4 — Décision du Juge (syllogisme judiciaire).}
\begin{itemize}
\item \textbf{Majeure :} selon l'article 745 du Code civil, applicable au défunt marié civilement, les enfants succèdent à parts égales sans distinction de sexe, et le conjoint survivant recueille sa part ; toute coutume contraire à l'égalité (Const., art. 1\textsuperscript{er}) est irrecevable.
\item \textbf{Mineure :} M. Essomba était marié civilement ; il laisse une épouse et trois enfants ; la coutume invoquée exclut les filles et l'épouse.
\item \textbf{Conclusion (dispositif) :} la succession est répartie entre l'épouse et les trois enfants selon le Code civil ; le neveu est débouté ; la chefferie, fonction non patrimoniale, peut être transmise selon la coutume, mais sans les terres litigieuses.
\end{itemize}

\textbf{Étape 5 — Discussion philosophique.} La loi a rendu justice au sens de l'égalité formelle, mais on peut discuter son \emph{effectivité} : menaces familiales, pressions sociales, coût de la procédure. La médiation aurait pu préserver la paix familiale (justice comme équité et harmonie), mais au risque de faire payer l'accord par le renoncement des plus faibles. Sous le voile d'ignorance de Rawls, chacun choisirait la règle égalitaire du Code civil, car nul n'accepterait rationnellement une règle qui le déshérite s'il naît fille ou épouse. La justice comme valeur dépasse donc la seule coutume : elle exige des principes universalisables.
\tcblower
\textbf{Commentaires méthodologiques.} (1) Ne jamais opposer « droit » et « justice » de façon manichéenne : ici le droit positif coïncide avec l'exigence de justice (égalité), mais le cours a montré que ce n'est pas toujours le cas (lois injustes). (2) Dans la rédaction, chaque affirmation juridique doit être appuyée par une référence précise (article, principe). (3) La nuance de la décision — distinguer la fonction de chef (transmissible par la coutume) du patrimoine foncier (soumis au Code civil) — illustre la différence entre la lettre de la loi et l'esprit d'équité, sans trahir la hiérarchie des normes. (4) Au débat, valoriser les élèves qui mobilisent Aristote (proportionnalité) et Rawls (universalisation) plutôt que des opinions personnelles.
\end{exoresolu}

\courssub{Bilan de l'activité}

Cette simulation a permis de mettre en évidence trois acquis essentiels. D'abord, le \cle{droit positif} n'est pas une donnée abstraite : il tranche des conflits réels de la vie camerounaise, et son application suppose de maîtriser la hiérarchie des normes (la coutume, respectable, demeure subordonnée à la Constitution et à la loi). Ensuite, la \cle{justice} apparaît comme une valeur qui dépasse la légalité : le débat final montre qu'appliquer la loi ne suffit pas toujours à apaiser, et que l'équité, la médiation et l'effectivité des droits posent des problèmes proprement philosophiques. Enfin, sur le plan méthodologique, l'élève a pratiqué la démarche complète exigée au Baccalauréat technique : analyser une situation, mobiliser des concepts et des auteurs (Aristote, Rawls), construire un raisonnement rigoureux (syllogisme judiciaire) et rédiger une conclusion justifiée — compétences directement transférables à la dissertation et à l'explication de texte.

\newpage

\coursec{Méthodes et exercices résolus}

\courssub{Méthodes et Exercices Résolus}

\begin{methode}[Analyser une situation concrète au prisme des concepts : Droit, Justice, Loi, Morale]
    \textbf{Objectif :} Appliquer les définitions philosophiques pour qualifier un fait social ou juridique.
    \begin{enumerate}
        \item \textbf{Identifier les faits :} Relever les actions, les règles invoquées, les acteurs et le contexte (camerounais de préférence).
        \item \textbf{Qualifier chaque élément :} 
            \begin{itemize}
                \item \cle{Droit} : Ensemble de règles obligatoires sanctionnées par la force publique (État).
                \item \cle{Loi} : Règle générale, abstraite, impersonnelle, votée par le Parlement.
                \item \cle{Morale} : Ensemble de règles intérieures, autonomes, sanctionnées par la conscience/opinion.
                \item \cle{Justice} : Idéal d'équité, de proportionnalité, valeur normative jugeant le Droit.
            \end{itemize}
        \item \textbf{Repérer les tensions :} Conflit Loi/Morale (légal mais illégitime ?), Droit positif/Justice (légal mais injuste ?), Justice formelle/Justice réelle.
        \item \textbf{Conclure philosophiquement :} La situation relève-t-elle du \emph{positivisme juridique} (le droit c'est la loi) ou du \cle{jusnaturalisme} (un droit injuste n'est pas du droit) ?
    \end{enumerate}
\end{methode}

\begin{exoresolu}[Qualification d'une situation : L'expulsion d'un quartier informel à Douala]
    \textbf{Énoncé :} La mairie de Douala fait démolir des habitations précaires construites sans titre foncier sur un domaine public, en application d'un arrêté préfectoral. Les habitants, installés depuis 15 ans, dénoncent une injustice : « Nous n'avons nulle part où aller, la loi nous écrase ». Un magistrat affirme : « La procédure est régulière, le droit est respecté ». 
    \textit{Analysez cette situation en distinguant le Droit, la Loi, la Morale et la Justice.}
    \tcblower
    \textbf{Solution détaillée :}
    \begin{enumerate}
        \item \textbf{Le Droit positif (Légalité formelle) :} L'arrêté préfectoral est une source de droit (règlement). La démolition est \cle{légale} : elle s'appuie sur une règle de droit (code de l'urbanisme, domaine public) et est exécutée par la force publique (police municipale). Le magistrat a raison sur le plan \emph{juridique strict} (positivisme kelsenien/hartien : la validité du droit ne dépend pas de sa valeur morale).
        \item \textbf{La Loi :} Si l'arrêté s'appuie sur une loi votée par le Parlement (ex: Loi sur la décentralisation ou code foncier), c'est l'application d'une \cle{loi} (générale, abstraite). Elle ne vise pas ces habitants précis, mais « tout occupant sans titre ».
        \item \textbf{La Morale (Légitimité subjective) :} Les habitants invoquent la \cle{morale} : l'exigence de dignité, le droit au logement (DUDH art. 25), la proportionnalité de la peine (perdre son toit pour une faute administrative). La sanction morale est la honte, l'indignation publique, le sentiment d'inhumanité. C'est une obligation \emph{autonome} (conscience), non hétéronome comme le droit.
        \item \textbf{La Justice (Équité / Valeur) :} Ici se joue le conflit \cle{Droit positif vs Justice}. 
            \begin{itemize}
                \item \emph{Justice commutative / légale :} Respect de la procédure (le juge dit : « le droit est respecté »).
                \item \emph{Justice distributive / sociale :} Inégalité de traitement (les riches régularisent, les pauvres sont expulsés). 
                \item \emph{Justice comme valeur (Aristote/Rawls) :} La loi appliquée aveuglément produit de l'injustice (\emph{summum jus, summa injuria}). La justice exige ici une \cle{équité} (correctif de la loi générale) : relogement, délai, indemnisation.
            \end{itemize}
        \item \textbf{Synthèse philosophique :} La situation illustre la thèse du \cle{jusnaturalisme} (Antigone, Thomas d'Aquin, Dworkin) : « \emph{Lex injusta non est lex} » (Une loi injuste n'est pas une loi). Le Droit positif est séparé de la Justice idéale. L'État de droit formel (procédure respectée) ne garantit pas l'État de droit matériel (respect des droits fondamentaux).
    \end{enumerate}
    \textbf{Conclusion :} L'expulsion est \textbf{légale} (conforme au Droit positif/Loi) mais \textbf{injuste} (contraire à la Justice comme équité et à la Morale sociale). Le défi camerounais est le passage de la \cle{légalité} à la \cle{légitimité}.
\end{exoresolu}

\begin{methode}[Distinguer et articuler Droit positif et Justice (Fondements du Droit)]
    \textbf{Objectif :} Maîtriser l'opposition classique (Positivisme vs Jusnaturalisme) et la synthèse moderne (État de droit, contrôle de constitutionnalité).
    \begin{enumerate}
        \item \textbf{Définir les deux pôles :}
            \begin{itemize}
                \item \cle{Positivisme juridique (Kelsen, Hart)} : Le droit = ce qui est posé (loi, coutume, jurisprudence). Validité = respect de la \cle{Grundnorm} (norme fondamentale) ou règle de reconnaissance. Séparation stricte \emph{Fait/Valeur}, \emph{Être/Devoir}. « Le droit est le droit, la morale est la morale ».
                \item \cle{Jusnaturalisme (Aristote, Stoïciens, Thomas d'Aquin, Locke, Dworkin)} : Existence d'un \cle{Droit Naturel} antérieur et supérieur au droit positif. Critères : nature humaine, raison, Dieu, droits de l'homme. Un droit positif injuste est invalide (\emph{lex injusta non est lex}).
            \end{itemize}
        \item \textbf{Identifier la thèse de l'auteur/du sujet :} Le sujet penche-t-il vers « Le droit c'est la loi » (positivisme) ou « Une loi injuste n'est pas une loi » (jusnaturalisme) ?
        \item \textbf{Articuler (Synthèse Bac) :} 
            \begin{itemize}
                \item Le Droit positif est \emph{nécessaire} (sécurité juridique, prévisibilité, paix civile).
                \item La Justice (Droit naturel/Droits de l'homme) est \emph{l'horizon normatif} (légitimité, contrôle de constitutionnalité, droits fondamentaux).
                \item \textbf{Institutionnalisation :} Conseil Constitutionnel (Cameroun), Cour Suprême, CEDH, CADHP. Le juge devient le gardien de la Justice contre la Loi.
            \end{itemize}
        \item \textbf{Exemple camerounais :} La Constitution de 1996 (Préambule) intègre les droits de l'homme \emph{au sommet de la hiérarchie des normes}. Le juge constitutionnel peut censurer la loi au nom de la Justice.
    \end{enumerate}
\end{methode}

\begin{exoresolu}[Dissertation type : « Le Droit positif suffit-il à réaliser la Justice ? »]
    \textbf{Énoncé :} \textit{Sujet de dissertation : « Le Droit positif suffit-il à réaliser la Justice ? »}
    \tcblower
    \textbf{Solution détaillée (Plan dialectique structuré) :}
    \textbf{Introduction :}
    \begin{itemize}
        \item \textbf{Accroche :} Citation de Montesquieu (\emph{De l'esprit des lois}) : « Il n'y a point de plus grand tyrannie que celle qui se perpétue à l'ombre des lois et avec les couleurs de la justice. »
        \item \textbf{Définitions :} Droit positif = ensemble des règles en vigueur dans un État (légalité formelle). Justice = valeur idéale d'équité, de respect des droits fondamentaux (légitimité matérielle).
        \item \textbf{Problématique :} La conformité à la loi (légalité) garantit-elle l'équité (justice) ? Le risque de \emph{summum jus, summa injuria}.
        \item \textbf{Annonce du plan :} 1. Le Droit positif comme condition nécessaire de la Justice (sécurité, égalité formelle). 2. L'insuffisance du Droit positif face à l'exigence de Justice (injustices légales, rigidité). 3. L'institutionnalisation de la Justice comme contrôle du Droit positif (État de droit matériel).
    \end{itemize}
    \textbf{Développement :}
    \paragraph{I. Le Droit positif : condition nécessaire de la Justice (Thèse positiviste / Kelsen, Hart).}
    \begin{itemize}
        \item \textbf{Argument 1 : Sécurité juridique et paix civile.} Sans règle posée, claire, générale, c'est l'arbitraire (loi du plus fort). Le Droit positif assure l'\cle{égalité formelle} devant la loi (Art. 1er Constitution Cameroun). Ex: Code pénal : même peine pour même délit, peu importe le statut social.
        \item \textbf{Argument 2 : Séparation des pouvoirs / Légitimité démocratique.} La loi émane du Parlement (représentant le peuple). Obéir à la loi, c'est obéir à soi-même (Rousseau). Le juge \emph{ne fait que} dire le droit (bouche de la loi), il ne le crée pas. Cela évite le gouvernement des juges.
    \end{itemize}
    \paragraph{II. L'insuffisance du Droit positif : la légalité peut être injuste (Antithèse jusnaturaliste / Aristote, Antigone, Radbruch, Dworkin).}
    \begin{itemize}
        \item \textbf{Argument 1 : La rigidité de la loi générale.} La loi est abstraite, le cas est concret. Aristote (\emph{Éthique à Nicomaque}, V, 10) : l'équité (\cle{epieikeia}) corrige la loi là où elle est défectueuse par sa généralité. Ex camerounais : Peine de prison ferme pour vol de nourriture par nécessité (état de nécessité moral non prévu par le code pénal rigide).
        \item \textbf{Argument 2 : Les lois injustes de l'Histoire.} Lois raciales (Vichy, Apartheid), lois totalitaires. Elles étaient \emph{légales} (positives) mais radicalement \emph{injustes}. Thèse de Radbruch (1946) : « Le droit positif qui nie les droits de l'homme n'est pas du droit ». Au Cameroun : Décrets d'exception ou lois liberticides peuvent être constitutionnels formellement mais injustes matériellement.
        \item \textbf{Argument 3 : Justice sociale vs Justice commutative.} Le Droit positif sanctuarise souvent la propriété (Code civil). La Justice distributive (Rawls, \emph{Théorie de la justice}) exige de corriger les inégalités de départ (fiscalité progressive, protection sociale) que le droit positif pur ignore.
    \end{itemize}
    \paragraph{III. Vers un État de droit matériel : La Justice comme norme supérieure de contrôle (Synthèse / Kelsen révisé, Dworkin, Conseil Constitutionnel).}
    \begin{itemize}
        \item \textbf{Argument 1 : Hiérarchie des normes et Bloc de constitutionnalité.} Au Cameroun, le Préambule de la Constitution de 1996 intègre la Déclaration des Droits de l'Homme. Le Conseil Constitutionnel (Loi n°96/06) censure les lois contraires aux droits fondamentaux. La Justice (Droits de l'homme) devient une \cle{norme juridique suprême}.
        \item \textbf{Argument 2 : Le juge, gardien de la Justice.} Le juge n'est plus « bouche de la loi » (Montesquieu) mais interprète les principes généraux du droit, les droits fondamentaux. Ex: Jurisprudence du Conseil Constitutionnel camerounais sur la liberté de presse, l'égalité homme/femme.
        \item \textbf{Argument 3 : Effectivité de la Justice.} Le Droit positif n'est pas assez : il faut des institutions indépendantes (Justice), des auxiliaires de justice, l'aide juridictionnelle (accès au droit pour les pauvres). La Justice n'est pas seulement dans le texte, mais dans l'\cle{effectivité}.
    \end{itemize}
    \textbf{Conclusion :}
    \begin{itemize}
        \item \textbf{Bilan :} Le Droit positif est le \emph{corps} (structure, procédure), la Justice en est l'\emph{âme} (finalité, valeur). Le premier sans la seconde est une coquille vide (tyrannie légale) ; la seconde sans le premier est un vœu pieux (utopie).
        \item \textbf{Ouverture :} Le défi camerounais : passer de l'\emph{État de droit formel} (textes votés) à l'\emph{État de droit matériel} (indépendance réelle de la justice, effectivité des droits sociaux, lutte contre la corruption judiciaire).
    \end{itemize}
\end{exoresolu}

\begin{methode}[Maîtriser la typologie aristotélicienne de la Justice (Formes de la Justice)]
    \textbf{Objectif :} Identifier, définir et appliquer les trois formes de la justice chez Aristote (\emph{Éthique à Nicomaque}, Livre V) et leurs prolongements modernes.
    \begin{enumerate}
        \item \textbf{Justice Générale (Légale) :} Respect de la loi = \cle{Justice légale}. Synonyme de vertu complète par rapport aux autres. \textbf{Mot-clé :} \emph{Légalité}.
        \item \textbf{Justice Particulière (Égale / Proportionnelle) :} Rapport entre personnes et biens/honneurs. Deux espèces :
            \begin{itemize}
                \item \cle{Justice Commutative (Corrective) :} Échanges entre égaux (contrats, délits). \textbf{Égalité arithmétique} (A = B). Rétablir l'équilibre rompu (sanction, réparation). \emph{Mot-clé : Égalité stricte, réparation.}
                \item \cle{Justice Distributive :} Partage des honneurs/biens/charges selon le \cle{mérite} (valeur, besoin, contribution). \textbf{Proportionnalité géométrique} : $\frac{A}{B} = \frac{C}{D}$. \emph{Mot-clé : Mérite, Proportionnalité, Critère de répartition.}
            \end{itemize}
        \item \textbf{Justice comme Équité (\emph{Epieikeia}) :} Correction de la loi générale trop rigide pour le cas particulier. \emph{Mot-clé : Humanité, Circonstances, Discrétion judiciaire.}
        \item \textbf{Prolongements modernes (Rawls) :} 
            \begin{itemize}
                \item Principe de \cle{Liberté égale} (Justice commutative moderne : droits de base).
                \item Principe de \cle{Différence} (Justice distributive moderne : inégalités acceptées si profit aux plus défavorisés).
                \item \cle{Justice procédurale} (pure, parfaite, imparfaite) : Justice du processus vs Justice du résultat.
            \end{itemize}
    \end{enumerate}
\end{methode}

\begin{exoresolu}[Application de la typologie : Conflit foncier et attribution de bourses au Cameroun]
    \textbf{Énoncé :} Deux situations : 
    \textbf{A.} Un litige oppose M. Fouda (propriétaire terrien) à M. Essomba (locataire) pour non-paiement de loyer (3 ans d'arriérés). Le tribunal condamne Essomba à payer et à libérer les lieux. 
    \textbf{B.} Le Ministère de l'Enseignement Supérieur attribue les bourses d'excellence. Deux étudiants ont 16/20 de moyenne. L'un est fils de ministre, l'autre est orphelin, premier de sa famille à l'université. La bourse est donnée au fils du ministre « car son dossier est complet plus tôt ».
    \textit{Qualifiez le type de justice en jeu dans chaque situation (Aristote / Moderne) et dites si elle est réalisée.}
    \tcblower
    \textbf{Solution détaillée :}
    \paragraph{Situation A : Litige locatif (Commutative).}
    \begin{itemize}
        \item \textbf{Type :} \cle{Justice Commutative} (Aristote) / \cle{Justice corrective} (Moderne).
        \item \textbf{Raisonnement :} Relation bilatérale contractuelle entre deux sujets de droit égaux (contractants). Il y a eu rupture d'égalité (Essomba a joui du bien sans payer). Le juge doit \textbf{rétablir l'égalité arithmétique} : restitution du bien + paiement de la dette (valeur égale à la perte subie).
        \item \textbf{Réalisation :} \textbf{OUI.} Le jugement rétablit l'équilibre \emph{ex ante}. C'est une justice \emph{procédurale parfaite} (règle claire, juge impartial, résultat juste).
    \end{itemize}
    \paragraph{Situation B : Attribution des bourses (Distributive).}
    \begin{itemize}
        \item \textbf{Type :} \cle{Justice Distributive} (Aristote) / \cle{Justice sociale / Distributive} (Rawls).
        \item \textbf{Raisonnement :} Répartition d'un bien rare (bourse) entre des membres d'une communauté (étudiants) selon un \textbf{critère de mérite} (moyenne = 16/20 pour les deux). La proportionnalité géométrique exige : $\frac{\text{Bourse}_1}{\text{Mérite}_1} = \frac{\text{Bourse}_2}{\text{Mérite}_2}$. Comme Mérite$_1$ = Mérite$_2$, on doit avoir Bourse$_1$ = Bourse$_2$ (ou tirage au sort / critères secondaires objectifs : situation sociale).
        \item \textbf{Critère utilisé par l'administration :} « Dossier complet plus tôt » (critère bureaucratique arbitraire) + népotisme implicite (fils de ministre).
        \item \textbf{Réalisation :} \textbf{NON. Injustice flagrante.}
            \begin{enumerate}
                \item Violation de l'\cle{égalité proportionnelle} (Aristote) : même mérite $\neq$ même traitement.
                \item Violation de l'\cle{équité} (Rawls) : le critère social (orphelin, primo-universitaire) aurait dû primer (principe de différence : favoriser le plus défavorisé).
                \item Violation de la \cle{Justice procédurale} : opacité, favoritisme, corruption (népotisme).
            \end{enumerate}
    \end{itemize}
    \textbf{Synthèse :} La situation A relève de la \textbf{justice commutative} (réparation d'un tort entre égaux) et est respectée. La situation B relève de la \textbf{justice distributive} (répartition selon le mérite/bisoin) et est violée par l'arbitraire et le népotisme, transformant un droit (bourse au mérite) en privilège.
\end{exoresolu}

\begin{methode}[Réussir l'explication de texte philosophique (Méthode Bac)]
    \textbf{Objectif :} Structurer une explication de texte (Durée : 2h-3h au Bac) selon les attendus MINESEC : Comprendre, Analyser, Discuter.
    \begin{enumerate}
        \item \textbf{Étape 1 : Lecture active \& Paraphrase (15-20 min).} Lire 3 fois. Souligner : Thèse principale, Concepts clés, Arguments, Exemples, Structure (connecteurs logiques : \emph{donc, or, mais, en effet}). Faire une \textbf{paraphrase} (reformulation simple du sens global) en marge.
        \item \textbf{Étape 2 : Introduction rédigée (20 min).} 
            \begin{itemize}
                \item \textbf{Accroche / Contexte :} Auteur, Œuvre, Date, Contexte historique/philosophique.
                \item \textbf{Thèse :} Une phrase : « L'auteur soutient que... »
                \item \textbf{Problématique :} La question philosophique à laquelle le texte répond (souvent : Comment X est-il possible ? / Faut-il Y ? / Quel est le rapport entre A et B ?).
                \item \textbf{Annonce du plan :} Suivre la structure du texte (généralement 2 ou 3 parties).
            \end{itemize}
        \item \textbf{Étape 3 : Développement structuré (1h - 1h30).} \textbf{Une partie = Un argument principal du texte.}
            \begin{itemize}
                \item \textbf{Titre de partie} (reprise d'un mot-clé du texte).
                \item \textbf{Explication :} Reprendre le fil de l'auteur (\emph{L'auteur explique que...}, \emph{Il argumente en disant que...}). Citer des \textbf{courtes citations} (guillemets) intégrées dans la phrase.
                \item \textbf{Analyse / Commentaire :} Définir les concepts, donner la référence doctrinale (autre auteur), donner un \textbf{exemple concret (Cameroun si possible)}. \textbf{Ne pas faire de hors-sujet.}
            \end{itemize}
        \item \textbf{Étape 4 : Discussion critique (20-30 min).} 
            \begin{itemize}
                \item \textbf{Portée :} Pourquoi ce texte est-il important ? (Avancée conceptuelle).
                \item \textbf{Limites / Objections :} Une objection argumentée (auteur opposé, contre-exemple, distinction oubliée). \emph{Attention :} Pas de « je pense que », mais « On peut objecter que... », « Cette thèse se heurte à... ».
                \item \textbf{Dépassement :} Comment la philosophie actuelle a résolu/dépassé le débat.
            \end{itemize}
        \item \textbf{Étape 5 : Conclusion (10 min).} Reprise thèse + problématique + ouverture.
        \item \textbf{Étape 6 : Relecture (10 min).} Orthographe, syntaxe, cohérence, citations exactes.
    \end{enumerate}
\end{methode}

\begin{exoresolu}[Explication de texte guidée : Extrait de John Rawls, \emph{Théorie de la justice} (1971) - Justice comme équité]
    \textbf{Énoncé :} \textit{Texte : « La justice comme équité commence... par l'idée que les principes de justice sont ceux que choisiraient des personnes rationnelles placées dans une situation initiale d'égalité... Le voile d'ignorance assure que nul n'est avantagé ou désavantagé par le hasard de la naissance ou la contingence sociale... » (Rawls, §11-12, résumé).}
    \textit{Expliquez ce texte en en faisant ressortir la thèse, les arguments et la portée.}
    \tcblower
    \textbf{Solution détaillée (Modèle de rédaction Bac) :}
    \paragraph{Introduction}
    John Rawls (1921-2002), philosophe libéral américain, publie \emph{Théorie de la justice} (1971). Il relance la philosophie politique normative contre l'utilitarisme dominant. Le texte extrait des §11-12 pose les fondations de sa théorie : le \cle{contractualisme} moderne (héritier de Locke, Rousseau, Kant) via la fiction du \cle{voile d'ignorance}.
    \textbf{Thèse :} Les principes de justice ne sont pas donnés par la nature ni par l'utilité, mais choisis rationnellement par des égaux dans une situation initiale équitable.
    \textbf{Problématique :} Comment fonder des principes de justice impartiaux, indépendants des intérêts particuliers et des hasards sociaux ?
    \textbf{Plan :} 1. La situation initiale : une procédure de choix rationnel sous contrainte d'impartialité. 2. Le voile d'ignorance : l'instrument épistémique de l'équité. 3. Les principes qui en découlent : priorité de la liberté et principe de différence.
    \paragraph{I. La situation initiale : contractualisme et rationalité}
    Rawls reprend l'idée de \cle{contrat social} (État de nature), mais la « situation initiale » n'est pas un état historique. C'est une \textbf{dispositif de représentation} (heuristique). Les parties sont \textbf{rationnelles} (cherchent à maximiser leurs « biens premiers » : droits, libertés, revenus, bases sociales de l'estime de soi) et \textbf{mutuellement désintéressées} (pas d'altruisme, pas d'envie). Elles négocient les règles de la société future. L'\textbf{égalité} entre parties est postulée (symétrie).
    \paragraph{II. Le voile d'ignorance : condition de l'impartialité (Argument central)}
    C'est l'innovation majeure. Les parties ignorent : leur place sociale, leur fortune, leurs talents, leur conception du bien, leur psychologie, l'époque historique. Elles savent seulement : faits généraux (psychologie, économie, politique).
    \textbf{Analyse :} Ce voile épistémique force à l'\textbf{universalisation} (Kant : \emph{impératif catégorique}). Je ne peux pas choisir un principe favorisant les riches si je risque d'être pauvre. Cela élimine le \cle{hasard de la naissance} (aristocratie, caste) et la \cle{contingence sociale} (classe, sexe, race). L'équité (\emph{fairness}) est dans la \textbf{procédure} de choix, pas dans le résultat direct.
    \paragraph{III. Les principes de justice choisis (Portée normative)}
    Rawls démontre (théorème de la maximin) que des rationnels averses au risque choisiraient deux principes (dans l'ordre lexicographique) :
    1. \textbf{Principe de liberté égale :} Chaque personne a un droit égal au système le plus étendu de libertés fondamentales compatibles avec même liberté pour tous (Priorité du droit sur le bien).
    2. \textbf{Principe de différence (Justice distributive) :} Les inégalités sociales/économiques sont justes seulement si : (a) elles profitent aux plus défavorisés (maximin) ET (b) les postes sont accessibles à tous (égalité des chances équitable).
    \paragraph{Discussion critique}
    \textbf{Portée :} Fondement laïc, rationnel, universaliste des droits de l'homme. Base théorique de l'État social de droit (Cameroun : Préambule Constitution 1996). Réfutation de l'utilitarisme (sacrifice minorité pour majorité).
    \textbf{Limites (Objections) :}
    \begin{itemize}
        \item \emph{Communautariens (Sandel, Walzer, Taylor) :} Le « moi » derrière le voile est déraciné, abstrait, sans histoire ni appartenance. La justice ne se décide pas hors contexte communautaire.
        \item \emph{Libéraux libertaires (Nozick) :} Le principe de différence viole les droits de propriété (vol légal). Justice = respect des acquisitions/transferts justes (théorie dotation), pas structure de résultat.
        \item \emph{Féministes (Okin, Pateman) :} Le voile ignore le genre et la sphère privée (famille), lieu majeur d'injustice.
        \item \emph{Réalistes (Sen) :} Comparer des « biens premiers » ne suffit pas ; il faut regarder les \emph{capabilités} réelles (convertir des ressources en fonctionnements).
    \end{itemize}
    \textbf{Conclusion :} Rawls offre la théorie la plus robuste de la \cle{Justice comme équité} : une justice procédurale pure qui garantit des principes substantiels (Liberté, Différence). Elle reste la référence pour penser la justice constitutionnelle et sociale aujourd'hui, y compris au Cameroun où le « voile » devrait inspirer la répartition des ressources (budget, bourses, postes) entre régions et groupes sociaux.
\end{exoresolu}

\begin{methode}[Réussir la dissertation philosophique (Méthode Bac - 3h à 4h)]
    \textbf{Objectif :} Produire une réponse argumentée, structurée, problématisée à un sujet de dissertation.
    \begin{enumerate}
        \item \textbf{Analyse du sujet (30 min) :} 
            \begin{itemize}
                \item \textbf{Définir les termes} (Dictionnaire + Sens philosophique précis : Droit, Justice, Loi, État, Légitimité...).
                \item \textbf{Repérer la structure :} Sujet-thèse (affirmation), Sujet-question, Sujet-citation. Identifier la \textbf{tension} (opposition concepts : Droit vs Justice, Légalité vs Légitimité, Formel vs Matériel).
                \item \textbf{Formuler la problématique :} Question unique, ouverte, philosophique, qui contient la tension. \emph{Ex: « La conformité à la loi suffit-elle à garantir la justice ? »}
            \end{itemize}
        \item \textbf{Recherche d'idées / Plan (30 min) :} 
            \begin{itemize}
                \item Brainstorming : Thèses d'auteurs (Kelsen, Hart, Aristote, Rawls, Antigone, Montesquieu, Kouchner, Sen...), Concepts, Exemples (Cameroun, Histoire, Actualité), Distinctions.
                \item \textbf{Choix du plan :} \textbf{Dialectique (Thèse / Antithèse / Synthèse)} est le standard Bac. \textbf{Analytique (Concepts : Définition, Conditions, Conséquences)} si sujet conceptuel pur. \textbf{Progressif (Genèse, Structure, Finalité)} pour sujets métaphysiques.
                \item Rédiger le \textbf{plan détaillé} (Arguments, Auteurs, Exemples par sous-partie).
            \end{itemize}
        \item \textbf{Rédaction Introduction (20 min) :} Accroche $\rightarrow$ Définitions $\rightarrow$ Problématique $\rightarrow$ Annonce plan (3 parties, 2 sous-parties chacune).
        \item \textbf{Rédaction Développement (1h30 - 2h) :} 
            \begin{itemize}
                \item \textbf{1 paragraphe = 1 idée force = 1 argument.}
                \item Structure paragraphe : \textbf{Phrase thèse} $\rightarrow$ \textbf{Explication/Argumentation} (Auteur, Concept, Raisonnement) $\rightarrow$ \textbf{Exemple concret / Illustration} $\rightarrow$ \textbf{Mini-transition/Conclusion partielle}.
                \item \textbf{Transitions explicites} entre parties (reprise problématique + annonce suivante).
            \end{itemize}
        \item \textbf{Rédaction Conclusion (15 min) :} Réponse synthétique à la problématique $\rightarrow$ Nuance finale $\rightarrow$ Ouverture (autre concept, actualité, auteur).
        \item \textbf{Relecture (15 min) :} Cohérence, orthographe, respect du sujet, pas de hors-sujet.
    \end{enumerate}
\end{methode}

\begin{exoresolu}[Dissertation type Bac : « La Justice n'est-elle que le respect de la Loi ? »]
    \textbf{Énoncé :} \textit{Sujet : « La Justice n'est-elle que le respect de la Loi ? »}
    \tcblower
    \textbf{Solution détaillée (Plan dialectique + Rédaction clés) :}
    \paragraph{Introduction}
    \textbf{Accroche :} Sophocle, \emph{Antigone} : « Je n'ai pas cru que tes édits eussent assez de force pour que, mortel, tu puisses outre-passer les lois non écrites et immuables des dieux. »
    \textbf{Définitions :} Loi = règle générale, obligatoire, sanctionnée par l'État (légalité formelle). Justice = valeur idéale d'équité, de proportionnalité, respect des droits (légitimité matérielle).
    \textbf{Problématique :} L'obéissance à la loi (positivisme) épuise-t-elle l'exigence de justice, ou la justice exige-t-elle de juger la loi elle-même (jusnaturalisme / contrôle de constitutionnalité) ?
    \textbf{Plan :} I. Le respect de la loi comme condition formelle de la justice (Thèse positiviste). II. L'insuffisance de la légalité : la loi peut être injuste (Antithèse jusnaturaliste). III. La justice comme norme supérieure fondant la légitimité de la loi (Synthèse : État de droit matériel).
    \paragraph{I. Thèse : La loi comme cadre nécessaire de la justice (Légalité = Justice formelle)}
    \textbf{A. Sécurité juridique et égalité formelle (Kelsen, Hart, Montesquieu).} La loi générale et abstraite protège contre l'arbitraire (jugement au cas par cas, favoritisme). « Le juge est la bouche qui prononce les paroles de la loi » (Montesquieu). Au Cameroun : Code pénal, Code civil, Procédures = prévisibilité. Justice commutative = application égale de la loi à tous (Art. 1er Const.).
    \textbf{B. Légitimité démocratique (Rousseau, Kelsen).} La loi émane du souverain (Peuple/Parlement). Obéir à la loi = s'obéir à soi-même (autonomie collective). Remettre en cause la loi au nom de sa « justice » personnelle = risque de tyrannie judiciaire, d'arbitraire subjectif. Séparation des pouvoirs.
    \paragraph{II. Antithèse : La légalité ne suffit pas, la loi peut être injuste (Justice > Loi)}
    \textbf{A. La rigidité de la loi vs l'équité du cas particulier (Aristote, \emph{Epieikeia}).} Loi = universel. Cas = particulier. \emph{Summum jus, summa injuria}. Ex: Vol par nécessité (faim), légitime défense mal qualifiée, prescription extinctive qui empêche la vérité. La justice commutative exige la réparation, la loi peut l'interdire.
    \textbf{B. Les lois injustes de l'Histoire (Antigone, Radbruch, Nuremberg).} Lois raciales, lois vichystes, apartheid. Etaient \emph{légales} (votées, promulguées), mais \emph{criminelles}. Thèse de Radbruch (1946) : « Le droit positif qui contredit les principes de la justice... n'est pas du droit ». Principe de non-rétroactivité pénal vs Crimes contre l'humanité (Justice naturelle/Droits de l'homme).
    \textbf{C. Justice sociale vs Justice légale (Rawls, Marx, Sen).} La loi sanctuarise souvent les inégalités de fait (propriété, héritage). Justice distributive (Rawls : principe de différence) exige de \emph{corriger} la loi du marché par la fiscalité, la protection sociale. Au Cameroun : Loi foncière (1974) vs droits coutumiers / populations autochtones (Baka, Bororo). Légalité formelle $\neq$ Justice matérielle.
    \paragraph{III. Synthèse : La Justice comme fondement et contrôle de la Loi (État de droit constitutionnel)}
    \textbf{A. Hiérarchie des normes : Le bloc de constitutionnalité (Kelsen révisé, Conseil Constitutionnel).} La Constitution (au Cameroun : Préambule 1996 + CADHP) intègre la Justice (Droits de l'homme) au sommet. La loi doit être \emph{conforme} à la Justice pour être valide. Le juge constitutionnel censure la loi au nom de la Justice.
    \textbf{B. Le juge, gardien de la justice (Dworkin, Garantie des droits).} Le juge n'applique pas mécaniquement. Il interprète les « principes » (justice, équité) qui sous-tendent les « règles ». Ex: Jurisprudence camerounaise sur la liberté d'association, l'égalité homme/femme (coutume vs loi).
    \textbf{C. Effectivité : De la Justice sur le papier à la Justice vécue.} Respecter la loi ne suffit pas si la justice est inaccessible (coûts, délais, corruption, langue). Justice = \cle{Effectivité} (Accès au juge, Aide juridictionnelle, Indépendance magistrature).
    \paragraph{Conclusion}
    \textbf{Bilan :} La Justice n'est \textbf{pas seulement} le respect de la Loi. La Loi est l'instrument, la Justice est la fin. Le respect de la loi assure la \textbf{paix} (ordre), la Justice assure la \textbf{légitimité} (valeur).
    \textbf{Ouverture :} Le défi camerounais : \textbf{L'effectivité.} Comment faire en sorte que la loi (texte) devienne justice (réalité) pour le justiciable lambda ? Lutte contre la corruption judiciaire, formation des magistrats, vulgarisation du droit, justice de proximité (chefferies, médiateurs). La Justice n'est pas un texte, c'est une \emph{pratique institutionnelle}.
\end{exoresolu}

\begin{methode}[Analyser l'effectivité de la Justice au Cameroun (Contexte national - Savoir-faire « Situation concrète »)]
    \textbf{Objectif :} Mobiliser les connaissances institutionnelles (organisation judiciaire) et les concepts philosophiques pour diagnostiquer l'état de la justice au Cameroun.
    \begin{enumerate}
        \item \textbf{Connaître l'architecture institutionnelle (Texte vs Réalité) :}
            \begin{itemize}
                \item \textbf{Sommet :} Cour Suprême (Judiciaire, Administrative, Comptes), Conseil Constitutionnel.
                \item \textbf{Ordinaires :} Cours d'Appel, Tribunaux de Grande Instance (TGI), Tribunaux de Première Instance (TPI), Tribunaux Militaires.
                \item \textbf{Spécialisés/Alternatifs :} Tribunaux Coutumiers (Chefferies - Droit coutumier), Médiation/Conciliation (Maisons de Justice, ONG).
                \item \textbf{Garanties :} Indépendance (CSM - Conseil Supérieur de la Magistrature), Inamovibilité, Budget.
            \end{itemize}
        \item \textbf{Identifier les écarts (Effectivité) :} 
            \begin{itemize}
                \item \textbf{Accès :} Coût (avocat, frais), Distance (déserts judiciaires), Langue (Français/Anglais vs langues locales), Analphabétisme juridique.
                \item \textbf{Indépendance :} Pression exécutive (nomination, avancement via CSM), Corruption (magistrats, greffiers, OPJ), Instruction à charge.
                \item \textbf{Délais :} Engorgement, lenteur procédure, détention préventive abusive (violation DDH).
                \item \textbf{Pluralisme juridique :} Conflit Loi moderne / Coutume (femmes, successions, terres). Non-application jugements (force publique absente).
            \end{itemize}
        \item \textbf{Articuler Concepts / Réalité :} 
            \begin{itemize}
                \item \cle{Séparation des pouvoirs} vs \cle{Subordination} réelle.
                \item \cle{Droit positif} (textes beaux) vs \cle{Justice réelle} (vécu).
                \item \cle{Justice formelle} (procédure respectée) vs \cle{Justice substantielle} (droits effectifs).
            \end{itemize}
    \end{enumerate}
\end{methode}

\begin{exoresolu}[Analyse critique : « La Justice au Cameroun : entre textes et réalités »]
    \textbf{Énoncé :} \textit{« Le Cameroun dispose d'un arsenal juridique complet garantissant l'indépendance de la justice et les droits de l'homme. Pourtant, les justiciables dénoncent une justice à deux vitesses. » Analysez ce paradoxe en mobilisant les concepts de Droit positif, Justice, Légalité, Légitimité et Effectivité.}
    \tcblower
    \textbf{Solution détaillée :}
    \paragraph{1. Le Droit positif : Un arsenal normatif complet (Légalité formelle)}
    \begin{itemize}
        \item \textbf{Textes fondateurs :} Constitution 1996 (Préambule : DUDH 1948, Charte Africaine DHP 1981). Lois organiques sur l'organisation judiciaire (Loi 2006/015), CSM (Loi 2006/016), Aide juridictionnelle (Décret 2005).
        \item \textbf{Garanties institutionnelles :} Indépendance (Art. 37 Const : « Le pouvoir judiciaire est indépendant »), Inamovibilité, CSM (organe de gestion carrière), Cour Suprême (Cassation + Contentieux administratif + Comptes), Conseil Constitutionnel (Contrôle lois).
        \item \textbf{Ratifications :} PIDCP, PIDESC, CEDAW, CDE, CADHP (Protocole Cour Africaine).
    \end{itemize}
    \textbf{Conclusion partielle :} Sur le plan du \cle{Droit positif} (légalité), l'État de droit formel existe. La \cle{Loi} est belle.
    \paragraph{2. La Justice réelle : Déficit d'effectivité et de légitimité (Justice matérielle)}
    \begin{itemize}
        \item \textbf{Indépendance relative (Légitimité entamée) :} CSM présidé par le Président de la République (Ministre de la Justice VP). Risque d'injonctions / pressions sur carrières (affectations, promotions). Perception : « Justice aux ordres » pour les affaires politiques/sensibles.
        \item \textbf{Corruption systémique (Injustice commutative/distributive) :} « Justice à deux vitesses » : Riche/Puissant = Issue favorable (arrêt exécution, non-lieu, liberté provisoire) ; Pauvre = Détention préventive longue, condamnation lourde, exécution forcée. \emph{Ex:} Affaire « Opération Épervier » (détournements) vs vols simples. Corruption aux barreaux, greffes, OPJ.
        \item \textbf{Accès au droit inégal (Justice distributive violée) :} 
            \begin{itemize}
                \item \emph{Géographique :} TGI/TPI concentrés chefs-lieux. Zones rurales (Extrême-Nord, Est, Nord) = déserts judiciaires. Justice coutumière seule accessible (non garantie DDH femmes/enfants).
                \item \emph{Financier :} Frais d'avocat, timbre, expertise. Aide juridictionnelle quasi inexistante en pratique.
                \item \emph{Linguistique/Culturel :} Procès en français/anglais. Justiciables analphabètes ou locuteurs locaux sans interprètes assermentés.
            \end{itemize}
        \item \textbf{Délais et Détention préventive (Violence DDH) :} Engorgement (TGI Yaoundé/Mfoundi : milliers dossiers). Détention préventive > peine encourue (violation PIDCP Art. 9). Justice lente = Justice niée.
        \item \textbf{Pluralisme juridique non résolu :} Conflit Loi 1981 (Personnes/Famille) vs Coutume. Ex: Successions (filles exclues coutumièrement), Mariage forcé/précoce, Terre (droit coutumier vs titre foncier). L'État ne tranche pas (ou mal).
        \item \textbf{Non-exécution des décisions :} Arrêts Cour Suprême / TGI contre l'Administration non exécutés (violation séparation pouvoirs, Art. 37 Const). « Juger c'est rien, exécuter c'est tout » (Proudhon).
    \end{itemize}
    \paragraph{3. Synthèse philosophique : Le fossé Légalité / Légitimité}
    \begin{itemize}
        \item \textbf{Positivisme déconnecté :} Le Droit positif (textes) flotte au-dessus de la société. Il n'est pas \emph{intériorisé} par les acteurs (magistrats corrompus) ni \emph{accessible} aux sujets.
        \item \textbf{Crise de la \cle{Légitimité} (Weber) :} L'obéissance au droit ne repose plus sur la croyance en sa valeur (Justice), mais sur la contrainte (Police) ou l'intérêt (Corruption).
        \item \textbf{Justice comme \cle{Valeur} vs Justice comme \cle{Service public} :} La Justice n'est pas un service public comme les autres (eau, électricité). Elle est le \textbf{dernier rempart} de la dignité. Son dysfonctionnement détruit le \cle{Contrat social}.
    \end{itemize}
    \textbf{Conclusion :} Le paradoxe s'explique par la confusion entre \textbf{Droit positif} (l'avoir-textes) et \textbf{Justice effective} (l'être-vécu). Le Cameroun est un \emph{État de droit formel} (textes) mais peine à devenir un \emph{État de droit matériel} (effectivité). La réforme ne est pas textuelle (plus de lois) mais \textbf{structurelle} (indépendance réelle CSM, budget justice, lutte corruption, justice de proximité, formation, exécution décisions). La Justice n'est pas dans le Code, elle est dans le \textbf{jugement rendu et exécuté équitablement}.
\end{exoresolu}

\begin{methode}[Argumenter une position éthique/juridique (Cas pratique / Dilemme) ]
    \textbf{Objectif :} Répondre à une question de type « Faut-il... ? », « Est-il juste que... ? » en structurant un raisonnement éthique rigoureux (type Grand Oral / Épreuve orale / Question de cours).
    \begin{enumerate}
        \item \textbf{Reformuler le dilemme :} Identifier le conflit de valeurs (ex: Sécurité vs Liberté, Ordre vs Droit, Majorité vs Minorité, Loi vs Conscience).
        \item \textbf{Poser le cadre normatif :} Quelle est la règle de droit positive ? Quel est le principe de justice supérieur (Droits de l'homme, Constitution, Droit naturel) ?
        \item \textbf{Construire l'argumentation (Dialectique courte) :}
            \begin{itemize}
                \item \textbf{Argument 1 (Pour/Thèse) :} Appui sur le Droit positif, l'utilité sociale, l'ordre public, la sécurité. Citer texte/auteur.
                \item \textbf{Argument 2 (Contre/Antithèse) :} Appui sur la Justice, les droits fondamentaux, la proportionnalité, la dignité humaine. Citer texte/auteur/contre-exemple.
                \item \textbf{Argument 3 (Synthèse/Arbitrage) :} Critère de décision : \cle{Proportionnalité} (But légitime ? Moyens adaptés ? Nécessaires ? Proportionnés stricto sensu ?). Ou \cle{Principe de subsidiarité} / \cle{Équité}.
            \end{itemize}
        \item \textbf{Conclusion tranchée :} Prendre position clairement (« En définitive, il est juste/juste de ne pas... car... ») en assumant la nuance.
    \end{enumerate}
\end{methode}

\begin{exoresolu}[Cas pratique : « Faut-il désobéir à une loi injuste ? » (Antigone / Désobéissance civile)]
    \textbf{Énoncé :} \textit{Une loi votée par le Parlement camerounais interdit toute manifestation non autorisée 72h à l'avance, sous peine de prison ferme. Une ONG de défense des droits de l'homme organise une marche pacifique non déclarée pour protester contre la vie chère. Les militants sont arrêtés. L'ONG appelle à la « désobéissance civile ». Est-ce légitime ?}
    \tcblower
    \textbf{Solution détaillée :}
    \paragraph{1. Qualification juridique et éthique}
    \begin{itemize}
        \item \textbf{Fait :} Violation d'une loi formelle (Loi sur les manifestations / Code pénal). \textbf{Illégalité} avérée (Positivisme : Kelsen/Hart disent « c'est illégal »).
        \item \textbf{Motif :} Protestation contre vie chère (Droit à la vie, Dignité, Liberté d'expression/réunion - Préambule Const., PIDCP Art. 19, 21). Délai 72h = obstacle déraisonnable (urgence sociale).
        \item \textbf{Nature de l'acte :} \cle{Désobéissance civile} (Rawls, Arendt, King, Gandhi) : Publique, non-violente, consciente, acceptation de la peine, visée changement loi/politique. $\neq$ Désobéissance ordinaire (égoïste, cachée) ou Révolte (violente).
    \end{itemize}
    \paragraph{2. Argumentation dialectique}
    \textbf{Thèse : L'obligation d'obéissance (Légalité / Ordre public).}
    \begin{itemize}
        \item Argument : Pacte social (Rousseau/Hobbes). La loi est la volonté générale. Désobéir = retour à l'état de nature, anarchie, insécurité juridique. Si chacun choisit ses lois, plus de société.
        \item Voie de droit : Recours au Conseil Constitutionnel (inconstitutionnalité), manifestation déclarée, pétition, vote.
    \end{itemize}
    \textbf{Antithèse : La légitimité de la désobéissance face à l'injustice (Justice / Droits de l'homme).}
    \begin{itemize}
        \item Argument Jusnaturaliste (Antigone, Thomas d'Aquin, Luther King) : \emph{Lex injusta non est lex}. Une loi qui viole les droits fondamentaux (liberté réunion, proportionnalité) perd son caractère obligatoire \emph{en conscience}.
        \item Argument Rawlsien : Désobéissance civile = « acte de protestation adressé au sens de la justice de la majorité ». Condition : minorité bloquée, voies légales épuisées/ineffectives, non-violence.
        \item Argument Proportionnalité (Droit constitutionnel moderne / CEDH) : Restriction liberté réunion doit être \textbf{nécessaire dans une société démocratique} (besoin social impérieux) et \textbf{proportionnée}. Délai 72h + prison ferme pour marche pacifique vie chère = \textbf{disproportionné}. La loi est \emph{inconstitutionnelle} matériellement.
    \end{itemize}
    \textbf{Synthèse / Arbitrage :}
    \begin{itemize}
        \item La loi, bien que \emph{légale} (votée), est \emph{injuste} (atteinte disproportionnée à droit fondamental).
        \item La désobéissance civile est \textbf{légitime éthiquement} (Justice > Loi positive) et \textbf{juridiquement défendable} (Vice de constitutionnalité / Exception d'illégalité / État de nécessité moral).
        \item \textbf{Mais :} Elle doit respecter le \textbf{cahier des charges} (Rawls) : Non-violence stricte, Publicité, Acceptation procès (tribune pour dire le droit), Visée abrogation loi.
    \end{itemize}
    \paragraph{3. Conclusion}
    \textbf{OUI, la désobéissance civile est légitime ici.} Elle n'est pas un refus de la Loi en général, mais l'affirmation que la \cle{Justice} (Droits de l'homme, Proportionnalité) est la norme suprême qui juge la \cle{Loi}. Au Cameroun, l'Art. 65 Const. (révision) et le Préambule fondent ce contrôle de constitutionnalité matériel. Les militants ne sont pas des délinquants, mais des \emph{lanceurs d'alerte démocratiques}. Le juge, saisi, doit relaxer au nom de la hiérarchie des normes (Constitution > Loi simple) ou de l'état de nécessité.
\end{exoresolu}

\begin{attention}[Les 5 erreurs qui coûtent des points au Bac (Philosophie - Droit/Justice)]
    \begin{enumerate}
        \item \textbf{Confondre \cle{Droit positif} et \cle{Justice} (ou \cle{Loi} et \cle{Morale}).} 
            \textit{Erreur :} « C'est juste car c'est la loi. » ou « C'est illégal donc c'est immoral. »
            \textit{Correctif :} Distinguer systématiquement \emph{Validité juridique} (procédure, source) et \emph{Légitimité morale} (équité, droits humains). Utiliser les concepts : \textbf{Légalité vs Légitimité}, \textbf{Fait vs Valeur}.
        \item \textbf{Ignorer la typologie aristotélicienne (Commutative vs Distributive) dans les exemples.}
            \textit{Erreur :} Traiter un problème de répartition (bourses, impôts, terres) comme un problème de contrat/réparation (vol, accident), ou inversement.
            \textit{Correctif :} Identifier le \textbf{type de rapport} : Échange entre égaux (Commutative / Égalité arithmétique) OU Partage commun selon mérite/besoin (Distributive / Proportionnalité géométrique). Citer \textbf{Aristote}.
        \item \textbf{Parler de « Justice » au singulier sans précision (Justice formelle / matérielle, procédurale / substantielle).}
            \textit{Erreur :} « La justice est respectée » alors que la procédure l'est mais le résultat est inique (ou inverse).
            \textit{Correctif :} Préciser : \textbf{Justice formelle} (respect règles procédure/loi) vs \textbf{Justice matérielle/réelle} (équité du résultat, effectivité droits). \textbf{Justice procédurale} vs \textbf{Justice distributive}.
        \item \textbf{Absence d'exemples concrets camerounais (ou exemples stéréotypés « Antigone » seulement).}
            \textit{Erreur :} Dissertation abstraite, exemples uniquement grecs/français (Vichy, Nuremberg) sans ancrage local.
            \textit{Correctif :} Mobiliser \textbf{le contexte camerounais} : Constitution 1996, Conseil Constitutionnel, Cour Suprême, CSM, Droit coutumier (chefferies), Baka/Bororo, Opération Épervier, Vie chère, Détention préventive, Aide juridictionnelle, Langues officielles. Cela prouve la maîtrise du programme « Philosophie de l'enseignement technique » (ancrage territorial).
        \item \textbf{Mauvaise gestion du temps / Structure invisible (Pas de plan apparent, pas de transitions, pas de problématique).}
            \textit{Erreur :} « Brasse-cousin » d'idées, pas de thèse/antithèse/synthèse lisible, introduction sans problématique, conclusion absente.
            \textit{Correctif :} \textbf{Respecter le canevas méthodologique} (Analyse $\rightarrow$ Problématique $\rightarrow$ Plan annoncé $\rightarrow$ Parties équilibrées $\rightarrow$ Transitions $\rightarrow$ Conclusion/Ouverture). Utiliser des \textbf{connecteurs logiques} (Premièrement, En effet, Cependant, Par conséquent, En définitive). \textbf{Relire} 10 min.
    \end{enumerate}
\end{attention}

\newpage

\coursec{Exercices}

\courssub{Je vérifie mes connaissances}

\begin{exercice}[QCM : notions fondamentales]
Pour chaque question, une seule réponse est correcte. Entoure la lettre correspondante.
\begin{enumerate}
\item Selon Kelsen, le droit positif se définit comme :
\begin{enumerate}
\item l'ensemble des lois morales universelles ;
\item un ordre normatif coercitif, valable dans un État donné ;
\item la volonté générale exprimée par le peuple ;
\item la coutume la plus ancienne du territoire.
\end{enumerate}
\item La justice distributive, chez Aristote, consiste à :
\begin{enumerate}
\item réparer un tort subi par une victime ;
\item répartir les biens et les honneurs proportionnellement au mérite de chacun ;
\item appliquer strictement la loi écrite ;
\item punir les criminels.
\end{enumerate}
\item Au Cameroun, en cas de conflit entre une coutume et la Constitution :
\begin{enumerate}
\item la coutume l'emporte car elle est plus ancienne ;
\item la Constitution l'emporte car elle est la norme suprême ;
\item le chef traditionnel tranche ;
\item les deux normes s'annulent.
\end{enumerate}
\item Le « voile d'ignorance » de Rawls désigne :
\begin{enumerate}
\item l'ignorance des lois par les citoyens ;
\item une situation hypothétique où chacun ignore la place qu'il occupera dans la société ;
\item le secret des délibérations judiciaires ;
\item l'illétrisme juridique des populations rurales.
\end{enumerate}
\end{enumerate}
\end{exercice}

\begin{exercice}[Vrai ou faux justifié]
Indique si chaque affirmation est vraie ou fausse, puis justifie ta réponse en deux ou trois lignes.
\begin{enumerate}
\item « Toute loi votée par le Parlement est nécessairement juste. »
\item « La justice commutative concerne les échanges et les contrats entre particuliers. »
\item « Le Conseil Supérieur de la Magistrature (CSM) participe à la gestion de la carrière des magistrats au Cameroun. »
\item « Le droit naturel est un droit écrit, adopté par les assemblées. »
\end{enumerate}
\end{exercice}

\begin{exercice}[Définitions express]
Définis en trois lignes maximum chacune des notions suivantes, en donnant pour chacune un exemple tiré de la vie camerounaise :
\begin{enumerate}
\item le droit positif ;
\item la loi ;
\item la justice réparatrice ;
\item la morale.
\end{enumerate}
\end{exercice}

\begin{exercice}[Distinctions conceptuelles]
Recopie et complète le tableau suivant en distinguant les quatre notions selon leur origine, leur sanction et leur domaine d'application :
\begin{center}
\begin{tabularx}{\linewidth}{|l|X|X|X|}\hline
\rowcolor{popblueL} \textbf{Notion} & \textbf{Origine} & \textbf{Sanction} & \textbf{Domaine} \\ \hline
Droit & & & \\ \hline
Justice & & & \\ \hline
Loi & & & \\ \hline
Morale & & & \\ \hline
\end{tabularx}
\end{center}
\end{exercice}

\courssub{Je m'entraîne}

\begin{exercice}[Analyse de situations concrètes]
Pour chacune des situations suivantes, indique s'il s'agit d'un problème relevant de la justice commutative, de la justice distributive ou de la justice réparatrice, et justifie ton choix :
\begin{enumerate}
\item L'État camerounais construit davantage de lycées techniques dans les régions septentrionales que dans les zones déjà bien équipées.
\item Un transporteur de la ville de Douala refuse de rembourser un commerçant dont la marchandise a été détruite pendant le trajet.
\item Un tribunal condamne un entrepreneur à indemniser les riverains dont les maisons ont été fissurées par ses travaux.
\item Un père partage son héritage entre ses enfants en donnant plus au fils aîné, selon la coutume.
\end{enumerate}
\end{exercice}

\begin{exercice}[Hiérarchie des normes au Cameroun]
Le juge d'un Tribunal d'Instance de Bafoussam est saisi d'un litige foncier. Il dispose des textes suivants : la Constitution de 1996, la Convention sur l'élimination de toutes les formes de discrimination à l'égard des femmes (CEDAW), une loi ordinaire de 1974, un décret d'application et une coutume locale.
\begin{enumerate}
\item Range ces normes de la plus haute à la plus basse dans la hiérarchie des normes.
\item Justifie la place de chaque norme.
\item Que doit faire le juge si la coutume contredit la CEDAW ratifiée par le Cameroun ?
\end{enumerate}
\end{exercice}

\begin{exercice}[Droit positif et justice]
Un maire fait appliquer un arrêté municipal interdisant aux jeunes de moins de 21 ans l'accès à un stade après 18 heures, sans aucune raison de sécurité établie.
\begin{enumerate}
\item Cet arrêté relève-t-il du droit positif ? Justifie.
\item Est-il pour autant juste ? Argumente en mobilisant la distinction entre légalité et justice.
\item Rédige un paragraphe argumenté (10 à 15 lignes) montrant qu'une norme juridique peut être légale sans être juste, en t'appuyant sur un autre exemple de ton choix.
\end{enumerate}
\end{exercice}

\begin{exercice}[Kelsen et Hart]
Après avoir rappelé en quelques lignes la conception du droit de Kelsen (norme fondamentale) et celle de Hart (règles primaires et règles secondaires), réponds aux questions suivantes :
\begin{enumerate}
\item Quelle est la différence essentielle entre les deux théories ?
\item Laquelle des deux permet le mieux de comprendre le fonctionnement du système juridique camerounais ? Justifie ta réponse par deux arguments.
\end{enumerate}
\end{exercice}

\courssub{Je me prépare au Bac}

\begin{exercice}[Sujet de dissertation : loi et justice]
\textbf{Sujet :} \emph{Le respect de la loi est-il la seule condition de la justice ?}

\medskip
Traite ce sujet en rédigeant une dissertation complète selon la méthode dialectique (thèse, antithèse, synthèse). Tu devras obligatoirement :
\begin{enumerate}
\item définir dans l'introduction les notions de \cle{loi}, de \cle{droit positif} et de \cle{justice}, et construire une problématique claire ;
\item mobiliser au moins trois références philosophiques étudiées dans la leçon (par exemple Aristote, Kelsen, Rawls) ;
\item illustrer ton propos par au moins deux exemples camerounais (loi sur la décentralisation, code pénal, coutume successorale, etc.) ;
\item rédiger une conclusion qui répond nettement à la question posée et propose une ouverture.
\end{enumerate}
\end{exercice}

\begin{exercice}[Explication de texte : John Rawls]
\textbf{Texte :} « La justice est la première vertu des institutions sociales, comme la vérité l'est des systèmes de pensée. Une théorie, si élégante et économique qu'elle soit, doit être rejetée ou révisée si elle est fausse ; de même, des lois et des institutions, si efficaces et bien organisées qu'elles soient, doivent être réformées ou abolies si elles sont injustes. [...] Chaque personne possède une inviolabilité fondée sur la justice que même le bien-être de la société dans son ensemble ne peut pas renverser. [...] Les principes de la justice sont choisis derrière un voile d'ignorance : nul ne connaît sa place dans la société, sa classe sociale ou son statut, ni sa fortune dans la distribution des talents et des capacités naturels. » (John Rawls, \emph{Théorie de la justice}, § 1-4, adapté.)

\medskip
\textbf{Questions :}
\begin{enumerate}
\item Dégage la thèse défendue par Rawls dans ce texte et explique-la en quelques lignes.
\item Analyse l'argument du « voile d'ignorance » : que signifie-t-il et quel rôle joue-t-il dans la justification des principes de justice ?
\item Explique pourquoi, selon Rawls, l'efficacité d'une institution ne suffit pas à la rendre légitime.
\item Discussion : cette théorie permet-elle de résoudre le conflit entre droit coutumier et droits des femmes au Cameroun ? Rédige un développement argumenté (20 à 30 lignes) prenant clairement position.
\end{enumerate}
\end{exercice}

\begin{exercice}[Cas pratique noté : l'affaire Mvondo]
Mme Mvondo, veuve habitant Yaoundé, se fait expulser de la concession familiale par son beau-frère, au nom d'une coutume qui réserve la succession au frère aîné du défunt. Elle saisit le Tribunal d'Instance en référé.

\medskip
\textbf{Travail à faire :} rédige l'ordonnance de référé du juge, comportant obligatoirement :
\begin{enumerate}
\item un exposé sommaire des faits et de la demande ;
\item les \textbf{motifs} de la décision, en mobilisant : l'article 1\textsuperscript{er} de la Constitution camerounaise (égalité de tous devant la loi), l'article 745 du Code civil (droits successoraux du conjoint survivant), la CEDAW ratifiée par le Cameroun, et la jurisprudence de la Cour Suprême de 2017 écartant les coutumes discriminatoires contraires à l'ordre public ;
\item une analyse philosophique (5 à 10 lignes) montrant en quoi cette décision relève à la fois de la justice distributive et de la justice réparatrice ;
\item le \textbf{dispositif} (ce que le juge ordonne concrètement).
\end{enumerate}
\end{exercice}

\newpage

\coursec{Corrigés des exercices}

\begin{corrige}[Exercice 1 — QCM : notions fondamentales]
\begin{enumerate}
\item \textbf{Réponse b.} Pour Hans Kelsen (\emph{Théorie pure du droit}), le droit positif est un \cle{ordre normatif coercitif} : un ensemble de normes hiérarchisées, valables sur le territoire d'un État donné, dont l'application peut être imposée par la force publique. Les réponses a et d confondent droit et morale ou droit et coutume ; la réponse c renvoie à Rousseau (volonté générale), qui définit la loi et non le droit positif.
\item \textbf{Réponse b.} Chez Aristote (\emph{Éthique à Nicomaque}, livre V), la \cle{justice distributive} répartit les biens, les honneurs et les charges \emph{proportionnellement} au mérite ou à la contribution de chacun (égalité géométrique). La réponse a décrit la justice commutative (réparation du tort), la réponse c la légalité stricte, la réponse d la justice pénale.
\item \textbf{Réponse b.} La Constitution de 1996 est la \cle{norme suprême} de l'ordre juridique camerounais : toute norme inférieure (loi, décret, coutume) qui lui est contraire doit être écartée. La coutume n'est recevable que si elle n'est pas contraire à la Constitution, aux lois et à l'ordre public.
\item \textbf{Réponse b.} Le \cle{voile d'ignorance} de Rawls est une situation hypothétique (la « position originelle ») où les partenaires choisissent les principes de justice sans connaître la place sociale, la fortune ou les talents qu'ils auront : ce dispositif garantit l'impartialité du choix.
\end{enumerate}
\textbf{Points de vigilance :} ne pas confondre justice distributive (répartition) et justice commutative (échanges et réparations) ; bien distinguer la loi (norme votée) du droit positif (ensemble des normes en vigueur).
\end{corrige}

\begin{corrige}[Exercice 2 — Vrai ou faux justifié]
\begin{enumerate}
\item \textbf{Faux.} Une loi votée est \emph{légale} mais pas nécessairement \emph{juste} : la légalité (conformité à la procédure) ne garantit pas la justice (conformité à l'équité et aux droits fondamentaux). L'histoire et la philosophie offrent des exemples de lois injustes (lois discriminatoires), ce qui fonde l'idée de droit naturel et le contrôle de constitutionnalité.
\item \textbf{Vrai.} La justice commutative (Aristote) règle les \emph{échanges} entre particuliers — contrats, ventes, réparations — selon une égalité arithmétique : elle vise à rétablir l'équilibre rompu, indépendamment du mérite des personnes.
\item \textbf{Vrai.} Le Conseil Supérieur de la Magistrature, placé sous la présidence du Président de la République, donne son avis sur les nominations, avancements et sanctions disciplinaires des magistrats : il participe donc bien à la gestion de leur carrière, garantie (imparfaite) d'indépendance de la justice.
\item \textbf{Faux.} Le droit naturel est un droit \emph{idéal}, non écrit, fondé sur la nature humaine et la raison, antérieur et supérieur aux lois positives. Ce qui est écrit et adopté par les assemblées relève du droit positif (lois).
\end{enumerate}
\textbf{Points de vigilance :} la justification doit mobiliser une distinction conceptuelle précise (légal/juste, naturel/positif) et non une simple opinion.
\end{corrige}

\begin{corrige}[Exercice 3 — Définitions express]
\begin{enumerate}
\item \textbf{Le droit positif} : ensemble des règles juridiques effectivement en vigueur dans un État à un moment donné, sanctionnées par la force publique. \emph{Exemple :} le Code pénal camerounais, qui définit et punit les infractions.
\item \textbf{La loi} : règle générale, abstraite et obligatoire, votée par le Parlement dans les conditions prévues par la Constitution. \emph{Exemple :} la loi d'orientation de l'éducation au Cameroun.
\item \textbf{La justice réparatrice} : forme de justice qui vise à réparer le tort causé à une victime (indemnisation, restitution, réhabilitation) plutôt qu'à seulement punir. \emph{Exemple :} un entrepreneur condamné par un tribunal de Douala à indemniser des riverains dont les maisons ont été endommagées.
\item \textbf{La morale} : ensemble de règles de conduite fondées sur la distinction du bien et du mal, dont la sanction est intérieure (conscience) ou sociale (blâme), et non étatique. \emph{Exemple :} le devoir moral de secourir un accidenté sur la route Yaoundé--Douala, même en l'absence d'obligation juridique précise.
\end{enumerate}
\textbf{Points de vigilance :} chaque définition doit comporter un genre proche (ensemble de règles...) et une différence spécifique (sanction, origine) ; l'exemple doit être réellement camerounais.
\end{corrige}

\begin{corrige}[Exercice 4 — Distinctions conceptuelles]
\begin{center}
\begin{tabularx}{\linewidth}{|l|X|X|X|}\hline
\rowcolor{popblueL} \textbf{Notion} & \textbf{Origine} & \textbf{Sanction} & \textbf{Domaine} \\ \hline
Droit & L'État, la société organisée (ensemble des normes juridiques) & Sanctions étatiques : amendes, peines, exécution forcée & Rapports sociaux extérieurs, régis objectivement \\ \hline
Justice & Idéal moral et valeur philosophique (équité) & Aucune sanction matérielle propre ; jugement de la conscience et de l'opinion & Évaluation des institutions, des lois et des conduites \\ \hline
Loi & Le Parlement (organe législatif), selon une procédure constitutionnelle & Sanctions juridiques prévues par les textes (pénales, civiles) & Norme générale et abstraite s'imposant à tous \\ \hline
Morale & La conscience individuelle, la tradition, la raison & Remords, approbation ou blâme social & Conduite intérieure et intentions des personnes \\ \hline
\end{tabularx}
\end{center}
\textbf{Points de vigilance :} ne pas écrire que la loi « vient de l'État » sans préciser l'organe (Parlement) ; bien marquer que la justice est d'abord une \emph{valeur} (ce à quoi le droit doit tendre) et non une norme sanctionnée.
\end{corrige}

\begin{corrige}[Exercice 5 — Analyse de situations concrètes]
\begin{enumerate}
\item \textbf{Justice distributive.} Il s'agit de répartir des biens publics (établissements scolaires) entre des groupes selon un critère d'équité : favoriser les régions défavorisées corrige une inégalité de départ (on peut y voir une application du principe de différence de Rawls : les inégalités sont justes si elles profitent aux plus désavantagés).
\item \textbf{Justice commutative.} Le litige naît d'un contrat de transport entre particuliers : la justice doit rétablir l'égalité rompue par l'inexécution du contrat (remboursement de la marchandise détruite), indépendamment du mérite des parties.
\item \textbf{Justice réparatrice.} Le tribunal n'ordonne pas seulement une punition mais la \emph{réparation} du dommage subi (indemnisation des riverains) : il s'agit de réparer le tort et de restaurer, autant que possible, la situation antérieure.
\item \textbf{Justice distributive (discutable).} Le partage de l'héritage est une répartition de biens entre des personnes selon un critère (ici l'ordre de naissance et le sexe, fixé par la coutume). On peut discuter la \emph{justice} de ce critère au regard de l'égalité constitutionnelle : c'est un cas où la distribution coutumière peut être jugée injuste car discriminatoire.
\end{enumerate}
\textbf{Points de vigilance :} toujours nommer le critère de répartition ou le type d'échange en cause ; pour la situation 4, la justification doit distinguer le \emph{type} de justice (distributive) de la \emph{valeur} de justice (équité), qui peut faire défaut.
\end{corrige}

\begin{corrige}[Exercice 6 — Hiérarchie des normes au Cameroun]
\begin{enumerate}
\item \textbf{Classement (de la plus haute à la plus basse) :} (1) la Constitution de 1996 ; (2) la CEDAW (traité international ratifié) ; (3) la loi ordinaire de 1974 ; (4) le décret d'application ; (5) la coutume locale.
\item \textbf{Justifications :}
\begin{itemize}
\item La \emph{Constitution} est la norme suprême : elle fonde et encadre toutes les autres (article 65 : toute loi contraire est inconstitutionnelle).
\item La \emph{CEDAW}, traité dûment ratifié et publié, a une valeur supra-législative : le préambule de la Constitution de 1996 affirme l'attachement du Cameroun aux libertés internationalement reconnues, et les traités ratifiés priment sur les lois.
\item La \emph{loi ordinaire} est votée par le Parlement dans le respect de la Constitution et des traités.
\item Le \emph{décret} est un acte réglementaire pris pour l'application des lois : il est subordonné à la loi qu'il met en œuvre.
\item La \emph{coutume} n'est recevable qu'à titre subsidiaire et à condition de n'être contraire ni à la loi, ni à la Constitution, ni à l'ordre public.
\end{itemize}
\item \textbf{Si la coutume contredit la CEDAW :} le juge doit écarter la coutume et appliquer la norme supérieure. La coutume discriminatoire (par exemple en matière successorale) est contraire à l'égalité proclamée par la Constitution et par la CEDAW ; la jurisprudence de la Cour Suprême écarte les coutumes contraires à l'ordre public. Le juge statue donc sur le fondement du droit positif étatique et international.
\end{enumerate}
\textbf{Barème indicatif (sur 10) :} classement exact 3 pts ; justification de chaque norme 4 pts ; réponse à la question 3 avec référence à l'ordre public 3 pts.
\textbf{Points de vigilance :} ne pas placer la coutume au-dessus du décret ; citer le principe de supra-légalité des traités ratifiés.
\end{corrige}

\begin{corrige}[Exercice 7 — Droit positif et justice]
\begin{enumerate}
\item \textbf{Oui, cet arrêté relève du droit positif.} Il est édicté par une autorité compétente (le maire), selon les formes prévues, et s'impose sur le territoire de la commune avec une sanction possible : il appartient donc à l'ordre normatif coercitif en vigueur (au sens de Kelsen).
\item \textbf{Non, il n'est pas pour autant juste.} La légalité (conformité formelle à la procédure) ne coïncide pas nécessairement avec la justice (conformité matérielle à l'équité). Cet arrêté établit une discrimination fondée sur l'âge, sans motif de sécurité établi : il viole le principe d'égalité et restreint arbitrairement une liberté. Une norme légale peut donc être injuste.
\item \textbf{Paragraphe argumenté (corrigé-type).} On doit distinguer la validité juridique d'une norme de sa valeur morale. Une norme est juridiquement valable dès lors qu'elle est produite selon la procédure prévue par l'ordre juridique : c'est la thèse du positivisme de Kelsen, pour qui le droit est un ordre de normes hiérarchisées dont la validité ne dépend pas de leur contenu moral. Mais l'histoire montre que des normes parfaitement légales ont été profondément injustes : les lois de ségrégation raciale aux États-Unis ou de l'apartheid en Afrique du Sud étaient votées selon les formes, et pourtant elles niaient l'égalité des personnes. De même, une coutume successorale qui exclut les filles de l'héritage peut être appliquée comme une règle sociale effective sans être juste, car elle contredit l'égalité proclamée par l'article 1\textsuperscript{er} de la Constitution camerounaise. C'est pourquoi la philosophie du droit maintient l'exigence d'un droit naturel ou de principes de justice (Rawls) permettant de juger les lois elles-mêmes : le droit positif doit être évalué à la lumière de la justice. La légalité est une condition de l'ordre social, mais elle n'est pas une condition suffisante de la justice.
\end{enumerate}
\textbf{Barème indicatif (sur 10) :} question 1 : 2 pts ; question 2 : 3 pts ; paragraphe : 5 pts (thèse claire 1 pt, distinction légal/juste 1 pt, référence philosophique 1 pt, exemple 1 pt, cohérence et langue 1 pt).
\textbf{Points de vigilance :} éviter la contradiction « c'est du droit donc c'est juste » ; citer explicitement la distinction légalité/justice ; l'exemple choisi doit être réel et contrôlé.
\end{corrige}

\begin{corrige}[Exercice 8 — Kelsen et Hart]
\textbf{Rappel.} Pour \textbf{Kelsen}, le droit est une hiérarchie de normes dont chacune tire sa validité d'une norme supérieure, jusqu'à la \cle{norme fondamentale} (\emph{Grundnorm}), présupposée, qui fonde la validité de la Constitution et de tout l'ordre juridique. Pour \textbf{Hart}, le droit est l'union de \cle{règles primaires} (qui imposent des obligations : ne pas voler, respecter les contrats) et de \cle{règles secondaires} (règles sur les règles : règles de reconnaissance pour identifier le droit, de changement pour le modifier, d'adjudication pour trancher les litiges).
\begin{enumerate}
\item \textbf{Différence essentielle :} Kelsen fonde l'unité du droit sur une norme unique présupposée (construction logique et pyramidale), tandis que Hart fonde l'existence du droit sur une \emph{pratique sociale} : la règle de reconnaissance existe parce qu'elle est effectivement acceptée et utilisée par les juges et les fonctionnaires. Kelsen raisonne en logicien de la norme, Hart en sociologue de la règle.
\item \textbf{Application au Cameroun.} On peut défendre les deux réponses, pourvu qu'elles soient argumentées. \emph{Exemple de réponse favorable à Hart :}
\begin{itemize}
\item \emph{Argument 1 :} le système camerounais est pluraliste (droit étatique, droit coutumier, droit international) ; la notion de règle de reconnaissance permet d'expliquer comment les juges identifient concrètement les normes applicables (par exemple en écartant les coutumes contraires à l'ordre public), ce que la pyramide unique de Kelsen explique moins bien.
\item \emph{Argument 2 :} les règles secondaires de changement et d'adjudication décrivent bien le fonctionnement réel des institutions camerounaises (Parlement qui modifie les lois, tribunaux qui tranchent les litiges), alors que la norme fondamentale kelsienne reste une hypothèse abstraite difficile à vérifier.
\end{itemize}
\emph{Réponse inverse possible :} la pyramide de Kelsen rend mieux compte de la suprématie de la Constitution de 1996 et du contrôle de constitutionnalité.
\end{enumerate}
\textbf{Barème indicatif (sur 10) :} rappel exact des deux théories 4 pts ; différence clairement formulée 2 pts ; deux arguments pertinents et développés 4 pts.
\textbf{Points de vigilance :} ne pas présenter la norme fondamentale comme une loi écrite ; ne pas confondre règles primaires (obligations) et secondaires (organisation du système).
\end{corrige}

\begin{corrige}[Exercice 9 — Sujet de dissertation : loi et justice]
\textbf{Corrigé-type détaillé (méthode dialectique).}

\textbf{Introduction.} Dans la vie sociale, nous obéissons aux lois en croyant servir la justice ; mais l'histoire enseigne que des lois ont pu être injustes. La \cle{loi} est une règle générale et obligatoire votée par le Parlement ; le \cle{droit positif} est l'ensemble des normes juridiques en vigueur dans un État, sanctionnées par la force publique ; la \cle{justice} est la valeur qui exige que soient rendus à chacun ses droits et que les biens soient équitablement répartis. Le problème est donc le suivant : la conformité à la loi suffit-elle à réaliser la justice, ou la justice exige-t-elle un critère supérieur à la loi ? Autrement dit : le légal est-il nécessairement le juste ?

\textbf{I. Thèse — Le respect de la loi est la condition de la justice.} Sans lois, point de justice possible : la justice a besoin de règles générales, identiques pour tous, pour éviter l'arbitraire et la vengeance privée. Aristote montre que l'homme injuste est celui qui « prend plus que sa part » ; la loi fixe précisément les parts. Kelsen ajoute que le droit positif, ordre normatif coercitif, garantit la paix sociale en imposant à tous les mêmes obligations. Au Cameroun, le Code pénal protège les personnes et les biens : celui qui respecte la loi ne vole pas, ne tue pas, honore ses contrats. Obéir à la loi, c'est donc déjà réaliser une justice minimale : l'égalité devant la règle commune.

\textbf{II. Antithèse — Le respect de la loi ne suffit pas : il existe des lois injustes.} La légalité n'est qu'une conformité formelle à la procédure ; elle ne dit rien du contenu. Les lois de l'apartheid étaient légales et pourtant profondément injustes. Au Cameroun, une coutume successorale qui déshérite la veuve et les filles peut être appliquée comme une règle effective sans être juste, car elle contredit l'égalité proclamée par l'article 1\textsuperscript{er} de la Constitution. C'est pourquoi la tradition du droit naturel affirme l'existence de droits antérieurs et supérieurs aux lois positives, et pourquoi la désobéissance civile (Antigone chez Sophocle, Gandhi, Martin Luther King) peut être un devoir moral face à une loi injuste. La justice est donc un critère qui juge la loi elle-même.

\textbf{III. Synthèse — La loi est une condition nécessaire mais non suffisante de la justice.} Il faut dépasser l'opposition : la justice pleine exige des lois \emph{et} des lois justes. Rawls propose un critère : des institutions sont justes si leurs principes seraient choisis derrière un « voile d'ignorance », où chacun ignore la place qu'il occupera ; ainsi, nul ne voterait une loi discriminatoire dont il pourrait être la victime. Concrètement, cela suppose des mécanismes de correction du droit par la justice : le contrôle de constitutionnalité, la primauté des traités (CEDAW) sur les coutumes discriminatoires, la jurisprudence de la Cour Suprême écartant les coutumes contraires à l'ordre public. La décentralisation camerounaise, lorsqu'elle rapproche les services publics des populations défavorisées, illustre une loi qui tend vers la justice distributive. Le respect de la loi est donc la première condition de la justice, mais la justice demeure l'idéal régulateur qui oblige à réformer les lois injustes.

\textbf{Conclusion.} Le respect de la loi n'est pas la seule condition de la justice : il en est la condition nécessaire, car sans règle commune règne l'arbitraire, mais non la condition suffisante, car la loi peut être injuste. La justice est à la fois le fondement et le juge de la loi. Ouverture : cette exigence d'une loi juste ne suppose-t-elle pas, au-delà des institutions, la vertu personnelle du citoyen et du juge ?

\textbf{Barème indicatif (sur 20) :} introduction (définitions, problématique, annonce du plan) 4 pts ; thèse argumentée 4 pts ; antithèse argumentée 4 pts ; synthèse dépassant l'opposition 4 pts ; conclusion et ouverture 2 pts ; langue, orthographe, cohérence 2 pts.
\textbf{Points de vigilance :} les trois références philosophiques (Aristote, Kelsen, Rawls) doivent être \emph{utilisées} dans l'argumentation et non simplement citées ; les deux exemples camerounais doivent être précis ; éviter la synthèse-faiblesse qui se contente de dire « les deux ont raison ».
\end{corrige}

\begin{corrige}[Exercice 10 — Explication de texte : John Rawls]
\begin{enumerate}
\item \textbf{Thèse.} Rawls affirme que la justice est la \emph{première} vertu des institutions sociales : elle est le critère absolu auquel toute loi et toute institution doivent être soumises. De même qu'une théorie fausse doit être rejetée quelle que soit son élégance, une institution injuste doit être réformée ou abolie quelle que soit son efficacité. La justice fonde en outre une « inviolabilité » de la personne : les droits de chacun ne peuvent être sacrifiés au bien-être général de la société.
\item \textbf{Le voile d'ignorance.} C'est un dispositif hypothétique : pour choisir les principes de justice, on imagine des partenaires qui ignorent leur future place sociale, leur classe, leur fortune et leurs talents naturels. Son rôle est de garantir l'\emph{impartialité} : ne sachant pas s'il sera riche ou pauvre, homme ou femme, puissant ou faible, chacun choisira des principes équitables qui protègent notamment les plus défavorisés. Le voile transforme l'intérêt personnel en garantie d'équité : on ne peut pas tricher en faveur d'une position dont on ignore si on l'occupera.
\item \textbf{Efficacité et légitimité.} L'efficacité est une qualité instrumentale : une institution efficace atteint bien ses fins (ordre, production, stabilité). Mais elle ne dit rien de la \emph{valeur} de ces fins ni de la manière dont les charges et les avantages sont répartis. Une administration efficace peut organiser l'injustice (exploitation, discrimination). La légitimité exige un critère moral indépendant : la justice. C'est pourquoi Rawls refuse de sacrifier les droits d'une personne au « bien-être de la société dans son ensemble » : le calcul utilitariste du plus grand bonheur du plus grand nombre peut justifier le sacrifice d'une minorité, ce que la justice interdit.
\item \textbf{Discussion (corrigé-type).} Le conflit entre droit coutumier et droits des femmes est réel au Cameroun : certaines coutumes successorales réservent l'héritage aux héritiers mâles et excluent la veuve de la concession familiale, au nom de la tradition et de la cohésion du lignage. La théorie de Rawls offre un instrument puissant pour trancher ce conflit. Plaçons-nous derrière le voile d'ignorance : un partenaire qui ignore s'il naîtra homme ou femme, fils aîné ou cadette, époux ou veuve, ne choisira jamais une règle qui déshérite la moitié de l'humanité, car il risquerait lui-même d'en être la victime. Le principe d'égalité et le principe de différence (les inégalités ne sont justes que si elles profitent aux plus désavantagés) conduisent donc à rejeter les coutumes discriminatoires : elles aggravent la situation des femmes, qui comptent parmi les plus défavorisées. De plus, la première thèse du texte — l'inviolabilité de la personne fondée sur la justice — interdit de sacrifier les droits des femmes au « bien-être » ou à la « cohésion » du groupe : l'argument de la tradition, qui est un argument d'efficacité sociale, ne suffit pas à légitimer l'institution coutumière. On objectera que le droit coutumier exprime une identité culturelle et que le droit positif ne doit pas tout uniformiser. L'objection est sérieuse : la justice ne doit pas devenir un prétexte au mépris des cultures. Mais Rawls ne demande pas d'abolir les coutumes ; il demande que toute norme, coutumière ou étatique, soit réformée si elle est injuste. Or une coutume qui prive une veuve de son toit viole un droit fondamental de la personne, que la Constitution camerounaise (article 1\textsuperscript{er}) et la CEDAW consacrent. La jurisprudence de la Cour Suprême, qui écarte les coutumes contraires à l'ordre public, applique concrètement cette exigence. En définitive, la théorie de Rawls ne règle pas mécaniquement tous les conflits, mais elle fournit le critère décisif : aucune tradition ne peut justifier que des personnes soient traitées comme de moindres personnes. Le dialogue entre droit coutumier et droit positif reste nécessaire, mais il doit se tenir sous l'autorité de la justice.
\end{enumerate}
\textbf{Barème indicatif (sur 20) :} thèse dégagée et expliquée 3 pts ; analyse du voile d'ignorance 4 pts ; distinction efficacité/légitimité 3 pts ; discussion argumentée (prise de position claire, objection traitée, application camerounaise) 8 pts ; langue 2 pts.
\textbf{Points de vigilance :} ne pas résumer le texte au lieu de l'expliquer ; dans la discussion, prendre \emph{clairement} position et traiter au moins une objection ; mobiliser le texte (citations courtes) à l'appui de chaque réponse.
\end{corrige}

\begin{corrige}[Exercice 11 — Cas pratique noté : l'affaire Mvondo]
\textbf{Corrigé-type : ordonnance de référé.}

\textbf{RÉPUBLIQUE DU CAMEROUN — Paix -- Travail -- Patrie}

\textbf{TRIBUNAL D'INSTANCE DE \ldots — ORDONNANCE DE RÉFÉRÉ N°\ldots}

\textbf{1. Exposé des faits et de la demande.} Mme Mvondo, veuve, demeurant à Yaoundé, a été expulsée de la concession familiale par son beau-frère, M. X., qui invoque une coutume réservant la succession au frère aîné du défunt. La requérante saisit le juge des référés aux fins de réintégration dans les lieux et de sauvegarde de ses droits.

\textbf{2. Motifs.}
\begin{itemize}
\item \emph{Sur le principe d'égalité :} l'article 1\textsuperscript{er} de la Constitution du 18 janvier 1996 proclame que tous les citoyens sont égaux en droits et en devoirs, et que la loi garantit à tous les citoyens les droits et libertés fondamentales. Une coutume qui prive une veuve de tout droit successoral en raison de son sexe institue une discrimination incompatible avec ce principe.
\item \emph{Sur les droits du conjoint survivant :} la loi n° 72-5 du 26 juin 1972 relative au régime matrimonial et aux successions reconnaît au conjoint survivant des droits successoraux sur la succession de son époux décédé. L'expulsion de Mme Mvondo de la concession familiale méconnaît directement ces droits.
\item \emph{Sur le droit international :} le Cameroun a ratifié la Convention sur l'élimination de toutes les formes de discrimination à l'égard des femmes (CEDAW), qui oblige l'État à éliminer les pratiques coutumières discriminatoires, notamment en matière de mariage et de succession. Ce traité, dûment ratifié, s'impose aux juridictions internes.
\item \emph{Sur la jurisprudence :} la Cour Suprême, dans sa jurisprudence constante (notamment arrêt de 2017), écarte les coutumes discriminatoires contraires à l'ordre public. La coutume invoquée par le défendeur est donc irrecevable.
\end{itemize}
Il en résulte que l'expulsion litigieuse repose sur une norme coutumière incompatible avec la hiérarchie des normes et qu'elle cause à la requérante un trouble manifestement illicite appelant une mesure urgente.

\textbf{3. Analyse philosophique.} Cette décision relève d'abord de la \cle{justice distributive} : il s'agit de répartir un bien (la succession, la concession) entre plusieurs ayants droit selon un critère juste — l'égalité des personnes et les droits reconnus par la loi — et non selon un critère arbitraire (le sexe, le rang de naissance). En refusant la répartition coutumière discriminatoire, le juge substitue une distribution fondée sur l'égalité de dignité, conforme au principe d'égalité et, chez Rawls, au principe de différence qui protège les plus défavorisés, telle la veuve. Elle relève ensuite de la \cle{justice réparatrice} : la requérante a subi un tort (l'expulsion de son domicile) ; la réintégration ordonnée répare ce tort et restaure, autant que possible, la situation antérieure. La justice apparaît ainsi comme la valeur qui corrige le droit coutumier par le droit positif et répare la victime : elle est à la fois idéal de répartition équitable et restitution concrète du droit violé.

\textbf{4. Dispositif.} \textbf{PAR CES MOTIFS,} le juge des référés, statuant publiquement et en premier ressort : ordonne la réintégration immédiate de Mme Mvondo dans la concession familiale ; fait interdiction à M. X. de troubler sa jouissance des lieux, sous astreinte de 50 000 FCFA par jour de retard ; dit que les droits successoraux définitifs seront tranchés au fond conformément au Code civil ; condamne le défendeur aux dépens.

\textbf{Barème indicatif (sur 20) :} exposé des faits clair 2 pts ; quatre motifs juridiques exacts et articulés 8 pts (2 pts chacun) ; analyse philosophique mobilisant correctement les deux formes de justice 6 pts ; dispositif concret et cohérent avec les motifs 3 pts ; présentation et langue 1 pt.
\textbf{Points de vigilance :} chaque motif doit citer le texte \emph{et} l'appliquer aux faits (ne pas se contenter d'énumérer les articles) ; l'analyse philosophique ne doit pas confondre justice distributive (critère de répartition) et justice réparatrice (réparation du tort) ; le dispositif doit découler logiquement des motifs et rester une mesure d'urgence (référé).
\end{corrige}

\newpage

\coursec{Fiche bilan}

\begin{popfigure}
\centering
\begin{tikzpicture}[mindmap, grow cyclic, every node/.style={concept, font=\small},
  root concept/.append style={concept color=popblue, font=\bfseries},
  level 1/.append style={level distance=3.4cm, sibling angle=72},
  level 2/.append style={level distance=2.2cm, font=\scriptsize}]
\node[root concept]{Droit et Justice}
  child[concept color=poporange]{ node {Définitions}
    child { node {Droit : ensemble de règles} }
    child { node {Justice : rendre à chacun son dû} }
  }
  child[concept color=popgreen]{ node {Distinctions}
    child { node {Loi / Morale / Droit} }
    child { node {Droit objectif / subjectif} }
  }
  child[concept color=poppurple]{ node {Fondements}
    child { node {Droit positif (Kelsen)} }
    child { node {Droit naturel (Aristote, Cicéron)} }
  }
  child[concept color=poppink]{ node {Formes de justice}
    child { node {Distributive} }
    child { node {Commutative} }
    child { node {Pénale, sociale} }
  }
  child[concept color=popgold]{ node {Justice au Cameroun}
    child { node {Institutions judiciaires} }
    child { node {Défis : accès, lenteur} }
  };
\end{tikzpicture}
\legende{Carte mentale de la leçon : le Droit et la Justice (5 branches essentielles à retenir).}
\end{popfigure}

\begin{bilan}
\textbf{Définitions clés à connaître par cœur}
\begin{itemize}
  \item \cle{Droit} : ensemble des règles juridiques qui régissent la vie en société et dont le respect est garanti par l'autorité publique.
  \item \cle{Justice} : vertu et principe consistant à \emph{rendre à chacun ce qui lui est dû} (\emph{suum cuique tribuere}, Ulpien) ; aussi institution chargée de dire le droit.
  \item \cle{Loi} : règle écrite, générale et obligatoire, votée par le Parlement.
  \item \cle{Morale} : ensemble de règles de conduite fondées sur la conscience individuelle, sanctionnées intérieurement (culpabilité), non par la contrainte publique.
  \item \cle{Droit objectif} : l'ensemble des règles juridiques en vigueur ; \cle{droit subjectif} : prérogative reconnue à un individu (droit de vote, droit de propriété).
\end{itemize}

\textbf{Thèses et auteurs à mobiliser}
\begin{center}
\begin{tabularx}{\linewidth}{|l|X|}
\hline
\rowcolor{popblueL} \textbf{Auteur} & \textbf{Thèse à retenir} \\
\hline
Aristote & Justice distributive (à chacun selon son mérite) et justice commutative (égalité dans les échanges et réparations). \\
\hline
Cicéron / stoïciens & Droit naturel : loi morale universelle, antérieure et supérieure aux lois humaines. \\
\hline
Kelsen & Positivisme juridique : le droit est un système de normes hiérarchisées ; sa validité ne dépend pas de sa moralité. \\
\hline
Hobbes & Sans État ni lois, la vie est «~solitaire, misérable, dangereuse, animale et brève~» : le droit fonde la paix. \\
\hline
Rawls & La justice comme équité : principes choisis sous le «~voile d'ignorance~». \\
\hline
\end{tabularx}
\end{center}

\textbf{Méthodes express}
\begin{itemize}
  \item \emph{Dissertation} : définir les termes du sujet dès l'introduction, problématiser (ex. : une loi injuste est-elle encore une loi ?), puis plan dialectique (thèse / antithèse / dépassement).
  \item \emph{Explication de texte} : situer l'auteur, dégager la thèse, suivre l'argumentation pas à pas, discuter avec un autre auteur.
  \item \emph{Application concrète} : toujours illustrer par une situation camerounaise (procès, héritage, corruption, accès au juge).
\end{itemize}
\end{bilan}

\begin{aretenir}[Checklist avant le Bac]
\begin{itemize}
  \item $\square$ Je sais définir le droit, la justice, la loi et la morale en une phrase chacun.
  \item $\square$ Je distingue droit objectif et droit subjectif avec un exemple.
  \item $\square$ Je sais opposer droit positif (Kelsen) et droit naturel (Cicéron).
  \item $\square$ Je connais la différence entre justice distributive et justice commutative chez Aristote.
  \item $\square$ Je sais citer au moins trois auteurs avec leur thèse exacte.
  \item $\square$ Je peux problématiser un sujet sous forme de question.
  \item $\square$ Je maîtrise le plan dialectique en trois parties.
  \item $\square$ Je sais illustrer chaque notion par un exemple camerounais concret.
  \item $\square$ Je connais les institutions de justice au Cameroun (tribunaux, Cour suprême, Conseil constitutionnel).
  \item $\square$ Je rédige une introduction en trois temps : accroche, définitions, problématique et annonce du plan.
  \item $\square$ Je relis ma copie : orthographe, connecteurs logiques, conclusion qui répond à la problématique.
\end{itemize}
\end{aretenir}

\findecours

\end{document}
